% ==== XI ====
\yoschapter{XI}{Normed Rings and Spectral Representation}
\label{chap:XI}

A linear space $A$ over a scalar field $(F)$ is said to be an algebra or a ring over $(F)$, if to each pair of elements $x, y \in A$ a unique product $xy \in A$ is defined with the properties: 
\begin{equation}\label{eq:XI.0.1}
\begin{yoslines}
&(x y) z = x (y z) (\text{ associativity }), x (y + z) = x y + x z (\text{ distributivity }), \alpha \beta (x y) = (\alpha x) (\beta y).
\end{yoslines}
\end{equation}

If there exists a unit element $e$ such that $ex = xe = x$ for every $x \in A$, then $A$ is said to be an algebra with a unit. A unit $e$ of $A$, if it exists, is uniquely determined. For, if $e'$ be another unit of $A$, then we must have $ee' = e = e'$. If the multiplication is commutative, i.e., $xy = yx$ for every pair $x, y \in A$, then $A$ is called a commutative algebra. Let $A$ be an algebra with a unit $e$. If, for an $x \in A$, there exists an $x' \in A$ such that $xx' = x'x = e$, then $x'$ is called an inverse of $x$. An inverse $x'$ of $x$, if it exists, is uniquely determined. For, if $x''$ be another inverse of $x$, then we must have 
\[
x ^{\prime \prime} (x x ^{\prime}) = x ^{\prime \prime} e = x ^{\prime \prime} = (x ^{\prime \prime} x) x ^{\prime} = e x ^{\prime} = x ^{\prime}.
\]

Thus we shall denote by $x^{-1}$ the inverse of $x$ if $x$ has an inverse. 

An algebra is called a Banach algebra, or in short a $B$-algebra if it is a $B$-space and satisfies 
\begin{equation}\label{eq:XI.0.2}
\left\| x y \right\| \leqq \left\| x \right\| \left\| y \right\|.
\end{equation}

The inequality 
\[
\begin{aligned}
&\left\| x _{n} y _{n} - x y \right\|  \leqq \left\| x _{n} (y _{n} - y) \right\| + \left\| (x _{n} - x) y \right\| \leqq \left\| x _{n} \right\| \left\| (y _{n} - y) \right\| + \left\| (x _{n} - x) \right\| \left\| y \right\|
\end{aligned}
\]
shows that $xy$ is a continuous function of both variables together. 

\begin{example}\label{ex:XI.0.Example.1} Let $X$ be a $B$-space. Then $L(X, X)$ is a $B$-algebra with a unit by the operator sum $T + S$ and operator product $TS$; the identity operator $I$ is the unit of this algebra $L(X, X)$, and the operator norm $\|T\|$ is the norm of the element $T$ of this algebra $L(X, X)$. 
\end{example}

\begin{example}\label{ex:XI.0.Example.2} Let $S$ be a compact topological space. Then $C(S)$ is a $B$-algebra by $(x_1 + x_2)(s) = x_1(s) + x_2(s), (\alpha x)(s) = \alpha x(s), (x_1x_2)(s) = x_1(s)x_2(s)$ and $\|x\| = \sup\limits_{s\in S}|x(s)|$. 
\end{example}

\begin{example}\label{ex:XI.0.Example.3} Let $B$ be the totality of continuous functions $x(s)$, $0 \leqq s \leqq 1$, which are representable as absolutely convergent Fourier series: 
\begin{equation}\label{eq:XI.0.3}
x (s) = \sum\limits_{n = - \infty} ^{\infty} c _{n} e ^{2 \pi i n s} \quad \text{ with } \quad \sum\limits_{n = - \infty} ^{\infty} | c _{n} | < \infty .
\end{equation}
\end{example}

Then it is easy to see that $B$ is a commutative $B$-algebra with a unit by the ordinary function sum and function multiplication, normed by 
\begin{equation}\label{eq:XI.0.4}
\left\| x \right\| = \sum\limits_{n = - \infty} ^{\infty} \left| c _{n} \right|.
\end{equation}

In the last two examples, the unit is given by the function $e(s) \equiv 1$ and $\|e\| = 1$. In the following sections, we shall be concerned with the commutative $B$-algebra with the unit $e$ such that 
\begin{equation}\label{eq:XI.0.5}
\left\| e \right\| = 1.
\end{equation}

Such an algebra is called a normed ring. 

A Historical Sketch. The notion of Banach algebras was introduced in analysis by M. \textsc{Nagumo} \cite{nagumo215}. He proved that Cauchy's complex function theory can be extended to functions with values in such an algebra, and applied it to the investigation of the resolvent of a bounded linear operator around an isolated singular point. The result is an abstract treatment of those given in our Chapter~\ref{chap:VIII}, \ref{sec:VIII.8}. K. \textsc{Yosida} \cite{yosida354} proved that a connected group embedded in a $B$-algebra is a Lie group iff the group is locally compact. This result is an extension of a result due to J. \textsc{von} \textsc{Neumann} \cite{neumann233} concerning matrix groups. Cf. E. \textsc{Hille}-R. S. \textsc{Phillips} \cite{hille97}, in which the result of K. \textsc{Yosida} \cite{yosida354} is reproduced. 

The ideal theory of normed rings was initiated by I. M. \textsc{Gelfand} \cite{gelfand76b}. He has shown that such a ring can be represented as the ring of continuous functions defined on the space of maximal ideals of the ring. By virtue of this representation, we can give an integration free treatment of the spectral resolution of bounded normal operators in a Hilbert space; see K. \textsc{Yosida} \cite{yosida355}. This result will be exposed in the following sections. The Gelfand representation may also be applied to a new proof of the Tauberian theorem of N. \textsc{Wiener} \cite{wiener340}. We shall expose this application in the last section of this chapter. For further details about $B$-algebras, see N. A. \textsc{Naimark} \cite{naimark221}, C. E. \textsc{Rickart} \cite{rickart271} and I. M. \textsc{Gelfand}-D. A. \textsc{Raikov}-G. E. \textsc{Šilov} \cite{gelfand79}.\yosstarfootnote{See also Supplementary Notes, p.~\pageref{notes:XI}.} 


\yossection{1}{Maximal Ideals of a Normed Ring}
\label{sec:XI.1}

We shall be concerned with a commutative $B$-algebra $B$ with a unit $e$ such that $\| e \| = 1$. 

\begin{definition}\label{def:XI.1.Definition.1} $A$ subset $J$ of $B$ is called an ideal of $B$ if $x, y \in J$ implies that $(\alpha x + \beta y) \in J$ and $zx \in J$ for every $z \in B$. $B$ itself and $\{0\}$ are ideals of $B$. Ideals other than $B$ and $\{0\}$ are called non-trivial ideals. A non-trivial ideal $J$ is said to be a maximal ideal if there exists no non-trivial ideal containing $J$ as a proper subset. 
\end{definition}

\begin{proposition}\label{prop:XI.1.Proposition.1} Each non-trivial ideal $J_0$ of $B$ is contained in a maximal ideal $J$. 
\end{proposition}

\begin{proof}
 Let $[J_0]$ be the set of all non-trivial ideals containing $J_0$. We order the ideals of $[J_0]$ by inclusion relation, that is, we denote $J_1 \prec J_2$ if $J_1$ is a subset of $J_2$. Suppose that $\{J_\alpha\}$ is a linearly ordered subset of $[J_0]$ and put $J_\beta = \bigcup\limits_{J_\alpha \in \{J_\alpha\}} J_\alpha$. We shall show that $J_\beta$ is an upper bound of $\{J_\alpha\}$. For, if $x, y \in J_\beta$, then there exist ideals $J_{\alpha_1}$ and $J_{\alpha_2}$ such that $x \in J_{\alpha_1}$ and $y \in J_{\alpha_2}$. Since $\{J_\alpha\}$ is linearly ordered, $J_{\alpha_1} \prec J_{\alpha_2}$ (or $J_{\alpha_1} > J_{\alpha_2}$) and so $x$ and $y$ both belong to $J_{\alpha_2}$; consequently $(x - y) \in J_{\alpha_2} \subseteq J_\beta$ and $zx \in J_{\alpha_2} \subseteq J_\beta$ for any $z \in B$. This proves that $J_\beta$ is an ideal. Since the unit element $e$ is not contained in any $J_\alpha, e$ is not contained in $J_\beta = \bigcup\limits_{J_\alpha \in \{J_\alpha\}} J_\alpha$. 

Thus $J_{\beta}$ is a non-trivial ideal containing every $J_{\alpha}$. Therefore, by Zorn's Lemma, there exists at least one maximal ideal which contains $J_{0}$. 
\end{proof}

\begin{corollaryu} An element $x$ of $B$ has the inverse $x^{-1} \in B$ such that $x^{-1}x = xx^{-1} = e$ iff $x$ is contained in no maximal ideal. 
\end{corollaryu}

\begin{proof}
 If $x^{-1} \in B$ exists, then any ideal $J \ni x$ must contain $e = xx^{-1}$ so that $J$ must coincide with $B$ itself. Let conversely, $x$ be contained in no maximal ideal. Then the ideal $xB = \{xb; b \in B\} \neq \{0\}$ must coincide with $B$ itself, since, otherwise, there exists at least one maximal ideal containing $xB \ni x = xe$. It follows that $xB = B$, and so there must exist an element $b \in B$ such that $xb = e$. By the commutativity of $B$, we have $xb = bx = e$, that is, $b = x^{-1}$. 
\end{proof}

\begin{proposition}\label{prop:XI.1.Proposition.2} $A$ maximal ideal $J$ is a closed linear subspace of $B$. 
\end{proposition}

\begin{proof}
 By the continuity of the algebraic operations (addition, multiplication and scalar multiplication) in $B$, the strong closure $J^a$ is also an ideal containing $J$. Suppose $J^a \neq J$. Then $J^a = B$, because of the maximality of the ideal $J$. Thus $e \in J^a$, and so there exists an $x \in J$ such that $\| e - x \| < 1$. $x$ has the inverse $x^{-1} \in B$ which is given by Neumann's series $e + (e - x) + (e - x)^2 + \cdots$ 
\[
e + (e - x) + (e - x) ^{2} + \dots \]

For, by $\|(e-x)^{n}\|\leqq\|e-x\|^{n}$, the series converges to an element $\in B$ which is the inverse of $x$, as may be seen by multiplying the series by $x=e-(e-x)$. Hence $e=x^{-1}x\in J$ and so $J$ cannot be a maximal ideal. 
\end{proof}

\begin{proposition}\label{prop:XI.1.Proposition.3} For any ideal $J$ of $B$, we write 

$x \equiv y (\bmod J)$ or $x \sim y (\bmod J)$ or in short $x \sim y$, if $(x - y) \in J$. Then $x \sim y$ is an equivalence relation, that is, we have 
\begin{equation}\label{eq:XI.1.1}
\begin{yoslines}
&x \sim x (\mathrm{reflexivity}), x \sim y \mathrm{implies} y \sim x (\mathrm{symmetry}), x \sim y \mathrm{and} y \sim z \mathrm{implies} x \sim z (\mathrm{transitivity}).
\end{yoslines}
\end{equation}
\end{proposition}

We denote by $\bar{x}$ the set $\{y; (y-x) \in J\}$; it is called the class (mod $J$) containing $x$. Then the classes $(\overline{x+y})$, $\alpha x$ and $(\overline{xy})$ are determined independently of the choice of elements $x$ and $y$ from the classes $\bar{x}$ and $\bar{y}$, respectively. 

Proof We have to show that $x \sim x'$, $y \sim y'$ implies that $(x + y) \sim (x' + y')$, $\alpha x \sim \alpha x'$ and $xy \sim x'y'$. These are clear from the condition that $J$ is an ideal. For instance, we have $xy - x'y' = (x - x')y + x'(y - y') \in J$ by $(x - x') \in J$ and $(y - y') \in J$. 

\begin{corollaryu} The set of classes $\overline{x}$ (mod $J$) thus constitutes an algebra by 
\begin{equation}\label{eq:XI.1.2}
\overline {{{x}}} + \overline {{{y}}} = \overline {{{x + y}}}, \alpha \overline {{{x}}} = \overline {{{\alpha x}}}, \overline {{{x}}} \overline {{{y}}} = \overline {{{x y}}}.
\end{equation}
\end{corollaryu}

\begin{definition}\label{def:XI.1.Definition.2} The above obtained algebra is called the residue class algebra of $B \pmod{J}$ and is denoted by $B / J$. Thus the mapping $x \to \overline{x}$ of the algebra $B$ onto $\overline{B} = B / J$ is a homomorphism, that is, relation \eqref{eq:XI.1.2} holds. 
\end{definition}

\begin{proposition}\label{prop:XI.1.Proposition.4} Let $J$ be a maximal ideal of $B$. Then $\overline{B}=B/J$ is a field, that is, each non-zero element $\bar{x}\in\overline{B}$ has an inverse $\bar{x}^{-1}\in\overline{B}$ such that $\bar{x}^{-1}\bar{x}=\bar{x}\bar{x}^{-1}=\bar{e}$. 
\end{proposition}

\begin{proof}
 Suppose the inverse $\bar{x}^{-1}$ does not exist. Then the set $\bar{x}\overline{B}=\{\bar{x}\bar{b};\bar{b}\in\overline{B}\}$ is an ideal of $\overline{B}$. It is non-trivial since it does not contain $\bar{e}$, but does contain $\bar{x}\neq0$. The inverse image of an ideal by the homomorphism is an ideal. Therefore, $B$ contains a non-trivial ideal containing $J$ as a proper subset, contrary to the maximality of the ideal $J$. 
\end{proof}

We are now able to prove 

\begin{theoremu} Let $B$ be a normed ring over the field of complex numbers, and $J$ a maximal ideal of $B$. Then the residue class algebra $\overline{B}=B/J$ is isomorphic to the complex number field, in the sense that each $\overline{x}\in\overline{B}$ is represented uniquely as $\overline{x}=\xi\overline{e}$, where $\xi$ is a complex number. 
\end{theoremu}

\begin{proof}
 We shall prove that $\overline{B}=B/J$ is a normed ring by the norm 
\begin{equation}\label{eq:XI.1.3}
\left\| \overline {{x}} \right\| = \inf\limits _ {x \in \overline {{x}}} \left\| x \right\|. \end{equation}

If this is proved, then $B / J$ is a normed field and so, by the Gelfand-Mazur theorem in Chapter~\ref{chap:V}, \ref{sec:V.3}, $\overline{B} = B / J$ is isomorphic to the complex number field. 

Now we have $\|\alpha \overline{x}\| = |\alpha \|\overline{x}\|$, and $\|\overline{x} + \overline{y}\| = \inf\limits_{x \in \overline{x}, y \in \overline{y}} \|x + y\| \leqq \inf\limits_{x \in \overline{x}} \|x\| + \inf\limits_{y \in \overline{y}} \|y\| = \|\overline{x}\| + \|\overline{y}\|; \|\overline{x} \overline{y}\| \leqq \|\overline{x}\| \|\overline{y}\|$ is proved similarly. If $\|\overline{x}\| = 0$, then there exists a sequence $\{x_n\} \subseteq \overline{x}$ such that $\slim\limits_{n \to \infty} x_n = 0$. Hence, for any $x \in \overline{x}$, $(x - x_n) \in J$ and so $\slim\limits_{n \to \infty} (x - x_n) = x$ which proves that $x \in J^a = J$, that is, $\overline{x} = \overline{0}$. Hence $\|\overline{x}\| = 0$ is equivalent to $\overline{x} = \overline{0}$. We have $\|\overline{e}\| \leqq \|e\| = 1$. If $\|\overline{e}\| < 1$, then there exists an element $x \in J$ such that $\|e - x\| < 1$. As in the proof of Proposition~\ref{prop:XI.1.Proposition.2}, the inverse $x^{-1}$ exists, which is contradictory to the Corollary of Proposition~\ref{prop:XI.1.Proposition.1}. Thus we must have $\|\overline{e}\| = 1$. Finally, since $B$ is a $B$-space and $J$ is a closed linear subspace by Proposition~\ref{prop:XI.1.Proposition.2}, the factor space $\overline{B} = B / J$ is complete with respect to the norm \eqref{eq:XI.1.3} (see Chapter~\ref{chap:I}, \ref{sec:I.11}). We have thus proved the Theorem. 
\end{proof}

\begin{corollaryu} We shall denote by $x(J)$ the number $\xi$ in the representation $\overline{x} = \xi \overline{e}$. Thus, for each $x \in B$, we obtain a complex-valued function $x(J)$ defined on the set $\{J\}$ of all the maximal ideals of $B$. Then we have 
\begin{equation}\label{eq:XI.1.4}
\begin{gathered}
(x + y) (J) = x (J) + y (J), (\alpha x) (J) = \alpha x (J), (x y) (J) = x (J) y (J), \text{ and } e (J) \equiv 1.
\end{gathered}
\end{equation}
\end{corollaryu}

We have, moreover, 
\begin{equation}\label{eq:XI.1.5}
\sup\limits _ {J \in \{J \}} | x (J) | \leqq \| x \|,
\end{equation}
and 
\begin{equation}\label{eq:XI.1.6}
\sup\limits _ {J \in \{J \}} | x (J) | = 0 \text{ implies } x = 0 \text{ iff } \bigcap\limits_ {J \in \{J \}} J = \{0 \}.
\end{equation}

\begin{proof}
 The mapping $x \to \bar{x} = x(J) \bar{e}$ of the algebra $B$ onto the residue class algebra $\overline{B} = B/J$ is a homomorphism, that is, relation \eqref{eq:XI.1.2} holds. Hence we have \eqref{eq:XI.1.4}. Inequality \eqref{eq:XI.1.5} is proved by 
\[
| \xi | = | \xi | \| \overline {{e}} \| = \| \overline {{x}} \| = \inf\limits _ {x \in \overline {{x}}} \| x \| \leqq \| x \|. \]

Property \eqref{eq:XI.1.6} is clear, since $x(J) = 0$ identically on $\{J\}$ iff $x \in \bigcap\limits_{J \in \{J\}} J$. 
\end{proof}

\begin{definition}\label{def:XI.1.Definition.3} The representation 
\begin{equation}\label{eq:XI.1.7}
x \rightarrow x (J)
\end{equation}
of the normed ring $B$, by the ring of functions $x(J)$ defined on the set $\{J\}$ of all the maximal ideals $J$ of $B$, is called the Gelfand-representation of $B$. 
\end{definition}


\yossection{2}{The Radical. The Semi-simplicity}
\label{sec:XI.2}

\begin{definition}\label{def:XI.2.Definition.1} Let $B$ be a normed ring over the complex number field, and $\{J\}$ the totality of the maximal ideals $J$ of $B$. Then the ideal $\bigcap\limits_{J \in \{J\}} J$ is called the radical of the ring $B$. $B$ is said to be semi-simple if its radical $R = \bigcap\limits_{J \in \{J\}} J$ reduces to the zero ideal $\{0\}$. 
\end{definition}

\begin{theorem}\label{thm:XI.2.Theorem.1} For any $x \in B$, $\lim\limits_{n \to \infty} \| x^n \|^{1/n}$ exists and we have 
\begin{equation}\label{eq:XI.2.1}
\lim\limits _{n \rightarrow \infty} \| x ^{n} \| ^{1/n} = \sup\limits _ {J \in \{J \}} | x (J) |.
\end{equation}
\end{theorem}

\begin{proof}
 Set $\alpha = \sup\limits_{J \in \{J\}} |x(J)|$. Then, by $\|x^n\| \geqq |x^n(J)| = |x(J)|^n$, we have $\|x^n\| \geqq \alpha^n$, and so $\lim\limits_{n \to \infty} \|x^n\|^{1/n} \geqq \alpha$. We have thus to prove $\varlimsup\limits_{n \to \infty} \|x^n\|^{1/n} \leqq \alpha$. 

Let $|\beta| > \alpha$. Then, for any $J \in \{J\}$, $x(J) - \beta \neq 0$, i.e., $(x - \beta e) \bar{\in} J$. Hence the inverse $(\beta e - x)^{-1}$ exists. Setting $\beta^{-1} = \lambda$, we see that the inverse $(\beta e - x)^{-1} = \lambda (e - \lambda x)^{-1}$ exists whenever $|\lambda| < \alpha^{-1}$. Moreover, as in Theorem~\ref{thm:VIII.2.Theorem.1} in Chapter~\ref{chap:VIII}, \ref{sec:VIII.2}, we see that $\lambda (e - \lambda x)^{-1}$ is, for $|\lambda| < \alpha^{-1}$, holomorphic in $\lambda$. Hence we have the Taylor expansion 
\[
\lambda (e - \lambda x) ^{-1} = \lambda (e + \lambda x _{1} + \lambda^{2} x _{2} + \dots + \lambda^{n} x _{n} + \dots).
\]

That $x_{n} = x^{n}$ may be seen by Neumann's series $(e - \lambda x)^{-1} = \sum\limits_{n=0}^{\infty} \lambda^{n} x^{n}$ which is valid for $\|\lambda x\| < 1$. By the convergence of the above Taylor series, we see that 
\[
\lim\limits _{n \rightarrow \infty} \| \lambda^{n} x ^{n} \| = 0 \text{ if } | \lambda | < \alpha^{-1}.
\]

Thus $\|x^n \| = |\lambda|^{-n}\|\lambda^n x^n \| < |\lambda|^{-n}$ for large $n$ when $|\lambda| < \alpha^{-1}$, and so $\varlimsup\limits_{n \to \infty} \|x^n \|^{1/n} \leqq |\lambda|^{-1}$ when $|\lambda|^{-1} > \alpha$, that is, $\varlimsup\limits_{n \to \infty} \|x^n \|^{1/n} \leqq \alpha$. 
\end{proof}

\begin{corollaryu} The radical $R = \bigcap\limits_{J \in \{J\}} J$ of $B$ coincides with the totality of the generalized nilpotent elements $x \in B$ which are defined by 
\begin{equation}\label{eq:XI.2.2}
\lim\limits _{n \rightarrow \infty} \| x ^{n} \| ^{1/n} = 0.
\end{equation}
\end{corollaryu}

\begin{definition}\label{def:XI.2.Definition.2} $A$ complex number $\lambda$ is said to belong to the spectrum of $x \in B$, if the inverse $(x - \lambda e)^{-1}$ does not exist in $B$. 
\end{definition}

If $\lambda$ belongs to the spectrum of $x$, then there exists a maximal ideal $J$ such that $(x - \lambda e) \in J$. Conversely, if $(x - \lambda e)$ belongs to a maximal ideal $J$, then the inverse $(x - \lambda e)^{-1}$ does not exist. Hence we obtain 

\begin{theorem}\label{thm:XI.2.Theorem.2} The spectrum of $x \in B$ coincides with the totality of values taken by the function $x(J)$ on the space $\{J\}$ consisting of all the maximal ideals $J$ of $B$. 
\end{theorem}

Application of Tychonov's theorem. We define, for any $J_0 \in \{J\}$, a fundamental system of neighbourhoods of $J_0$ by 
\begin{equation}\label{eq:XI.2.3}
\{J \in \{J \}; | x _{i} (J) - x _{i} (J _{0}) | < \varepsilon_{i} (i = 1, 2, \dots , n) \},
\end{equation}
where $\varepsilon_{i} > 0$, $n$ and $x_{i} \in B$ are arbitrary. Then $\{J\}$ becomes a topological space and each $x(J)$, $x \in B$, becomes a continuous function on $\{J\}$. We have only to verify that if $J_{0} \neq J_{1}$, then there exist a neighbourhood $V_{0}$ of $J_{0}$ and a neighbourhood $V_{1}$ of $J_{1}$ with empty intersection. This may be done as follows. Let $x_{0} \in J_{0}$ and $x_{0} \bar{\in} J_{1}$ so that $x_{0}(J_{0}) = 0$ and $x_{0}(J_{1}) = \alpha \neq 0$. Then $V_{0} = \{J \in \{J\}; |x_{0}(J)| < |\alpha| / 2\}$ and $V_{1} = \{J \in \{J\}; |x_{0}(J) - x_{0}(J_{1})| < |\alpha| / 2\}$ have an empty intersection. 

\begin{theorem}\label{thm:XI.2.Theorem.3} The space $\{J\}$, topologized as above, is a compact space. 
\end{theorem}

\begin{proof}
 We attach, to each $x \in B$, the compact set 
\[
K _{x} = \{z; | z | \leqq \| x \| \}
\]
of the complex $z$-plane. Then the topological product 
\[
S = \prod\limits_{x \in B} K _{x}
\]
is a compact space, by virtue of Tychonov's theorem. See Chapter~\ref{chap:0}. To any maximal ideal $J_0 \in \{J\}$, we assign the point 
\[
\prod\limits_{x \in B} x (J _{0}) = s (J _{0}) \in S.
\]
$\{J\}$ is in one-one correspondence with a subset $S_{1}$ of $S$ by the above correspondence $J_{0} \rightarrow s(J_{0})$. Moreover, the topology of $\{J\}$ is the same as the relative topology of $S_{1}$ as a subset of $S$. Hence, if we can show that $S_{1}$ is a closed subset of the compact space $S$, then its topological image $\{J\}$ is compact. 

To prove that $S_{1}$ is a closed set, we consider an accumulation point $\omega = \prod\limits_{x \in B} \lambda_x \in S$ of the set $S_{1}$ in $S$. We shall show that the mapping $x \to \lambda_x$ is a homomorphism of the algebra $B$ into the complex number field $(K)$. Then $J_0 = \{x; \lambda_x = 0\}$ becomes, as may be seen from the isomorphism of $B / J_0$ with $(K)$, a maximal ideal of $B$ and $(x - \lambda_x e) \in J_0$, that is, $x(J_0) = \lambda_x$. This proves that the point $\omega = \prod\limits_{x \in B} \lambda_x = \prod\limits_{x \in B} x(J_0)$ belongs to $S_{1}$. 

We thus have to show that 
\[
\lambda_{x+y} = \lambda_{x} + \lambda_{y}, \lambda_{\alpha x} = \alpha \lambda_{x}, \lambda_{xy} = \lambda_{x} \lambda_{y}, \lambda_{e} = 1.
\]

We shall, for instance, prove that $\lambda_{x+y} = \lambda_x + \lambda_y$. Since $\omega = \prod\limits_{x \in B} \lambda_x$ is an accumulation point of $S_1$, there exists, for any $\varepsilon > 0$, a maximal ideal $J$ such that 
\begin{equation}\label{eq:XI.2.4}
\left| \lambda_{x} - x (J) \right| < \varepsilon , \left| \lambda_{y} - y (J) \right| < \varepsilon , \left| \lambda_{x+y} - (x + y) (J) \right| < \varepsilon .
\end{equation}

By $(x + y)(J) = x(J) + y(J)$ and the arbitrariness of $\varepsilon > 0$, we easily see that $\lambda_{x+y} = \lambda_x + \lambda_y$ is true. 
\end{proof}

We can now state the fundamental facts about the Gelfand representation $x \to x(J)$ of the normed ring $B$ in the form of 

\begin{theorem}\label{thm:XI.2.Theorem.4} $A$ normed ring $B$ over the complex number field is represented homomorphically by the ring of functions $x(J)$ on the compact space $\{J\}$ of all the maximal ideals $J$ of $B$. The radical $R$ of $B$ consists of those and only those elements which are represented by functions identically equal zero on $\{J\}$. The representation $x \to x(J)$ is isomorphic iff the ring $B$ is semi-simple. 
\end{theorem}

Application of the Stone-Weierstrass Theorem (of Chapter~\ref{chap:0}). The above obtained ring of functions is dense in the space of all complex-valued continuous functions on $\{J\}$ with uniform convergence topology if the ring $B$ is symmetric (or involutive) in the following sense: 

For any $x \in B$, there exists an $x^{*} \in B$ such that $x^{*}(J) = \overline{x\langle J\rangle}$ on $\{J\}$. ($4$) 

\yossub{Examples of Gelfand Representations}

\begin{example}\label{ex:XI.2.Example.1} Let $B = C(S)$ where $S$ is a compact topological space, and $J_0$ a maximal ideal of $C(S)$. Then there exists a point $s_0 \in S$ such that $x(s_0) = 0$ for all $x \in J_0$. Otherwise, for any $s_\alpha \in S$, there exists an $x_\alpha \in J_0$ such that $x_\alpha(s_\alpha) \neq 0$. $x_\alpha(s)$ being a continuous function, there exists a neighbourhood $V_\alpha$ of $s_\alpha$ such that $x_\alpha(s) \neq 0$ in $V_\alpha$. Since $S$ is compact, there exists a finite system, say, $V_{\alpha_1}, V_{\alpha_2}, \ldots, V_{\alpha_n}$ such that $\bigcup\limits_{j=1}^{n} V_{\alpha_j} = S$. Hence the function 
\[
x (s) = \sum\limits_{i=1} ^{n} \overline {{{{x _ {\alpha_{i}} (s)}}}} x _ {\alpha_{i}} (s) \in J _{0}
\]
does not vanish over $S$, and the inverse $x^{-1}$, $x^{-1}(s) = x(s)^{-1}$, of $x \in J_0$ exists, contrary to the maximality of the ideal $J_0$. Thus we see that $J_0$ is contained in the maximal ideal $J' = \{x \in B; x(s_0) = 0\}$. By the maximality of $J_0$, we must have $J_0 = J'$. In this way we see that the space $\{J\}$ of the maximal ideals $J$ of $B$ is in one-one correspondence with the points $s$ of $S$. 
\end{example}

\begin{example}\label{ex:XI.2.Example.2} Let $B$ be the totality of functions $x(s)$, $0 \leqq s \leqq 1$, which can be represented by absolutely convergent Fourier series: 
\[
x (s) = \sum\limits_{n = - \infty} ^{\infty} c _{n} e ^{2 \pi i s n}, \sum\limits_{n = - \infty} ^{\infty} | c _{n} | < \infty .
\]
\end{example}

$B$ is a normed ring by $(x + y)(s) = x(s) + y(s), (xy)(s) = x(s)y(s)$ and $\|x\| = \sum\limits_{j} |c_j|$. Let $J_0$ be a maximal ideal of $B$. Set $e^{2\pi is} = x_1$. Then $x_1^{-1} = e^{-2\pi is}$, and so, by $|x_1(J_0)| \leqq \|x_1\| = 1, |x_1^{-1}(J_0)| = |x_1(J_0)^{-1}| \leqq \|x_1^{-1}\| = 1$, we see that $|x_1(J_0)| = 1$. Hence there exists a point $s_0, 0 \leqq s_0 \leqq 1$, such that $x_1(J_0) = e^{2\pi is_0}$. Thus $x_n = e^{2\pi isn} = x_1^n$ satisfies $x_n(J_0) = e^{2\pi is_n n}$, and so $x(J_0) = \sum\limits_{n=-\infty}^{\infty} c_n e^{2\pi is_n n} = x(s_0)$. In this way, we see that for any maximal ideal $J_0$ of $B$, there exists a point $s_0, 0 \leqq s_0 \leqq 1$, such that the homomorphism $x \to x(J_0)$ is given by $x(J_0) = x(s_0)$, for all $x \in B$. It is also clear that the mapping $x \to x(s_0)$ gives a homomorphism of the algebra $B$ into the complex number field. Therefore we see that the maximal ideal space of $B$ coincides with $\{e^{2\pi is}; 0 \leqq s \leqq 1\}$. 

\begin{corollaryn}[N. \textsc{Wiener}'s theorem] If an absolutely convergent Fourier series $x(s) = \sum\limits_{n=-\infty}^{\infty} c_n e^{2\pi isn}$ does not vanish on $[0,1]$, then the function $1/x(s)$ is also representable as an absolutely convergent Fourier series. For, $x$ does not belong to any of the maximal ideals of the normed ring of the above Example~\hyperref[ex:XI.2.Example.2]{$2$}.

\end{corollaryn}
\begin{example}\label{ex:XI.2.Example.3} We take $B_{1} = C[0,1]$, and define, for $x, y \in B_{1}$, 
\[
\begin{aligned}
(x + y) (s) &= x (s) + y (s), (\alpha x) (s) = \alpha x (s), (x y) (s) =
\\
\int_{0} ^{s} x (s - t) y (t) \dd t \text{ and } \| x \| &= \sup\limits _{s \in [ 0, 1 ]} | x (s) |.
\end{aligned}
\]
$B_{1}$ is a commutative $B$-algebra without unit. Adjoining formally a unit $e$ by the rule $ex = xe = x$, $\|e\| = 1$, the set $B = \{z = \lambda e + x; x \in B_{1}\}$ becomes a normed ring by the operations 
\[
\begin{yoslines}
&(\lambda_{1} e + x _{1}) + (\lambda_{2} e + x _{2}) = (\lambda_{1} + \lambda_{2}) e + (x _{1} + x _{2}), \alpha (\lambda e + x) \\
&= \alpha \lambda e + \alpha x, (\lambda_{1} e + x _{1}) (\lambda_{2} e + x _{2}) = \lambda_{1} \lambda_{2} e + \lambda_{1} x _{2} + \lambda_{2} x _{1} + x _{1} x _{2} \text{ and } \| \lambda e + x \| = | \lambda | + \| x \|.
\end{yoslines}
\]
\end{example}

We have, by induction, 
\[
\left| x ^{2} (s) \right| \leqq M ^{2} s, \left| x ^{3} (s) \right| \leqq M ^{3} \frac {s ^{2}}{2}, \dots , \left| x ^{n} (s) \right| \leqq M ^{n} \frac {s ^{n-1}}{(n - 1) !}, \dots
\]
where $M = \sup\limits_{s \in S} |x(s)| = \|x\|$. Thus every $x \in B_1$ is a generalized nilpotent element of $B$, due to the fact that $\lim\limits_{n \to \infty} (n!)^{1/n} = \infty$. 


\yossection{3}{The Spectral Resolution of Bounded Normal Operators}
\label{sec:XI.3}

Let $X$ be a Hilbert space, and let a system $M$ of bounded normal operators $\in L(X, X)$ satisfy the conditions: 
\begin{align}
T, S& \in M \text{ implies } T S = S T (\text{ commutativity }), \label{eq:XI.3.1} \\
T& \in M \text{ implies } T ^{*} \in M. \label{eq:XI.3.2}
\end{align}

A system $M$ consisting of a bounded normal operator $T \in L(X, X)$ and its adjoint $T^{*}$ surely satisfies \eqref{eq:XI.3.1} and \eqref{eq:XI.3.2}. 

Let $M'$ be the totality of operators $\in L(X, X)$ which commute with every $T \in M$, and let $B = M'' = (M')'$ be the totality of operators $\in L(X, X)$ commutative with every operator $S \in M'$. 

\begin{proposition}\label{prop:XI.3.Proposition.1} Every element of $B$ is a normal operator. $B$ is a normed ring over the complex number field by the operator sum, the operator product, the unit $I$ (the identity operator) and the operator norm $\|T\|$. 
\end{proposition}

\begin{proof}
 $M \subseteq M'$ by \eqref{eq:XI.3.1}, and so $M' \supseteq M''$. Hence $M''' = (M'')' \supseteq M''$ and so $B = M''$ is a commutative ring. The identity operator $I$ belongs to $B$ and is the unit of this algebra $B$. By \eqref{eq:XI.3.2}, we easily see that every operator $\in B$ is normal. Since the multiplication $TS$ and the adjoint formation $T \to T^*$ in the algebra $B$ are continuous with respect to the norm of the operator, it is easy to see that the ring $B$ is complete with respect to the operator norm. 
\end{proof}

\begin{theorem}\label{thm:XI.3.Theorem.1} By the Gelfand representation 
\begin{equation}\label{eq:XI.3.3}
B \ni T \rightarrow T (J),
\end{equation}
the ring $B$ is represented isomorphically by the algebra $C(\{J\})$ of all continuous functions $T(J)$ on the compact space $\{J\}$ of all the maximal 
ideals $J$ of $B$ in such a way that 
\begin{equation}\label{eq:XI.3.4}
\|T\|=\sup_{J\in\{J\}}|T(J)|,
\end{equation}
\begin{equation}\label{eq:XI.3.5}
\begin{minipage}{0.92\linewidth}
\everymath{}\let\limits\nolimits
$T(J)$ is real-valued on $\{J\}$ iff $T$ is self-adjoint,
\end{minipage}
\end{equation}
\begin{equation}\label{eq:XI.3.6}
\begin{minipage}{0.92\linewidth}
\everymath{}\let\limits\nolimits
$T(J)\geqq0$ on $\{J\}$ iff $T$ is self-adjoint and positive, that is, $(Tx,x)\geqq0$ for all $x\in X$.
\end{minipage}
\end{equation}

\end{theorem}

\begin{proof}
 We first show that, for any bounded normal operator $T$, 
\begin{equation}\label{eq:XI.3.7}
\left\| T ^{2} \right\| = \left\| T \right\| ^{2}. \end{equation}

By the normality of $T$, we see that $H = TT^{*} = T^{*}T$ is self-adjoint. Hence, by Theorem~\ref{thm:VII.3.Theorem.3} in Chapter~\ref{chap:VII}, \ref{sec:VII.3}, 
\[
\left\| T \right\| ^{2} = \sup\limits _{\| x \| \leqq 1} (T x, T x) = \sup\limits _{\| x \| \leqq 1} \left| \left(T ^{*} T x, x\right) \right| = \| H \| = \| T ^{*} T \| = \| T T ^{*} \|.
\]

Since $(T^{*})^{2} = (T^{2})^{*}$, $T^{2}$ is normal with $T$. Thus, as above, we obtain $\| T^2\|^2 = \| T^{*2}T^2\|$, which is, by the commutativity $TT^{*} = T^{*}T$, equal to $\| (T^{*}T)^{2}\| = \| H^{2}\|$. Since $H^{2}$ is self-adjoint, we obtain, again by Theorem~\ref{thm:VII.3.Theorem.3} in Chapter~\ref{chap:VII}, \ref{sec:VII.3}, 
\[
\left\| H \right\| ^{2} = \sup\limits _{\| x \| \leqq 1} (H x, H x) = \sup\limits _{\| x \| \leqq 1} \left| \left(H ^{2} x, x\right) \right| = \left\| H ^{2} \right\|.
\]

Therefore, $\|T^2\|^2 = \|H^2\| = \|H\|^2 = (\|T\|^2)^2$, that is, $\|T^2\| = \|T\|^2$. We have, by \eqref{eq:XI.3.7}, $\|T\| = \lim\limits_{n\to \infty} \|T^n\|^{1/n}$, because we know already that the right hand limit exists (see \eqref{eq:VIII.2.3} in Chapter~\ref{chap:VIII}, \ref{sec:VIII.2}). Hence, by Theorem~\hyperref[thm:XI.2.Theorem.4]{$4$} of the preceding section, the representation \eqref{eq:XI.3.3} is isomorphic and \eqref{eq:XI.3.4} is true. 
\end{proof}

\begin{proof}[Proof of \eqref{eq:XI.3.5}.]
Let a self-adjoint $T \in B$ satisfy, for a certain $J_0 \in \{J\}$, $T(J_0) = a + ib$ with $b \neq 0$. Then the self-adjoint operator $S = (T - aI)/b \in B$ satisfies $(I + S^2)(J_0) = 1 + i^2 = 0$, and so $(I + S^2)$ does not have an inverse in $B$. But, by Theorem~\ref{thm:VII.3.Theorem.2} in Chapter~\ref{chap:VII}, \ref{sec:VII.3}, $(I + S^2)$ has an inverse which surely belongs to $B$. Thus, if $T \in B$ is self-adjoint, $T(J)$ must be real-valued. Let $T \in B$ be not self-adjoint, and put
\[
T = \frac {T + T ^{*}}{2} + i \frac {T - T ^{*}}{2 i}.
\]

Then, since the first term on the right is self-adjoint, the self-adjoint operator $(T - T^{*}) / 2i$ must be $\neq 0$. Thus, by the isomorphism of representation \eqref{eq:XI.3.3}, there must exist a $J_0 \in \{J\}$ such that $\frac{T - T^*}{2i}(J_0) \neq 0$. Hence $T(J_0) = \frac{T + T^*}{2}(J_0) + i\frac{T - T^*}{2i}(J_0)$ is not real. For, as proved above, self-adjoint operators are represented by real-valued functions. 
\end{proof}

\begin{proof}[Proof of \eqref{eq:XI.3.6}.]
We first show that
\begin{equation}\label{eq:XI.3.8}
T ^{*} (J) = \overline {{{{T (J)}}}} \text{ on } \{J \}.
\end{equation}

This is clear, since self-adjoint operators $(T + T^{*})/2$ and $(T - T^{*})/2i$ are represented by real-valued functions. Therefore, by \eqref{eq:XI.3.4} and the result of the preceding section, the ring $B$ is represented by the ring of all continuous complex-valued functions on $\{J\}$ satisfying \eqref{eq:XI.3.5} and \eqref{eq:XI.3.8}. Let $T(J) \geqq 0$ over $\{J\}$. Then $S(J) = T(J)^{1/2}$ is a continuous function on $\{J\}$. Hence, by the isomorphism of the representation \eqref{eq:XI.3.3}, $S^{2} = T$. By \eqref{eq:XI.3.5}, we have $S = S^{*}$. Hence $(Tx, x) = (S^{2}x, x) = (Sx, Sx) \geqq 0$. To prove, conversely, that the condition $(Tx, x) \geqq 0$ for all $x \in X$ implies $T(J) \geqq 0$ over $\{J\}$, we set $T_{1}(J) = \max(T(J), 0)$ and $T_{2}(J) = T_{1}(J) - T(J)$. Then, by what we have proved above, $T_{1}$ and $T_{2}$ are both $\in B$, self-adjoint and positive: $(T_{j}x, x) \geqq 0$ for all $x \in X$ ($j = 1$, $2$). Moreover, we have $T_{1}T_{2} = 0$ and $T_{2} = T_{1} - T$. The former equality is implied by $T_{1}(J)T_{2}(J) = 0$. 

Therefore, we have 
$0 \leqq (TT_{2}x, T_{2}x) = (-T_{2}^{2}x, T_{2}x) = -(T_{2}^{3}x, x) = -(T_{2}T_{2}x, T_{2}x) \leqq 0$. Thus $(T_{2}^{3}x, x) = 0$ and so, by Theorem~\ref{thm:VII.3.Theorem.3} in Chapter~\ref{chap:VII}, \ref{sec:VII.3}, we must have $T_{2}^{3} = 0$. Hence, by $\|T_{2}\| = \lim\limits_{n \to \infty} \|T_{2}^{n}\|^{1/n}$, we obtain $T_{2} = 0$. We have thus proved $T = T_{1}$ and hence $T(J) \geqq 0$ on $\{J\}$. 

We shall thus write $T \geqq 0$ if $T$ is self-adjoint and positive. We also write $S \geqq T$ if $(S - T) \geqq 0$. 
\end{proof}

\begin{theorem}\label{thm:XI.3.Theorem.2} Let $\{T_n\} \subseteq B$ be a sequence of self-adjoint operators such that 
\begin{equation}\label{eq:XI.3.9}
0 \leqq T _{1} \leqq T _{2} \leqq \dots \leqq T _{n} \leqq \dots \leqq S \in B.
\end{equation}
\end{theorem}

Then, for any $x \in X$, $\slim\limits_{n \to \infty} T_n x = Tx$ exists, i.e., $\slim\limits_{n \to \infty} T_n = T$ exists and $T \in B$, $S \geqq T \geqq T_n$ ( $n = 1, 2, \ldots$). 

\begin{proof}
 We first remark that, by \eqref{eq:XI.3.6}, 
\begin{equation}\label{eq:XI.3.10}
E, F \in B \text{ and } E \geqq 0, F \geqq 0 \text{ imply } E + F \geqq 0 \text{ and } EF \geqq 0.
\end{equation}
Thus $0 \leqq T_1^2 \leqq T_2^2 \leqq \cdots \leqq T_n^2 \leqq \cdots \leqq S^2$. Hence, for any $x \in X$, a finite $\lim\limits_{n \to \infty} (T_n^2 x, x)$ exists. Since, by \eqref{eq:XI.3.6}, $T_{n+k}^2 \geqq T_{n+k} T_n \geqq T_n^2$, we also have 
\[
\lim\limits _{n, k \rightarrow \infty} (T _{n+k} ^{2} x, x) = \lim\limits _{n, k \rightarrow \infty} (T _{n+k} T _{n} x, x) = \lim\limits _{n \rightarrow \infty} (T _{n} ^{2} x, x). \]

Therefore, $\lim\limits_{n,m\to \infty}((T_n - T_m)^2 x,x) = \lim\limits_{n,m\to \infty}\|T_nx - T_mx\|^2 = 0$ so that $s\text{- }\lim\limits_{n\to \infty}T_nx = Tx$ exists. That $T\in B$ and $S\geqq T\geqq T_n$ is clear from the process of the proof. 
\end{proof}

\begin{theorem}\label{thm:XI.3.Theorem.3} Let a sequence of real-valued functions $\{T_n(J)\}$, where $T_n \in B$, satisfy the condition 
\begin{equation}\label{eq:XI.3.11}
0 \leqq T_{1}(J) \leqq T_{2}(J) \leqq \cdots \leqq T_{n}(J) \leqq \cdots \leqq a \text{ finite constant on } \{J\}.
\end{equation}
\end{theorem}

Then, by \eqref{eq:XI.3.6} and Theorem~\ref{thm:XI.3.Theorem.2}, $\slim\limits_{n \to \infty} T_n = T$ exists. In such a case, we can prove that 
\[
D = \{J \in \{J \}; T (J) \neq \lim\limits _{n \rightarrow \infty} T _{n} (J) \}
\]
is a set of the first category, and so $D^{C} = \{J\} - D$ is dense in $\{J\}$. 

\begin{proof}
 By Theorem~\ref{thm:XI.3.Theorem.2}, $T \geqq T_n$ and so $T(J) \geqq \lim\limits_{n \to \infty} T_n(J)$ on $\{J\}$. By Baire's theorem in Chapter~\ref{chap:0}, \ref{sec:0.2}, the set of points of discontinuity of the function $\lim\limits_{n \to \infty} T_n(J)$ is of the first category. Hence, if the set $D$ is not of the first category, then there exists at least one point $J_0 \in D$ at which $\lim\limits_{n \to \infty} T_n(J)$ is continuous. In other words, there exists a positive number $\delta$ and an open set $V(J_0) \ni J_0$ of $\{J\}$ such that 
\[
T (J) \geqq \delta + \lim\limits _{n \rightarrow \infty} T _{n} (J) \text{ whenever } J \in V (J _{0}). \]

Since the compact space $\{J\}$ is normal, and since $T(J) \geqq \lim\limits_{n \to \infty} T_n(J)$ on $\{J\}$, we may construct, by Urysohn's theorem an open set $V_1(J_0) \ni J_0$ and a function $W(J) \in C(\{J\})$ such that $0 \leqq W(J) \leqq \delta$ on $\{J\}$, $V_1(J_0)^a \subseteq V(J_0)$, $W(J) = \delta/2$ on $V_1(J_0)$ and $W(J) = 0$ on $V(J_0)^C$. Hence $T(J) - W(J) \geqq \lim\limits_{n \to \infty} T_n(J)$ on $\{J\}$, and so, by \eqref{eq:XI.3.6}, $T - W \geqq T_n$ ( $n = 1, 2, \ldots$). $W(J) \not\equiv 0$ implies, by the isomorphism \eqref{eq:XI.3.3}, that $W \neq 0$, $W \geqq 0$. Thus, again by \eqref{eq:XI.3.6}, $T - W \geqq s\text{- }\lim\limits_{n \to \infty} T_n$, contrary to $T = s\text{- }\lim\limits_{n \to \infty} T_n$. Finally, since $\{J\}$ is a compact space, the complement $D^C = \{J\} - D$ of the set $D$ of the first category must be dense in $\{J\}$. 
\end{proof}

We are now able to prove (K. \textsc{Yosida} \cite{yosida355}) 

\yossub{The spectral resolution or the spectral representation of operators $\in B$.}

Consider the set $C'(\{J\})$ of all complex-valued bounded functions $T'(J)$ on $\{J\}$ such that $T'(J)$ is different from a continuous function $T(J)$ only on a set of the first category. We identify two functions from $C'(\{J\})$ if they differ only on a set of the first category. Then $C'(\{J\})$ is divided into classes. Since the complement of a set of the first category is dense in the compact space $\{J\}$, each class $T'$ contains exactly one continuous function $T(J)$ which corresponds, by the isomorphism $B \leftrightarrow C(\{J\})$, to an element $T \in B$. 

For any $T \in B$ and for any complex number $z = \lambda + i\mu$, we put $E_z =$ the element $\in B$ which corresponds to the class $E_z'$ containing the defining function $E_z'(J)$ of the set $\{J \in \{J\}; \operatorname{Re} T(J) < \lambda, \operatorname{Im} T(J) < \mu\}$. It is clear that there exists a monotone increasing sequence of continuous functions $f_n$ of complex argument such that $E_z'(J) = \lim\limits_{n \to \infty} f_n(T(J))$ and 
so $E_{z}^{\prime}(J)\in C^{\prime}(\{J\})$. We have then 
\[
\begin{gathered}
| T (J) - \sum\limits_{j=2} ^{n} (\lambda_{j} + i \mu_{j}) \left(E _ {\lambda_{j} + i \mu_{j}} ^{\prime} (J) + E _ {\lambda_{j-1} + i \mu_{j-1}} ^{\prime} (J) - E _ {\lambda_{j-1} + i \mu_{j}} ^{\prime} (J) \right.
\\
\left. - E _ {\lambda_{j} + i \mu_{j-1}} ^{\prime} (J)\right) | \leqq \varepsilon \text{ on } \{J \} \text{ if }
\end{gathered}
\]
\[
\begin{gathered}
\lambda_{1} = - \alpha - \frac {\varepsilon}{\sqrt {2}} \leqq \lambda_{2} \leqq \dots \leqq \lambda_{n} = \alpha = \sup\limits _ {J \in \{J \}} | \operatorname{Re} T (J) |,
\\
\mu_{1} = - \beta - \frac {\varepsilon}{\sqrt {2}} \leqq \mu_{2} \leqq \dots \leqq \mu_{n} = \beta = \sup\limits _ {J \in \{J \}} | \operatorname{Im} T (J) |,
\\
\left(\sup\limits _{j} (\lambda_{j} - \lambda_{j-1}) ^{2} + \sup\limits _{j} (\mu_{j} - \mu_{j-1}) ^{2}\right) ^{1/2} \leqq \varepsilon .
\end{gathered}
\]

Thus, by the definition of $E_{z}(J)$, we have 
\[
\begin{gathered}
\left| T (J) - \sum\limits_{j=2} ^{n} (\lambda_{j} + i \mu_{j}) \left(E _ {\lambda_{j} + i \mu_{j}} (J) + E _ {\lambda_{j-1} + i \mu_{j-1}} (J) \right. \right. \left. \left. - E _ {\lambda_{j-1} + i \mu_{j}} (J) - E _ {\lambda_{j} + i \mu_{j-1}} (J)\right) \right| \leqq \varepsilon
\end{gathered}
\]
on $\{J\}$, since the complement of a set of the first category is dense in the compact space $\{J\}$. Therefore, by \eqref{eq:XI.3.4}, we have 
\[
\left\| T - \sum\limits_{j=2} ^{n} \left(\lambda_{j} + i \mu_{j}\right) \left(E _ {\lambda_{j} + i \mu_{j}} + E _ {\lambda_{j-1} + i \mu_{j-1}} - E _ {\lambda_{j-1} + i \mu_{j}} - E _ {\lambda_{j} + i \mu_{j-1}}\right) \right\| \leqq \varepsilon ,
\]
which we shall write as 
\begin{equation}\label{eq:XI.3.12}
T = \iint z \dd E _{z},
\end{equation}
and it is called the spectral resolution of the normal operator $T$. 


\yossection{4}{The Spectral Resolution of a Unitary Operator}
\label{sec:XI.4}

If $T$ is a unitary operator $\in B$, then, by 
\begin{equation}\label{eq:XI.4.1}
T (J) T ^{*} (J) = T (J) \overline {{T (J)}} = 1,
\end{equation}
we see that the values taken by the function $T(J)$ on $\{J\}$ are complex numbers of absolute value $1$. From this fact, we can simplify the spectral resolution $\iint z dE_z$ of $T$. 

The defining function $E_{\theta}'(J)$ of the set $\{J \in \{J\}; \arg(T(J)) \in (0, \theta)\}$, $0 < \theta < 2\pi$, belongs to $C'(\{J\})$, and we have, by setting $E_0'(J) = 0$, $E_{2\pi}'(J) = 1$, 
\[
\begin{gathered}
\left| T (J) - \sum\limits_{j=1} ^{n} e ^ {i \theta_{j}} (E _ {\theta_{j}} ^{\prime} (J) - E _ {\theta_{j-1}} ^{\prime} (J)) \right| \leqq \max\limits _{j} \left| e ^ {i \theta_{j}} - e ^ {i \theta_{j-1}} \right| (0 = \theta_{0} < \theta_{1} < \dots < \theta_{n} = 2 \pi).
\end{gathered}
\]

Let $E_{\theta}(J)$ be the continuous function on $\{J\}$ which is different from $E_{\theta}'(J)$ only on a set of the first category, and let $E_{\theta}$ be the operator $\in B$ which corresponds to $E_{\theta}(J)$ by the isomorphic representation $B \ni T \leftrightarrow T(J)$. Then, as in the preceding section, 
\[
\left\| T - \sum\limits_{j=1} ^ {\dot {n}} e ^ {i \theta_{j}} (E _ {\theta_{j}} - E _ {\theta_{j-1}}) \right\| \leqq \max\limits _{j} \left| e ^ {i \theta_{j}} - e ^ {i \theta_{j-1}} \right|.
\]

This we write, in view of the fact $e^{2\pi i} = 1$, 
\begin{equation}T = \int_{0} ^{2 \pi} e ^{i \theta} \dd F (\theta), \text{ where } F (\theta) = E _{\theta +0} - E _{+0} \text{ for } 0 < \theta < 2 \pi , F (0) = 0, F (2 \pi) = I. \label{eq:XI.4.2}\end{equation}

Here $E_{\theta +0}$ is defined by $E_{\theta +0}x = \slim\limits_{\theta \downarrow 0} E_{\theta}x$, the existence of this limit will be proved below. 

\begin{theorem}\label{thm:XI.4.Theorem.1} The system of operators $F(\theta), 0 \leqq \theta \leqq 2\pi$, satisfies the conditions: 
\begin{equation}\label{eq:XI.4.3}
\parbox{0.90\linewidth}{each $F(\theta)$ is a projection operator, commutative with every bounded linear operator commutative with $T$,}
\end{equation}
\end{theorem}

\begin{gather}F (\theta) F \left(\theta^{\prime}\right) = F \left(\min \left(\theta , \theta^{\prime}\right)\right), \label{eq:XI.4.4} \\F (0) = 0, F (2 \pi) = I, \label{eq:XI.4.5} \\F (\theta + 0) = F (\theta), \quad 0 \leqq \theta < 2 \pi , \text{ in the sense that } \slim _ {\theta^{\prime} \downarrow \theta} F (\theta^{\prime}) x = F (\theta) x \text{ for every } x \in X. \label{eq:XI.4.6}
\end{gather}

\begin{proof}
 It is sufficient to prove that $E_{\theta}, 0 \leqq \theta \leqq 2\pi$, satisfies the conditions: 
\begin{align}
\text{ each } E _{\theta} \text{ is a projection operator }& \in B , \tag{3'}\label{eq:XI.4.3'} \\
E _{\theta} E _ {\theta^{\prime}}& = E _ {\min (\theta , \theta^{\prime})}, \tag{4'}\label{eq:XI.4.4'} \\
E _{0}& = 0, E _{2 \pi} = I, \tag{5'}\label{eq:XI.4.5'} \\
E _{\theta} x& = \slim _ {\theta^{\prime} \downarrow \theta} E _ {\theta^{\prime}} x \text{ for every } x \in X \text{ and } 0 \leqq \theta < 2 \pi .  \tag{6'}\label{eq:XI.4.6'} \end{align}

We have $E_{\theta}^{\prime}(J) = \overline{E_{\theta}^{\prime}(J)}$ and $E_{\theta}^{\prime}(J)^{2} = E_{\theta}^{\prime}(J)$. Hence, by the result of the preceding section, we obtain $E_{\theta} = E_{\theta}^{*}$ and $E_{\theta}^{2} = E_{\theta}$. This proves ($3$'). \eqref{eq:XI.4.4} is proved similarly from $E_{\theta}^{\prime}(J)E_{\theta^{\prime}}^{\prime}(J) = E_{\min (\theta ,\theta^{\prime})}^{\prime}(J)$, and \eqref{eq:XI.4.5} is proved similarly. Next let $\theta_{n}\downarrow \theta$. Then $E_{\theta_n}^{\prime}(J)\geqq E_{\theta_{n+1}}^{\prime}(J)\geqq E_{\theta}^{\prime}(J)$, and so, by the result in the preceding section, $s\text{- }\lim\limits_{n\to \infty}E_{\theta_n} = E$ exists and $E(J) = E_{\theta}^{\prime}(J) = \lim\limits_{\theta_{n}\downarrow \theta}E_{\theta_{n}}^{\prime}(J)$ on $\{J\}$ except possibly on a set of the first category. Thus $E = E_{\theta}$. 
\end{proof}

\begin{example}\label{ex:XI.4.Example.1} Let a linear operator $T$ defined by 
\[
T y (s) = e ^{is} y (s), \text{ where } y (s) \in L ^{2} (- \infty , \infty).
\]
$T$ is unitary. We define, when $2n\pi < s \leqq 2(n + 1)\pi$, 
\begin{equation}F (\theta) y (s) = y (s) \quad \text{ for } \quad s \leqq \theta + 2 n \pi \leqq 2 (n + 1) \pi , F (\theta) y (s) = 0 \quad \text{ for } \quad \theta + 2 n \pi < s. \notag\end{equation}
\end{example}

It is easy to see that $T = \int_{0}^{2\pi} e^{i\theta} dF(\theta)$. 


\begin{example}\label{ex:XI.4.Example.2} Let a linear operator $T_{1}$ be defined by 
\[
T _{1} x (t) = x (t + 1) \quad \text{ in } \quad L ^{2} (- \infty , \infty).
\]
$T_{1}$ is unitary. By the Fourier transformation 
\[
y (s) = U x (t) = \underset {n \rightarrow \infty} {\text{ l . i . m . }} (2 \pi) ^{-1/2} \int_{-n} ^{n} e ^{-ist} x (t) \dd t,
\]
we obtain $Ux(t + 1) = e^{is}Ux(t) = e^{is}y(s)$. Thus 
\[
T _{1} x (t) = x (t + 1) = U ^{-1} e ^{is} y (s) = U ^{-1} T y (s) = U ^{-1} T U x (t),
\]
that is, $T_{1} = U^{-1}TU$. Therefore, we have 
\[
T _{1} = \int_{0} ^{2 \pi} e ^{i \theta} \dd F _{1} (\theta), \text{ where } F _{1} (\theta) = U ^{-1} F (\theta) U.
\]
\end{example}

The uniqueness of the spectral resolution. Since $T^{-1} = T^{*}$ and $T^{-1}(J) = T^{*}(J) = T(J)^{-1}$, we easily see that 
\begin{equation}\label{eq:XI.4.7}
T ^{-1} = \int_{0} ^{2 \pi} e ^{- i \theta} \dd F (\theta).
\end{equation}

Let $\max\limits_{j}|e^{i\theta_j} - e^{i\theta_{j-1}}| < \varepsilon$. Then, from 
\[
T = \sum\limits_{j} e ^ {i \theta_{j}} (F (\theta_{j}) - F (\theta_{j-1})) + \delta , \| \delta \| < \varepsilon ,
\]
we obtain, by \eqref{eq:XI.4.4}, 
\[
\begin{gathered}
T ^{2} = \sum\limits_{j} e ^ {2 i \theta_{j}} (F (\theta_{j}) - F (\theta_{j-1})) + \delta^{\prime}, \text{ where } \| \delta^{\prime} \| \leqq \| (T - \delta) \delta \| + \| \delta (T - \delta) \| + \| \delta^{2} \|
\\
\leqq (\| T \| + \varepsilon) \varepsilon + \varepsilon (\| T \| + \varepsilon) + \varepsilon^{2}.
\end{gathered}
\]

Hence we have $T^2 = \int_0^{2\pi} e^{2i\theta} dF(\theta)$, and, more generally, 
\begin{equation}\label{eq:XI.4.8}
T ^{n} = \int_{0} ^{2 \pi} e ^{i n \theta} \dd F (\theta) \quad (n = 0, \pm 1, \pm 2, \dots).
\end{equation}

Therefore, if there exists another spectral resolution $T = \int_{0}^{2\pi} e^{i\theta} dF_1(\theta)$ satisfying ($3$) to \eqref{eq:XI.4.6}, then, for any polynomial $p(\theta)$ in $e^{i\theta}$ and $e^{-i\theta}$, 
\[
\int_{0} ^{2 \pi} p (\theta) d ((F (\theta) x, y) - (F _{1} (\theta) x, y)) = 0, (x, y \in X).
\]

Thus, by continuity, the above equality holds for any continuous function $p(\theta)$ with $p(0) = p(2\pi)$. Let $0 < \theta_{0} < \theta_{1} < 2\pi$, and take 
\[
\begin{aligned}
&p _{n} (\theta)  = 0 \text{ for } 0 \leqq \theta \leqq \theta_{0} \text{ and for } \theta_{1} + \frac {1}{n} \leqq \theta \leqq 2 \pi , = 1 \text{ for } \theta_{0} + \frac {1}{n} \leqq \theta \leqq \theta_{1}, \\
&= \text{ linear for } \theta_{0} \leqq \theta \leqq \theta_{0} + \frac {1}{n} \text{ and for } \theta_{1} \leqq \theta \leqq \theta_{1} + \frac {1}{n}.
\end{aligned}
\]

Then letting $n \to \infty$, we obtain, by \eqref{eq:XI.4.6}, 
\[
\lim\limits _{n \rightarrow \infty} \int_{0} ^{2 \pi} p _{n} (\theta) d [ (F (\theta) x, y) - (F _{1} (\theta) x, y) ] = [ (F (\theta) x, y) - (F _{1} (\theta) x, y) ] _ {\theta_{0}} ^ {\theta_{1}} = 0,
\]
which is valid for all $x, y \in X$. Hence, letting $\theta_{0} \downarrow 0$ and making use of conditions \eqref{eq:XI.4.5} and \eqref{eq:XI.4.6}, we see that $F(\theta_{1}) = F_{1}(\theta)$. Therefore, the spectral resolution for a unitary operator is uniquely determined. 


\yossection{5}{The Resolution of the Identity}
\label{sec:XI.5}

\begin{definition}\label{def:XI.5.Definition.1} $A$ family of projections $E(\lambda), -\infty < \lambda < \infty$, in a Hilbert space $X$ is called a (real) resolution of the identity if it satisfies the conditions: 
\begin{align}
E (\lambda) E (\mu)& = E (\min (\lambda , \mu)), \label{eq:XI.5.1} \\
E (-& \infty) = 0, E (+ \infty) = I, \text{ where } E (- \infty) x = \slim _{\lambda \downarrow - \infty} E (\lambda) x \text{ and } \notag \\
E (+& \infty) x = \slim _{\lambda \uparrow \infty} E (\lambda) x, \label{eq:XI.5.2} \\
E (\lambda + 0)& = E (\lambda), \text{ where } E (\lambda + 0) x = s \underset {\mu \downarrow \lambda} {- \lim} E (\mu) x. \label{eq:XI.5.3}
\end{align}
\end{definition}

\begin{proposition}\label{prop:XI.5.Proposition.1} For any $x, y \in X$, the function $(E(\lambda)x, y)$ is, as a function of $\lambda$, of bounded variation. 
\end{proposition}

\begin{proof}
 Let $\lambda_1 < \lambda_2 < \cdots < \lambda_n$. Then, by \eqref{eq:XI.5.1}, $E(\alpha, \beta] = E(\beta) - E(\alpha)$ is a projection. Thus we have, by Schwarz' inequality, 
\[
\begin{aligned}
&\sum\limits_{j} | (E (\lambda_{j-1}, \lambda_{j} ] x, y) | = \sum\limits_{j} | (E (\lambda_{j-1}, \lambda_{j} ] x, E (\lambda_{j-1}, \lambda_{j} ] y) |
\\
&\quad \leqq \sum\limits_{j} \| E (\lambda_{j-1}, \lambda_{j} ] x \| \cdot \| E (\lambda_{j-1}, \lambda_{j} ] y \| \leqq \left(\sum\limits_{j} \| E (\lambda_{j-1}, \lambda_{j} ] x \| ^{2}\right) ^{1/2}
\\
&\quad \cdot \left(\sum\limits_{j} \| E (\lambda_{j-1}, \lambda_{j} ] y \| ^{2}\right) ^{1/2} = (\| E (\lambda_{1}, \lambda_{n} ] x \| ^{2}) ^{1/2} \cdot (\| E (\lambda_{1}, \lambda_{n} ] y \| ^{2}) ^{1/2} \leqq \| x \| \cdot \| y \|.
\end{aligned}
\]

For, by the orthogonality 
\begin{equation}\label{eq:XI.5.4}
E \left(\lambda_{j-1}, \lambda_{j} \right] \cdot E \left(\lambda_{i-1}, \lambda_{i} \right] = 0 \quad (i \neq j)
\end{equation}
implied by \eqref{eq:XI.5.1}, we have, for $m > n$, 
\begin{equation}\label{eq:XI.5.5}
\left\| x \right\| ^{2} \geqq \left\| E \left(\lambda_{n}, \lambda_{m} \right] x \right\| ^{2} = \sum\limits_{i=n} ^{m-1} \left\| E \left(\lambda_{i}, \lambda_{i+1} \right] x \right\| ^{2}.
\end{equation}
\end{proof}

\begin{corollaryu} For any $\lambda, -\infty < \lambda < \infty$, the operators $E(\lambda + 0) = s\text{- }\lim\limits_{\lambda' \downarrow \lambda} E(\lambda')$ and $E(\lambda - 0) = s\text{- }\lim\limits_{\lambda' \uparrow \lambda} E(\lambda')$ do exist. 
\end{corollaryu}

\begin{proof}
 From \eqref{eq:XI.5.5}, we see that, if $\lambda_{n} \uparrow \lambda$, then 
\[
\lim\limits _{j, k \rightarrow \infty} \| E (\lambda_{j}, \lambda_{k} ] x \| ^{2} = 0,
\]
and the same is true for the case $\lambda_{n} \downarrow \lambda$. 
\end{proof}

\begin{proposition}\label{prop:XI.5.Proposition.2} Let $f(\lambda)$ be a complex-valued continuous function on $(-\infty, \infty)$, and let $x \in X$. Then we can define, for $-\infty < \alpha < \beta < \infty$, 
\[
\int_{\alpha} ^{\beta} f (\lambda) \dd E (\lambda) x
\]
as the \slim of the Riemann sums 
$\sum\limits_{j}f(\lambda_{j}^{\prime})E(\lambda_{j},\lambda_{j+1}]x,\text{ where }\alpha=\lambda_{1}<\lambda_{2}<\cdots<\lambda_{n}=\beta,x_{j}\in(\lambda_{j},\lambda_{j+1}],\text{ when the max }|\lambda_{j+1}-\lambda_{j}| \text{ tends to zero. }$ 
\end{proposition}

\begin{proof}
 $f(\lambda)$ is uniformly continuous on the compact interval $[\alpha, \beta]$. Let $|f(\lambda) - f(\lambda')| \leqq \varepsilon$ whenever $|\lambda - \lambda'| \leqq \delta$. We consider two partitions of $[\alpha, \beta]$: 
\begin{equation}\alpha = \lambda_{1} < \dots < \lambda_{n} = \beta , \max\limits _{j} | \lambda_{j+1} - \lambda_{j} | \leqq \delta , \alpha = \mu_{1} < \dots < \mu_{m} = \beta , \max\limits _{j} | \mu_{j+1} - \mu_{j} | \leqq \delta , \notag\end{equation}
and let 
\[
\alpha = \nu_{1} < \dots < \nu_{p} = \beta , p \leqq m + n,
\]
be the superposition of these two partitions. Then, if $\mu_k' \in (\mu_k, \mu_{k+1}]$, we have 
\[
\begin{aligned}
&\quad \sum\limits_{j} f (\lambda_{j} ^{\prime}) E (\lambda_{j}, \lambda_{j+1} ] x - \sum\limits_{k} f (\mu_{k} ^{\prime}) E (\mu_{k}, \mu_{k+1} ] x = \sum\limits_{s} \varepsilon_{s} E (\nu_{s}, \nu_{s+1} ] x, \text{ with } | \varepsilon_{s} | \leqq 2 \varepsilon ,
\end{aligned}
\]
and so the square of the norm of the left side is, as in \eqref{eq:XI.5.5}, 
\[
\leqq \varepsilon^{2} \left\| \sum\limits_{s} E (\nu_{s}, \nu_{s+1} ] x \right\| ^{2} = \varepsilon^{2} \left\| E (\alpha , \beta ] x \right\| ^{2} \leqq \varepsilon^{2} \left\| x \right\| ^{2}. \]
\end{proof}

\begin{corollaryu} We may define $\int_{\infty}^{\infty}f(\lambda)dE(\lambda)x$ as the $\slim\limits_{\alpha \downarrow -\infty ,\beta \uparrow \infty}\int_{\alpha}^{\beta}f(\lambda)dE(\lambda)x,$ 
when the right side limit exists. 
\end{corollaryu}

\begin{theorem}\label{thm:XI.5.Theorem.1} For a given $x \in X$, the following three conditions are mutually equivalent: 
\begin{align}
&\int_{- \infty} ^{\infty} f (\lambda) \dd E (\lambda) x \text{ exists }, \label{eq:XI.5.6} \\
&\int_{- \infty} ^{\infty} | f (\lambda) | ^{2} \dd \| E (\lambda) x \| ^{2} < \infty , \label{eq:XI.5.7} \\
F (y)& = \int_{- \infty} ^{\infty} f (\lambda) \dd (E (\lambda) y, x) \text{ defines a bounded linear functional. } \label{eq:XI.5.8}
\end{align}
\end{theorem}

\begin{proof}
 We shall prove the implications $\eqref{eq:XI.5.6} \to \eqref{eq:XI.5.8} \to \eqref{eq:XI.5.7} \to \eqref{eq:XI.5.6}$. 


\eqref{eq:XI.5.6} $\rightarrow$ \eqref{eq:XI.5.8}. The scalar product of $y$ with the approximate Riemann sum of $\int_{-\infty}^{\infty} f(\lambda) dE(\lambda)x$ is a bounded linear functional of $y$. Hence, by $(y, E(\lambda)x) = (E(\lambda)y, x)$ and the resonance theorem, we obtain \eqref{eq:XI.5.8}. 

\eqref{eq:XI.5.8} $\rightarrow$ \eqref{eq:XI.5.7}. We apply the operator $E(\alpha, \beta]$ to the approximate Riemann sum of $y = \int_{\alpha}^{\beta} f(\overline{\lambda}) dE(\lambda)x$. We then see, by \eqref{eq:XI.5.1}, that $y = E(\alpha, \beta]y$. Thus, again by \eqref{eq:XI.5.1}, 
\[
\begin{aligned}
&\overline {{F (y)}} = \int_{- \infty} ^{\infty} \overline {{f (\lambda)}} \dd (E (\lambda) x, y) = \lim\limits _ {\alpha^{\prime} \downarrow - \infty , \beta^{\prime} \uparrow \infty} \int_ {\alpha^{\prime}} ^ {\beta^{\prime}} \overline {{f (\lambda)}} \dd (E (\lambda) x, y)
\\
&= \lim\limits _ {\alpha^{\prime} \downarrow - \infty , \beta^{\prime} \uparrow \infty} \int_ {\alpha^{\prime}} ^ {\beta^{\prime}} \overline {{f (\lambda)}} \dd (E (\lambda) x, E (\alpha , \beta ] y) = \lim\limits _ {\alpha^{\prime} \downarrow - \infty , \beta^{\prime} \uparrow \infty} \int_ {\alpha^{\prime}} ^ {\beta^{\prime}} \overline {{f (\lambda)}} \dd (E (\alpha , \beta ] E (\lambda) x, y)
\\
&= \int_{\alpha} ^{\beta} \overline {{f (\lambda)}} \dd (E (\lambda) x, y) = \| y \| ^{2}.
\end{aligned}
\]

Hence $\| y\| ^2 \leqq \| F\| \cdot \| y\|$, i.e., $\| y\| \leqq \| F\|$. On the other hand, by approximating $y = \int_{\alpha}^{\beta}\overline{f(\lambda)} dE(\lambda)x$ by Riemann sums, we obtain, by \eqref{eq:XI.5.1}, 
\[
\left\| y \right\| ^{2} = \left\| \int_{\alpha} ^{\beta} \overline {{f (\lambda)}} \dd E (\lambda) x \right\| ^{2} = \int_{\alpha} ^{\beta} | f (\lambda) | ^{2} d \left\| E (\lambda) x \right\| ^{2},
\]
so that $\int_{\alpha}^{\beta}|f(\lambda)|^2 \dd\|E(\lambda)x\|^2 \leqq \|F\|^2$. Therefore, by letting $\alpha \downarrow -\infty$, $\beta \downarrow \infty$, we see that \eqref{eq:XI.5.7} is true. 

\eqref{eq:XI.5.7} $\rightarrow$ \eqref{eq:XI.5.6}. We have, for $\alpha' < \alpha < \beta < \beta'$, 
\[
\begin{aligned}
&\quad \left\| \int_ {\alpha^{\prime}} ^ {\beta^{\prime}} f (\lambda) \dd E (\lambda) x - \int_{\alpha} ^{\beta} f (\lambda) \dd E (\lambda) x \right\| ^{2} = \int_ {\alpha^{\prime}} ^{\alpha} | f (\lambda) | ^{2} \dd \| E (\lambda) x \| ^{2} + \int_{\beta} ^ {\beta^{\prime}} | f (\lambda) | ^{2} \dd \| E (\lambda) x \| ^{2}
\end{aligned}
\]
as above. Hence \eqref{eq:XI.5.7} implies \eqref{eq:XI.5.6}.
\end{proof}

\begin{theorem}\label{thm:XI.5.Theorem.2} Let $f(\lambda)$ be a real-valued continuous function. Then, a self-adjoint operator $H$ with $D(H) = D$ is defined through 
\begin{equation}\label{eq:XI.5.9}
\begin{gathered}
(H x, y) = \int_{- \infty} ^{\infty} f (\lambda) \dd (E (\lambda) x, y), \text{ where }
\\
x \in D = \left\{x; \int_{- \infty} ^{\infty} | f (\lambda) | ^{2} \dd \| E (\lambda) x \| ^{2} < \infty \right\} \text{ and any } y \in X,
\end{gathered}
\end{equation}
and we have $HE(\lambda) \supseteq E(\lambda)H$, that is, $HE(\lambda)$ is an extension of $E(\lambda)H$. 
\end{theorem}

\begin{proof}
 For any $y \in X$ and for any $\varepsilon > 0$, there exist $\alpha$ and $\beta$ with $-\infty < \alpha < \beta < \infty$ such that $\|y - E(\alpha, \beta] y\| < \varepsilon$. Moreover, we have 
\[
\int_{- \infty} ^{\infty} | f (\lambda) | ^{2} \dd \| E (\lambda) E (\alpha , \beta ] y \| ^{2} = \int_{\alpha} ^{\beta} | f (\lambda) | ^{2} \dd \| E (\lambda) y \| ^{2}. \]

Hence $E(\alpha, \beta] y \in D$ and so, by \eqref{eq:XI.5.2}, $D^a = X$. $H$ is symmetric by 
\[
f (\lambda) = \overline {{f (\lambda)}}, (E (\lambda) x, y) = \overline {{(E (\lambda) y , x)}}.
\]

If $y \in D(H^{*})$ and $H^{*}y = y^{*}$, then, by $E(\alpha, \beta]z \in D$ and \eqref{eq:XI.5.1}, 
\[
(z, E (\alpha , \beta ] y ^{*}) = (E (\alpha , \beta ] z, H ^{*} y) = (H E (\alpha , \beta ] z, y) = \int_{\alpha} ^{\beta} f (\lambda) \dd (E (\lambda) z, y).
\]

Thus, by the resonance theorem, 
\[
\lim\limits _{\alpha \downarrow - \infty , \beta \uparrow \infty} (z, E (\alpha , \beta ] y ^{*}) = \int_{- \infty} ^{\infty} f (\lambda) \dd (E (\lambda) z, y) = F (z)
\]
is a bounded linear functional. Hence, by the preceding theorem, 
\[
\int_{- \infty} ^{\infty} | f (\lambda) | ^{2} \dd \| E (\lambda) y \| ^{2} < \infty , \text{ that is, } y \in D.
\]

Therefore, $D = D(H) \supseteq D(H^*)$. Since $H$ is a symmetric operator, we have $H \subseteq H^*$ and so $H$ must be self-adjoint, i.e., $H = H^*$. 


Finally, let $x \in D(H)$. Then, by applying $E(\mu)$ to the approximate Rie-
mann sums of $Hx = \int_{-\infty}^{\infty} f(\lambda) dE(\lambda)x$, we obtain, by \eqref{eq:XI.5.1}, 
\[
\begin{aligned}
&E (\mu) H x  = \int_{- \infty} ^{\infty} f (\lambda) \dd (E (\mu) E (\lambda) x) = \int_{- \infty} ^{\infty} f (\lambda) \dd (E (\lambda) E (\mu) x) = H E (\mu) x.
\end{aligned}
\]
\end{proof}

\begin{corollary}\label{cor:XI.5.Corollary.1} In the particular case $f(\lambda) = \lambda$, we have 
\begin{equation}\label{eq:XI.5.10}
(H x, y) = \int_{- \infty} ^{\infty} \lambda \dd (E (\lambda) x, y), \quad \text{ for } \quad x \in D (H), y \in X.
\end{equation}
\end{corollary}

We shall write it symbolically 
\[
H = \int_{- \infty} ^{\infty} \lambda \dd E (\lambda),
\]
and call it the spectral resolution or the spectral representation of the self-adjoint operator $H$. 

\begin{corollary}\label{cor:XI.5.Corollary.2} We have, for $H = \int_{-\infty}^{\infty} f(\lambda) dE(\lambda)$ given by \eqref{eq:XI.5.9}, 
\begin{equation}\label{eq:XI.5.11}
\left\| H x \right\| ^{2} = \int_{- \infty} ^{\infty} | f (\lambda) | ^{2} d \left\| E (\lambda) x \right\| ^{2} \quad \text{ whenever } \quad x \in D (H).
\end{equation}
\end{corollary}

In particular, if $H$ is a bounded self-adjoint operator, then 
\begin{equation}\label{eq:XI.5.12}
(H ^{n} x, y) = \int_{- \infty} ^{\infty} f (\lambda) ^{n} \dd (E (\lambda) x, y) \quad \text{ for } \quad x, y \in X (n = 0, 1, 2, \dots).
\end{equation}

\begin{proof}
 Since $E(\lambda) Hx = HE(\lambda)x$ for $x \in D(H)$, we have, by \eqref{eq:XI.5.1}, 
\[ \begin{aligned} (H x, H x) & = \int f (\lambda) \dd (E (\lambda) x, H x) = \int f (\lambda) d (H E (\lambda) x, x) \\ & = \int f (\lambda) d _{\lambda} \left\{\int f (\mu) d _{\mu} (E (\mu) E (\lambda) x, x) \right\} \\ & = \int f (\lambda) d _{\lambda} \left\{\int_{- \infty} ^{\lambda} f (\mu) \dd (E (\mu) x, x) \right\} = \int f (\lambda) ^{2} \dd \| E (\lambda) x \| ^{2}. \end{aligned}
\]


The last part of the Corollary may be proved similarly.
\end{proof}

\begin{examplen} It is easy to see that the multiplication operator
\[
H x (t) = t x (t) \quad \text{ in } \quad L ^{2} (- \infty , \infty)
\]
admits the spectral resolution $H = \int_{-\infty}^{\infty}\lambda dE(\lambda)$, where 
\begin{equation}\label{eq:XI.5.13}
\begin{aligned} E (\lambda) x (t) & = x (t) \quad \text{ for } \quad t \leqq \lambda , \\ & = 0 \quad \text{ for } \quad t > \lambda . \end{aligned}
\end{equation}

For, 
\[
\begin{aligned}
\int_{- \infty} ^{\infty} \lambda^{2} \dd \| E (\lambda) x \| ^{2} &= \int_{- \infty} ^{\infty} \lambda^{2} d _{\lambda} \int_{- \infty} ^{\lambda} | x (t) | ^{2} \dd t = \int_{- \infty} ^{\infty} t ^{2} | x (t) | ^{2} \dd t = \| H x \| ^{2},
\\
\int_{- \infty} ^{\infty} \lambda \dd (E (\lambda) x, y) &= \int_{- \infty} ^{\infty} \lambda d _{\lambda} \int_{- \infty} ^{\lambda} x (t) \overline {{y (t)}} \dd t = \int_{- \infty} ^{\infty} t \cdot x (t) \overline {{y (t)}} \dd t = (H x, y).
\end{aligned}
\]


\end{examplen}
\yossection{6}{The Spectral Resolution of a Self-adjoint Operator}
\label{sec:XI.6}

\begin{theorem}\label{thm:XI.6.Theorem.1} $A$ self-adjoint operator $H$ in a Hilbert space $X$ admits a uniquely determined spectral resolution. 
\end{theorem}

\begin{proof}
 The Cayley transform $U = U_{H} = (H - iI)(H + iI)^{-1}$ of the self-adjoint operator $H$ is unitary (see Chapter~\ref{chap:VII}, \ref{sec:VII.4}). Let $U = \int_{0}^{2\pi} e^{i\theta} dF(\theta)$ be the spectral resolution of $U$. Then we have 
\[
F (2 \pi - 0) = \slim _{\theta \downarrow 0} F (2 \pi - \theta) = F (2 \pi) = I. \]

If otherwise, the projection $F(2\pi)-F(2\pi-0)$ would not be equal to the zero operator. Thus there exists an element $y \neq 0$ such that 
\[
(F (2 \pi) - F (2 \pi - 0)) y = y.
\]

Hence, by $F(\theta)F(\theta') = F(\min (\theta ,\theta'))$, 
\[
\begin{aligned}
&U y  = \int_{0} ^{2 \pi} e ^{i \theta} \dd (F (\theta) (F (2 \pi) - F (2 \pi - 0))) y = (F (2 \pi) - F (2 \pi - 0)) y = y.
\end{aligned}
\]

Thus $(y, z) = (Uy, Uz) = (y, Uz)$ and so $(y, z - Uz) = 0$ for every $z \in X$. $U$ being the Cayley transform of a self-adjoint operator $H$, we know (see Chapter~\ref{chap:VII}, \ref{sec:VII.4}) that the range $R(I - U)$ is dense in $X$. Hence we must have $y = 0$, which is a contradiction. 

Thus, if we set 
\[
\lambda = - \cot \theta , E (\lambda) = F (\theta),
\]
then $0 < \theta < 2\pi$ and $-\infty < \lambda < \infty$ are in a topological correspondence. Hence $E(\lambda)$ is a resolution of the identity with $F(\theta)$. We shall show that the self-adjoint operator 
\[
H ^{\prime} = \int_{- \infty} ^{\infty} \lambda \dd E (\lambda)
\]
is equal to $H$. Since $H = i(I + U)(I - U)^{-1}$, we have only to show that 
\[
(H ^{\prime} (y - U y), x) = (i (y + U y), x) \quad \text{ for all } \quad x, y \in X.
\]

But, since $D(H')^a = X$, we may restrict $x$ to be in the domain $D(H')$. Now, by $F(\theta) \cdot F(\theta') = F(\min(\theta, \theta'))$, 
\[
\begin{aligned}
&(y - U y, F (\theta) x)  = \int_{0} ^{2 \pi} (1 - e ^ {i \theta^{\prime}}) d _ {\theta^{\prime}} (F (\theta^{\prime}) y, F (\theta) x) = \int_{0} ^{2 \pi} (1 - e ^ {i \theta^{\prime}}) d _ {\theta^{\prime}} (F (\theta) F (\theta^{\prime}) y, x) \\
&= \int_{0} ^{\theta} (1 - e ^ {i \theta^{\prime}}) \dd (F (\theta^{\prime}) y, x).
\end{aligned}
\]

Hence 
\[
\begin{aligned}
&(y - U y, H ^{\prime} x)  = \int_{- \infty} ^{\infty} \lambda d (y - U y, E (\lambda) x) = \int_{0} ^{2 \pi} - \cot \theta d \left\{\int_{0} ^{\theta} (1 - e ^ {i \theta^{\prime}}) \dd (F (\theta^{\prime}) y, x) \right\} \\
&= \int_{0} ^{2 \pi} i (1 + e ^{i \theta}) \dd (F (\theta) y, x) = (i (y + U y), x).
\end{aligned}
\]

The uniqueness of the spectral representation. Suppose $H = \int_{-\infty}^{\infty} \lambda dE(\lambda)$ admits another spectral representation $H = \int_{-\infty}^{\infty} \lambda dE'(\lambda)$ such that 
$E^{\prime}(\lambda_{0})\neq E(\lambda_{0})$ for some $\lambda_0$. Then, by setting 
\[
\lambda = - \cot \theta , E ^{\prime} (\lambda) = F ^{\prime} (\theta),
\]
we have $F'(\theta_0) \neq F(\theta_0)$, where $\lambda_0 = -\cot \theta_0$. By a similar calculation to the above, we can prove that the Cayley transform of $\int_{-\infty}^{\infty} \lambda dE'(\lambda)$ is equal to $\int_{0}^{2\pi} e^{i\theta} dF'(\theta)$. Hence the unitary operator $U$ admits two different spectral representations $U = \int_{0}^{2\pi} e^{i\theta} dF(\theta)$ and $U = \int_{0}^{2\pi} e^{i\theta} dF'(\theta)$, contrary to what we have proved in Chapter~\ref{chap:XI}, \ref{sec:XI.4}. 
\end{proof}

We have thus proved (see Chapter~\ref{chap:VII}, \ref{sec:VII.3} and $4$) the fundamental result due to J. \textsc{von} \textsc{Neumann} \cite{neumann228}: 


\begin{theorem}\label{thm:XI.6.Theorem.2} $A$ symmetric operator $H$ has a closed symmetric extension $H^{**}$. A closed symmetric operator $H$ admits a uniquely determined spectral representation iff $H$ is self-adjoint. $H$ is self-adjoint iff its Cayley transform is unitary. 
\end{theorem}

\begin{remarku} It sometimes happens, in applications, that $H$ is not self-adjoint but $H^{*}$ is self-adjoint. In such a case, $H$ is said to be essentially self-adjoint. In this connection, see T. \textsc{Kato} \cite{kato140} concerning Schrödinger's operators in quantum mechanics. 
\end{remarku}

The spectral representation of the momentum operator $H_{1}$: 
\[
H _{1} x (t) = \frac {1}{i} \frac {d}{\dd t} x (t) \quad \text{ in } \quad L ^{2} (- \infty , \infty).
\]

The Fourier transform $U$ defined by 
\[
x (t) = U \cdot y (s) = \underset {n \rightarrow \infty} {\text{ l . i . m . }} (2 \pi) ^{-1/2} \int_{-n} ^{n} e ^{ist} y (s) \dd s
\]
is unitary and $U^{-1}x(t) = U^{*}x(t) = Ux(-t)$. Hence, denoting by $E(\lambda)$ the resolution of the identity given by \eqref{eq:XI.5.13} in Chapter~\ref{chap:XI}, \ref{sec:XI.5}, we obtain a resolution of the identity $\{E'(\lambda)\}, E'(\lambda) = UE(\lambda)U^{-1}$. If both $y(s)$ and $sy(s)$ belong to $L^2(-\infty, \infty) \cap L^1(-\infty, \infty)$, then 
\[
\begin{aligned}
\frac {1}{i} \frac {d}{\dd t} x (t) &= \frac {1}{i} \frac {d}{\dd t} \left((2 \pi) ^{-1/2} \int_{- \infty} ^{\infty} e ^{ist} y (s) \dd s\right)
\\
&= (2 \pi) ^{-1/2} \int_{- \infty} ^{\infty} e ^{ist} s y (s) \dd s = U (s y (s)) = U s U ^{-1} x (t),
\end{aligned}
\]
or, symbolically, 
\begin{equation}\label{eq:XI.6.1}
\frac {1}{i} \frac {d}{\dd t} = U s U ^{-1}.
\end{equation}

Hence, for the self-adjoint operator $H = s \cdot = \int_{-\infty}^{\infty} \lambda dE(\lambda)$, we have 
\begin{equation}\int U ^{-1} H _{1} U \cdot y (s) = s \cdot y (s) = H y (s) \text{ whenever } y (s), L ^{2} (- \infty , \infty) \cap L ^{1} (- \infty , \infty) \notag\end{equation}

For any $y(s) \in D(H) = D(s \cdot)$, set 
\[
y _{n} (s) = y (s) \text{ or } = 0 \text{ according as } | s | \leqq n \text{ or } | s | > n.
\]

Then surely $y_{n}(s)$, $s y_{n}(s)$ both $\in L^{2}(-\infty, \infty) \cap L^{1}(-\infty, \infty)$ and, moreover, \slim $y_{n} = y$, \slim $H y_{n} = H y$. Thus, since the self-adjoint operators $U^{-1}H_{1}U$ and $H$ are closed, we have, by $(U^{-1}H_{1}U)y_{n} = Hy_{n}$, 
\[
\left(U ^{-1} H _{1} U\right) y = H y \quad \text{ whenever } \quad y \in D (H),
\]
that is, $U^{-1}H_{1}U$ is a self-adjoint extension of the self-adjoint operator $H$. Hence, by taking the adjoint, we see that $H^{*}=H$ is also an extension of $(U^{-1}H_{1}U)^{*}=U^{-1}H_{1}U$; consequently $U^{-1}H_{1}U=H$ and so 
\[
H _{1} = U H U ^{-1} = \int_{- \infty} ^{\infty} \lambda d (U E (\lambda) U ^{-1}) = \int_{- \infty} ^{\infty} \lambda \dd E ^{\prime} (\lambda).
\]


\yossection{7}{Real Operators and Semi-bounded Operators. Friedrichs' Theorem}
\label{sec:XI.7}

Real operators and semi-bounded operators, defined below, have self-adjoint extensions. Thus we can apply von Neumann's theorem to these extensions to the effect that they admit spectral resolutions. 

\begin{definition}\label{def:XI.7.Definition.1} Let $X = L^2(S, \mathfrak{B}, m)$ and let $H$ be a symmetric operator defined in $X$ into $X$. $H$ is said to be a real operator, if \begin{enumerate}[i)]
\item $x(s) \in D(H)$ then $\overline{x(s)} \in D(H)$
\item $H$ maps real-valued functions into real-valued functions.
\end{enumerate}
\end{definition}

\begin{examplen} Let $f(s)$ be a real-valued continuous function in $(-\infty, \infty)$. Then, for $X = L^2(-\infty, \infty)$, the operator of multiplication by $f(s)$ is a real operator.

\end{examplen}
\begin{theorem}[J. von Neumann {\cite{neumann228}}]\label{thm:XI.7.Theorem.1} $A$ real operator $H$ admits a self-adjoint extension. 
\end{theorem}

\begin{proof}
 Let $U = U_H$ be the Cayley transform of $H$. Then $D(U) = \{(H + iI)x; x \in D(H)\}$ consists of the functions obtained by taking the complex-conjugate of the functions of $R(U) = \{(H - iI)x; x \in D(H)\}$. Thus, if we define an extension $U_1$ of $U$ through 
\begin{equation}U _{1} = U \text{ in } D (U), U _{1} \left(\sum\limits_{\alpha} c _{\alpha} \varphi_{\alpha}\right) = \sum\limits_{\alpha} c _{\alpha} \bar {\varphi} _{\alpha}, \text{ where } \{\varphi_{\alpha} \} \text{ is a complete } \notag\end{equation}
orthonormal system of the Hilbert space $D(U)^{\perp}$, 
then $U_{1}$ is a unitary extension of $U$. Therefore the self-adjoint extension $H_{1}$ of $H$ exists such that $U_{1} = U_{H_{1}}$ (see Chapter~\ref{chap:VII}, \ref{sec:VII.4}). 
\end{proof}

\begin{definition}\label{def:XI.7.Definition.2} $A$ symmetric operator $H$ is said to be upper semi-bounded (or lower semi-bounded) if there exists a real constant $\alpha$ such that 
\[
(H x, x) \leqq \alpha \| x \| ^{2} (\text{ or } (H x, x) \geqq \alpha \| x \| ^{2}) \text{ for all } x \in D (H)
\]
\end{definition}

If $(H x, x) \geqq 0$ for all $x \in D(H)$, then $H$ is said to be a positive operator. 

\begin{examplen} Let $q(s)$ be continuous and non-negative in $(-\infty, \infty)$. Consider the operator $H$ defined for $C^{2}$ functions $x(s)$ with compact support by
\[
(H x) (s) = - x ^{\prime \prime} (s) + q (s) x (s).
\]

Then $H$ is a positive operator in the Hilbert space $L^{2}(-\infty,\infty)$, as may be verified by partial integration. 

\end{examplen}
\begin{theorem}[K. \textsc{Friedrichs} {\cite{friedrichs71}}]\label{thm:XI.7.Theorem.2} $A$ semi-bounded operator $H$ admits a self-adjoint extension. 
\end{theorem}

Proof (due to H. \textsc{Freudenthal} \cite{freudenthal66}). $-H$ is lower semi-bounded if $H$ is upper semi-bounded. If $H$ is lower semi-bounded as above, then $H_{1} = H + (1 - \alpha)I$ satisfies the condition that $(H_{1}x, x) \geqq \|x\|^{2}$ for all $x \in D(H_{1})$. Therefore, since $\alpha I$ is self-adjoint, we may assume that the symmetric operator $H$ satisfies the condition 
\begin{equation}\label{eq:XI.7.1}
(H x, x) \geqq \| x \| ^{2} \text{ for all } x \in D (H).
\end{equation}

We introduce a new norm $\|x\|'$ and the associated new scalar product $(x, y)'$ in $D(H)$ through 
\begin{equation}\label{eq:XI.7.2}
\left\| x \right\| ^{\prime} = (H x, x), (x, y) ^{\prime} = (H x, y).
\end{equation}

Since $H$ is symmetric and satisfies \eqref{eq:XI.7.1}, it is easy to see that $D(H)$ becomes a pre-Hilbert space with respect to $\|x\|'$ and $(x, y)'$. We denote by $D(H)'$ the completion of this pre-Hilbert space. 

We shall show that $D(H)'$ is, as an abstract set without topology, a subset of the set $X$ which is the original Hilbert space. Proof: A Cauchy sequence $\{x_n\}'$ of the pre-Hilbert space $D(H)$ satisfies $\| x_n - x_m \|' \geqq \| x_n - x_m \|$ and $\lim\limits_{n,m \to \infty} \| x_n - x_m \|' = 0$; consequently $\{x_n\}$ is also a Cauchy sequence of the original Hilbert space $X$. If we can show, for a Cauchy sequence $\{y_n\}$ of $D(H)'$, that 
\begin{equation}\label{eq:XI.7.3}
\lim\limits _{n \rightarrow \infty} \| y _{n} \| ^{\prime} \neq 0 \text{ does not imply } \lim\limits _{n \rightarrow \infty} \| y _{n} \| = 0,
\end{equation}
then the correspondence 
\begin{equation}\label{eq:XI.7.4}
\left\{x _{n} \right\} ^{\prime} \rightarrow \left\{x _{n} \right\}
\end{equation}
is one-one from the Cauchy sequences of $D(H)$ into the Cauchy sequences of $X$. Two Cauchy sequences $\{x_n\}', \{z_n\}'$ of $D(H)$ (of $X$) being identified if $\lim\limits_{n \to \infty} \|x_n - z_n\|' = 0$ (if $\lim\limits_{n \to \infty} \|x_n - z_n\| = 0$). Since $X$ is complete, we may thus identify its Cauchy sequence $\{x_n\}$ with the element $x \in X$ such that $\lim\limits_{n\to \infty}\|x_n - x\| = 0$. Therefore, $D(H)'$ may, as an abstract set without topology, be identified with a subset of $X$ by the correspondence \eqref{eq:XI.7.4}. The proof of \eqref{eq:XI.7.3} is obtained, remembering the continuity of the scalar product in $D(H)'$ and in $X$, as follows. $\lim\limits_{n,m\to \infty}\|x_n - x_m\|' = 0$, $\lim\limits_{n\to \infty}\|x_n\|' = \alpha > 0$ and $\lim\limits_{n\to \infty}\|x_n\| = 0$ imply a contradiction 
\[
\alpha^{2} = \lim\limits _{n, m \rightarrow \infty} (x _{n}, x _{m}) ^{\prime} = \lim\limits _{n, m \rightarrow \infty} (H x _{n}, x _{m}) = \lim\limits _{n \rightarrow \infty} (H x _{n}, 0) = 0.
\]

We next set 
\begin{equation}\label{eq:XI.7.5}
\tilde {D} = D (H ^{*}) \cap D (H) ^{\prime}.
\end{equation}

Since $D(H) \subseteq D(H^{*})$, we must have $D(H) \subseteq \tilde{D} \subseteq D(H^{*})$. Hence we can define an extension $\tilde{H}$ of $H$ by restricting $H^{*}$ to the domain $\tilde{D} = D(\tilde{H})$. We have to show that $\tilde{H}$ is self-adjoint. 

We first show that $\tilde{H}$ is symmetric. Suppose $x, y \in \tilde{D}$; there exist two sequences $\{x_n\}'$, $\{y_n\}'$ of $D(H)$ such that $\|x - x_n\|' \to 0$, $\|y - y_n\|' \to 0$ as $n \to \infty$. Hence, by the continuity of the scalar product in $D(H)'$, we see that a finite $\lim\limits_{n,m \to \infty} (x_n, y_m)' = \lim\limits_{n,m \to \infty} (Hx_ny_m), = \lim\limits_{n,m \to \infty} (x_n, \tilde{H}y_m)$ exists. This limit is equal to 
\[
\lim\limits _{n \rightarrow \infty , m \rightarrow \infty} (H x _{n}, y _{m}) = \lim\limits _{n \rightarrow \infty} (H x _{n}, y) = \lim\limits _{n \rightarrow \infty} (x _{n}, \tilde {H} y) = (x, \tilde {H} y)
\]
and also to 
\[
\lim\limits _{m \rightarrow \infty} \lim\limits _{n \rightarrow \infty} (H x _{n}, y _{m}) = \lim\limits _{m \rightarrow \infty} (x, H y _{m}) = \lim\limits _{m \rightarrow \infty} (\tilde {H} x, y _{m}) = (\tilde {H} x, y).
\]

Hence $\tilde{H}$ is symmetric, that is, $\tilde{H} \subseteq (\tilde{H})^*$. 

Next let $x \in D(H)$, $y \in X$. Then, by 
\[
\left| (x, y) \right| \leqq \left\| x \| \right. \left\| y \right\| \leqq \left\| x \| ^{\prime} \cdot \right. \left\| y \right\|,
\]
we see that $f(x) = (x, y)$ is a bounded linear functional on the pre-Hilbert space $D(H)$. Hence $f(x)$ can, by continuity, be extended to a bounded linear functional on the Hilbert space $D(H)'$. Hence, by F. \textsc{Riesz}' representation theorem as applied to the Hilbert space $D(H)'$, there exists a uniquely determined $y' \in D(H)'$ such that 
\[
f (x) = (x, y) = \left(x, y ^{\prime}\right) ^{\prime} = \left(H x, y ^{\prime}\right) \quad \text{ for all } \quad x \in D (H).
\]

This proves that $y' \in D(H^*)$ and $H^*y' = y$. Hence $y' \in \tilde{D}$ and $\tilde{H}y' = y$. We have thus proved that $R(\tilde{H}) = X$, and so, by the Corollary of Theorem~\ref{thm:VII.3.Theorem.1} in Chapter~\ref{chap:VII}, \ref{sec:VII.3}, $\tilde{H}$ must be self-adjoint. 


\yossection[8. The Spectrum of a Self-adjoint Operator]{8}{The Spectrum of a Self-adjoint Operator. Rayleigh's Principle and the Krylov-Weinstein Theorem. The Multiplicity of the Spectrum}
\label{sec:XI.8}

\begin{theorem}\label{thm:XI.8.Theorem.1} Let $H = \int \lambda dE(\lambda)$ be a self-adjoint operator in a Hilbert space $X$. Let $\sigma(H), P_{\sigma}(H), C_{\sigma}(H)$ and $R_{\sigma}(H)$ be the spectrum, the point spectrum, the continuous spectrum and the residual spectrum of $H$, respectively. Then \begin{enumerate}[(i)]
\item $\sigma(H)$ is a set on the real line;
\item $\lambda_0 \in P_{\sigma}(H)$ is equivalent to the condition $E(\lambda_0) \neq E(\lambda_0 - 0)$ and the eigenspace of $H$ corresponding to the eigenvalue $\lambda_0$ is $R(E(\lambda_0) - E(\lambda_0 - 0))$;
\item $\lambda_0 \in C_{\sigma}(H)$ is equivalent to the condition $E(\lambda_0) = E(\lambda_0 - 0)$ in such a way that $E(\lambda_1) \neq E(\lambda_2)$ whenever $\lambda_1 < \lambda_0 < \lambda_2$;
\item $R_{\sigma}(H)$ is void.
\end{enumerate}
\end{theorem}

\begin{proof}
 We know already that, for a self-adjoint operator $H$, the resolvent set $\varrho(H)$ of $H$ comprises all the complex numbers $\lambda$ with $\operatorname{Im}(\lambda) \neq 0$ (see Chapter~\ref{chap:VIII}, \ref{sec:VIII.1}). Hence (i) is clear. We have $\lambda_0 I = \lambda_0 \int_{-\infty}^{\infty} dE(\lambda)$ by the definition of the resolution of the identity $\{E(\lambda)\}$ and so $(H - \lambda_0 I) = \int_{-\infty}^{\infty} (\lambda - \lambda_0) dE(\lambda)$. Hence, as in Corollary~\hyperref[cor:XI.5.Corollary.2]{$2$} of Theorem~\ref{thm:XI.5.Theorem.2} in Chapter~\ref{chap:XI}, \ref{sec:XI.5}, we obtain 
\begin{equation}\label{eq:XI.8.1}
\left\| (H - \lambda_{0} I) x \right\| ^{2} = \int_{- \infty} ^{\infty} (\lambda - \lambda_{0}) ^{2} d \left\| E (\lambda) x \right\| ^{2}, x \in D (H). \end{equation}

Thus, by $E(-\infty) = 0$ and the right continuity of $\|E(\lambda)x\|^2$ in $\lambda$, we see that $Hx = \lambda_0x$ iff 
\[
\begin{yoslines}
&E (\lambda) x = E (\lambda_{0} + 0) x = E (\lambda_{0}) x \quad \text{ for } \quad \lambda \geqq \lambda_{0} E (\lambda) x = E (\lambda_{0} - 0) x = 0 \quad \text{ for } \quad \lambda < \lambda_{0},
\end{yoslines}
\]
that is, $Hx = \lambda_0 x$ iff $(E(\lambda_0) - E(\lambda_0 - 0))x = x$. This proves (ii). Next we shall prove (iv). If $\lambda_0 \in R_\sigma(H)$, then, by (i), $\lambda_0$ is a real number. By the condition $R(H - \lambda_0 I)^a = D((H - \lambda_0 I)^{-1})^a \neq X$, we see that there exists a $y \neq 0$ which is orthogonal to $R(H - \lambda_0 I)$, i.e., $((H - \lambda_0 I)x, y) = 0$ for all $x \in D(H)$. Hence $(Hx, y) = (\lambda_0 x, y) = (x, \lambda_0 y)$ and so $y \in D(H^*)$, $H^*y = \lambda_0 y$. This proves that $Hy = \lambda_0 y$, i.e., $\lambda_0$ is an eigenvalue of $H$. Hence we have obtained a contradiction $\lambda_0 \in R_\sigma(H) \cap P_\sigma(H)$, and so $R_\sigma(H)$ must be a void set. 

Let $\lambda_0$ be a real number not belonging to $\sigma(H)$. Then the resolvent $(\lambda_0 I - H)^{-1}$ exists. Hence $H_{\lambda_0} = (H - \lambda_0 I)$ has a continuous inverse $(H - \lambda_0 I)^{-1}$. The last condition is, by (iv), equivalent to $\lambda_0 \in \varrho(H)$ and to the condition that there exists a positive number $\alpha$ such that 
\[
\left\| (H - \lambda_{0} I) x \right\| \geqq \alpha \left\| x \right\| \quad \text{ for all } \quad x \in D (H).
\]

This condition is by \eqref{eq:XI.8.1} equivalent to 
\begin{equation}\label{eq:XI.8.2}
\int_{- \infty} ^{\infty} (\lambda - \lambda_{0}) ^{2} \dd \| E (\lambda) x \| ^{2} \geqq \alpha^{2} \| x \| ^{2} \quad \text{ for all } \quad x \in D (H).
\end{equation}

Suppose that we have, for $\lambda_1 < \lambda_0 < \lambda_2$ with $\lambda_0 - \lambda_1 = \lambda_2 - \lambda_0 < \alpha$, $E(\lambda_1) = E(\lambda_2)$. Then, contrary to \eqref{eq:XI.8.2}, we obtain 
\[
\int_{- \infty} ^{\infty} (\lambda - \lambda_{0}) ^{2} \dd \| E (\lambda) x \| ^{2} < \alpha^{2} \int_{- \infty} ^{\infty} \dd \| E (\lambda) x \| ^{2} = \alpha^{2} \| x \| ^{2}.
\]

This proves (iii) by (i), (ii), (iv) and \eqref{eq:XI.8.2}.
\end{proof}

\begin{remarku} The Example in Chapter~\ref{chap:XI}, \ref{sec:XI.5} gives a self-adjoint operator $H$ for which all real numbers are in the continuous spectrum of $H$. 
\end{remarku}

\begin{theorem}\label{thm:XI.8.Theorem.2} Let $H$ be a bounded self-adjoint operator. Then 
\begin{equation}\label{eq:XI.8.3}
\sup\limits _{\lambda \in \sigma (H)} \lambda = \sup\limits _{\| x \| \leqq 1} (H x, x), \inf\limits _{\lambda \in \sigma (H)} \lambda = \inf\limits _{\| x \| \leqq 1} (H x, x).
\end{equation}
\end{theorem}

\begin{proof}
 Since $(Hx, x) = (x, Hx) = (\overline{Hx, x}) = \text{ real number }$, we can consider 
\[
\alpha_{1} = \inf\limits _{\| x \| \leqq 1} (H x, x) \quad \text{ and } \quad \alpha_{2} = \sup\limits _{\| x \| \leqq 1} (H x, x). \]

Let $\lambda_0 \in \sigma(H)$. Then, by Theorem~\ref{thm:XI.8.Theorem.1}, there exists, for every pair $(\lambda_1, \lambda_2)$ of real numbers with $\lambda_1 < \lambda_0 < \lambda_2$, a $y = y_{\lambda_1, \lambda_2} \neq 0$ such that $(E(\lambda_2) - E(\lambda_1))y = y$. We may assume that $\|y\| = 1$. Hence 
\[
\begin{aligned}
&(H y, y)  = \int \lambda \dd (E (\lambda) y, y) = \int \lambda \dd \| E (\lambda) y \| ^{2} = \int \lambda \dd \| E (\lambda) (E (\lambda_{2}) - E (\lambda_{1})) y \| ^{2} \\
&= \int_ {\lambda_{1}} ^ {\lambda_{2}} \lambda \dd \| (E (\lambda) - E (\lambda_{1})) y \| ^{2}.
\end{aligned}
\]

Thus, by letting $\lambda_1 \uparrow \lambda_0$, $\lambda_2 \downarrow \lambda_0$, we obtain $\lim (H y_{\lambda_1, \lambda_2}, y_{\lambda_1, \lambda_2}) = \lambda_0$. This proves that $\sup\limits_{\lambda \in g(H)} \lambda = \sup\limits_{\lambda_0} \leqq \alpha_2$. 

Let us assume that $\alpha_{2} \in \sigma(H)$. Then, by Theorem~\ref{thm:XI.8.Theorem.1}, there exists a pair $(\lambda_{1}, \lambda_{2})$ of real numbers such that $\lambda_{1} < \alpha_{2} < \lambda_{2}$ and $E(\lambda_{2}) = E(\lambda_{1})$. Hence $I = I - E(\lambda_{2}) + E(\lambda_{1})$, $(I - E(\lambda_{2}))E(\lambda_{1}) = E(\lambda_{1})(I - E(\lambda_{2})) = 0$ and so either $(I - E(\lambda_{2}))$ or $E(\lambda_{1})$ is not equal to the zero operator. If $(I - E(\lambda_{2})) \neq 0$, then there exists a $y$ with $\|y\| = 1$, $(I - E(\lambda_{2}))y = y$. In this case, we have 
\[
\begin{aligned}
&(H y, y)  = \int \lambda \dd \| E (\lambda) y \| ^{2} = \int \lambda \dd \| E (\lambda) (I - E (\lambda_{2})) y \| ^{2} = \int_ {\lambda_{1}} ^{\infty} \lambda \dd \| E (\lambda) y \| ^{2} \geqq \lambda_{2} > \alpha_{2};
\end{aligned}
\]
and in the case $E(\lambda_1) \neq 0$, we obtain 
\[
(H z, z) \leqq \lambda_{1} < \alpha_{2} \text{ for a } z \text{ with } \| z \| = 1, E (\lambda_{1}) z = z.
\]

Therefore the assumption $\alpha_{2} \bar{\in} \sigma(H)$ is absurd and so we have proved that $\sup\limits_{\lambda = \sup\limits_{\lambda \in \sigma(H)} H(x, x)} \inf\limits_{\|x\| \leqq 1} \inf\limits_{\lambda \in \sigma(H)} \lambda = \inf\limits_{\|x\| \leqq 1} (Hx, x)$. 
\end{proof}

\begin{theorem}[\textsc{Krylov}-\textsc{Weinstein}]\label{thm:XI.8.Theorem.3} Let $H$ be self-adjoint, and define, for any $x \in D(H)$ with $\| x \| = 1$, 
\begin{equation}\label{eq:XI.8.4}
\alpha_{x} = (H x, x), \beta_{x} = \| H x \|.
\end{equation}
\end{theorem}

Then, for any $\varepsilon > 0$, we can find a $\lambda_0 \in \sigma(H)$ satisfying the inequality 
\begin{equation}\label{eq:XI.8.5}
\alpha_{x} - (\beta_{x} ^{2} - \alpha_{x} ^{2}) ^{1/2} - \varepsilon \leqq \lambda_{0} \leqq \alpha_{x} + (\beta_{x} ^{2} - \alpha_{x} ^{2}) ^{1/2} + \varepsilon .
\end{equation}

\begin{proof}
 We have 
\[
\begin{aligned}
&\beta_{x} ^{2} = (H x, H x) = (H ^{2} x, x) = \int \lambda^{2} \dd \| E (\lambda) x \| ^{2}, \alpha_{x} = (H x, x) = \int \lambda \dd \| E (\lambda) x \| ^{2},
\\
&\| x \| ^{2} = \int \dd \| E (\lambda) x \| ^{2},
\end{aligned}
\]
and so 
\[ \begin{aligned} \beta_{x} ^{2} - \alpha_{x} ^{2} & = \int \lambda^{2} \dd \| E (\lambda) x \| ^{2} - 2 \alpha_{x} \int \lambda \dd \| E (\lambda) x \| ^{2} + \alpha_{x} ^{2} \int \dd \| E (\lambda) x \| ^{2} \\ & = \int (\lambda - \alpha_{x}) ^{2} \dd \| E (\lambda) x \| ^{2}. \end{aligned}
\]

Therefore, if $\left\|E(\lambda)x\right\|^{2}$ does not vary in the interval given by \eqref{eq:XI.8.5}, we would obtain a contradiction 
\[
\beta_{x} ^{2} - \alpha_{x} ^{2} \geqq ((\beta_{x} ^{2} - \alpha_{x} ^{2}) ^{1/2} + \varepsilon) ^{2} > \beta_{x} ^{2} - \alpha_{x} ^{2}.
\]
\end{proof}
\begin{remarku} The so-called Rayleigh principle consists in taking $\alpha_{x}$ as an approximation to the spectrum of the operator $H$. If we calculate $\beta_{x}$, then Theorem~\ref{thm:XI.8.Theorem.3} gives an upper bound of the error when we take $\alpha_{x}$ as an approximation of the spectrum of $H$. For a concrete application of such error estimate, we refer the reader to K. \textsc{Yosida} \cite{yosida344}. 
\end{remarku}

The Multiplicity of the Spectrum. We shall begin with the case of a self-adjoint matrix $H = \int \lambda dE(\lambda)$ in an $n$-dimensional Hilbert space $X_{n}$. Let $\lambda_1, \lambda_2, \ldots, \lambda_p (p \leqq n)$ be the eigenvalues of $H$ with the multiplicity $m_1, m_2, \ldots, m_p$, respectively $\left( \sum\limits_{j=1}^{p} m_j = n \right)$. Let $x_{j_1}, x_{j_2}, \ldots, x_{j_{m_j}}$ be orthonormal eigenvectors of $H$ belonging to the eigenvalue $\lambda_j (Hx_{j_{s}} = \lambda_j x_{j_{s}})$ so that $\{x_{j_{s}}; s = 1, 2, \ldots, m_j\}$ spans the eigenspace $E_{\lambda_j} = R(E(\lambda_j) - E(\lambda_j - 0))$ of $H$ belonging to the eigenvalue $\lambda_j$. Then the set $\{x_{j_{s}}; j = 1, 2, \ldots, p \text{ and } s = 1, 2, \ldots, m_j\}$ is a complete orthonormal system of vectors of the space $X_n$, and hence every vector $y$ of $X_n$ is represented uniquely as the linear combination of $x_{j_{s}}$'s: 
\begin{equation}\label{eq:XI.8.6}
y = \sum\limits_{j=1} ^{p} \sum\limits_{s=1} ^ {m _{j}} x _ {j _{s}} x _ {j _{s}}.
\end{equation}

Thus, denoting by $P_{\lambda_{j}}$ the projection $(E(\lambda_{j})-E(\lambda_{j}-0))$ upon the eigenspace $E_{\lambda_{j}}$, we have, for any $\alpha<\beta$, 
\begin{equation}\label{eq:XI.8.7}
(E (\beta) - E (\alpha)) y = \sum\limits_ {\alpha < \lambda_{j} \leqq \beta} \left(\sum\limits_{s=1} ^ {m _{j}} \alpha_ {j _{s}} x _ {j _{s}}\right) = \sum\limits_ {\alpha < \lambda_{j} \leqq \beta} P _ {\lambda_{j}} y,
\end{equation}
\begin{equation}\label{eq:XI.8.8}
P_{\lambda_j}(E(\beta) - E(\alpha))y = \sum\limits_{s=1}^{m_j}\alpha_{j_{s}}x_{j_{s}} \text{ or } = 0 \text{ according as } \alpha < \lambda_j\leqq \beta \text{ or not.}
\end{equation}

Hence, for fixed $\alpha < \beta$ and a fixed linear subspace $M$ of $X$, the set 
\[
\{(E (\beta) - E (\alpha)) y; y \in M \}
\]
does not contain $E_{\lambda_j}$ if the dimension of $M$, $\dim(M)$, is $< m_j$. Moreover, for a suitable $M$ with $\dim(M) = m_j$, the set $\{(E(\beta) - E(\alpha))y; y \in M\}$ with $\alpha < \lambda_j \leqq \beta$ contains $E_{\lambda_j}$. In fact, the statement is true for $M$ containing $x_{j_1}, x_{j_2}, \ldots, x_{j_{m_j}}$. In particular, $m_1 = m_2 = \cdots = m_p = 1$ with $p = n$ iff there exists a fixed vector $y \in X_n$ such that the set of vectors 
\[
\{(E (\beta) - E (\alpha)) y; \alpha < \beta \}
\]
spans the whole space $X_{n}$. 

These considerations lead to the following definitions: 

\begin{definition}\label{def:XI.8.Definition.1} The spectrum of a self-adjoint operator $H = \int \lambda dE(\lambda)$ in a Hilbert space $X$ is said to be simple, if there exists a fixed vector $y \in X$ such that the set of vectors $\{(E(\beta) - E(\alpha))y; \alpha < \beta\}$ spans a linear subspace strongly dense in $X$. 
\end{definition}

\begin{definition}\label{def:XI.8.Definition.2} Let $H = \int \lambda dE(\lambda)$ be a self-adjoint operator in a Hilbert space $X$. For fixed $\alpha < \beta$, consider linear subspaces $M$ of $(E(\beta) - E(\alpha)) \cdot X$ such that 
\begin{equation}\label{eq:XI.8.9}
(E (\beta) - E (\alpha)) \cdot M = (E (\beta) - E (\alpha)) \cdot X.
\end{equation}
\end{definition}

We may take, for example, $M = (E(\beta) - E(\alpha)) \cdot X$ to meet condition \eqref{eq:XI.8.9}. The minimum of the set of values $\dim(M)$, where $M$ satisfies condition \eqref{eq:XI.8.9}, will be called the total multiplicity of the spectrum of $H$ contained in the interval $(\alpha, \beta]$. 

\begin{definition}\label{def:XI.8.Definition.3} The multiplicity of the spectrum of a self-adjoint operator $H = \int \lambda dE(\lambda)$ at $\lambda = \lambda_0$ is defined as the limit, as $n \to \infty$, of the total multiplicity of the spectrum of $H$ contained in the interval $(\lambda_0 - n^{-1}, \lambda_0 + n^{-1}]$. 
\end{definition}

\begin{examplen} The coordinate operator $H$, i.e., the operator $H$ defined by $H \cdot x(t) = t \cdot x(t)$ in $L^2(-\infty, \infty)$, is of the simple spectrum.

\end{examplen}
\begin{proof}
 We know that the spectral resolution $H = \int \lambda dE(\lambda)$ is given by 
\[
E (\lambda) x (t) = x (t) \text{ or } = 0 \text{ according as } t \leqq \lambda \text{ or } t > \lambda . \]

Let $y(t)$ be defined by 
\[
y (t) = c _{k} > 0 \text{ for } k - 1 < t \leqq k \quad (k = 0, \pm 1, \pm 2, \dots)
\]
with $\sum\limits_{k} c_k^2 < \infty$ so that $y(t) \in L^2(-\infty, \infty)$. 

Then it is easy to see that the linear combinations of vectors of the form $(E(\beta)-E(\alpha))y,\alpha<\beta$, are strongly dense in the totality of step functions with compact support and consequently strongly dense in $L^{2}(-\infty,\infty)$. 
\end{proof}

The Problem of the Unitary Equivalence. Two self-adjoint operators $H_{1}$ and $H_{2}$ in an $n$-dimensional Hilbert space $X_{n}$ are said to be unitarily equivalent to each other if there exists a unitary matrix $U$ in $X_{n}$ such that $H_{1}=UH_{2}U^{-1}$. It is well-known that $H_{1}$ and $H_{2}$ are unitarily equivalent to each other iff they have the same system of eigenvalues with respectively the same multiplicities. Thus the eigenvalues together with the respective multiplicities are the unitary invariants of a self-adjoint matrix. 

The investigation of the unitary invariants for a self-adjoint operator in an infinite dimensional Hilbert space goes back to a paper by E. \textsc{Hellinger} \cite{hausdorff85} published in $1909$. See, e.g., M. H. \textsc{Stone} \cite{stone311}. There the Hilbert space is assumed to be separable. For the non-separable Hilbert space case, see F. \textsc{Wecken} \cite{wecken333} and H. \textsc{Nakano} \cite{nakano224} and also P. R. \textsc{Halmos} \cite{halmos93}. K. \textsc{Yosida} \cite{yosida356} proved the following theorem: 

Let $H$ be a self-adjoint operator in a Hilbert space $X$, and let us denote by $(H)'$ the totality of bounded linear operators $\in L(X, X)$ which are commutative with $H$. Then two self-adjoint operators $H_1$ and $H_2$ in $X$ are unitarily equivalent to each other iff there exists a one-to-one mapping $T$ of $(H_1)'$ onto $(H_2)'$ such that $T$ defines a ring-isomorphism of the ring $(H_1)'$ with the ring $(H_2)'$ in such a way that $(T \cdot B)^* = T \cdot B^*$ for every $B \in (H_1)'$. 

Thus the algebraic structure of the ring $(H_{\mathbf{1}})'$ is the unitary invariant of $H_{\mathbf{1}}$. 


\yossection{9}{The General Expansion Theorem. A Condition for the Absence of the Continuous Spectrum}
\label{sec:XI.9}

Let $H = \int \lambda dE(\lambda)$ be a self-adjoint operator in a Hilbert space $X$. Then, by $E(+\infty) = I$ and $E(-\infty) = 0$, we have the representation 
\begin{equation}\label{eq:XI.9.1}
x = \underset {\lambda_{1} \downarrow - \infty , \lambda_{2} \uparrow \infty} {\slim} \int_ {\lambda_{1}} ^ {\lambda_{2}} \dd E (\lambda) x = \underset {\lambda_{1} \downarrow - \infty , \lambda_{2} \uparrow \infty} {\slim} (E (\lambda_{2}) - E (\lambda_{1})) x \text{ for every } x \in X.
\end{equation}

We shall call \eqref{eq:XI.9.1} the general expansion theorem associated with the self-adjoint operator $H$. In concrete cases, it sometimes happens that the resolvent $(\lambda I - H)^{-1}$ is obtainable more easily than the spectral resolution $H = \int \lambda dE(\lambda)$. In such a case, the general expansion theorem \eqref{eq:XI.9.1}  
may be replaced by 
\begin{equation}\tag{1'}\label{eq:XI.9.1'}
\begin{aligned}
&x = \underset {\alpha \downarrow - \infty , \beta \uparrow \infty} {\slim} s - \underset {v \downarrow 0} {\lim} \frac {1}{2 \pi i} \left[ \int_{\alpha} ^{\beta} ((u - i v) I - H) ^{-1} x \dd u \right. \left. + \int_{\beta} ^{\alpha} ((u + i v) I - H) ^{-1} x \dd u \right], \quad x \in X.
\end{aligned}
\end{equation}

\begin{proof}
 If $v \neq 0$, then 
\[
((u + i v) I - H) ^{-1} x = \int_{- \infty} ^{\infty} \frac {1}{u + i v - \lambda} \dd E (\lambda) x, \text{ for every } x \in X. \]

For, by approximating the integral by Riemann sums and remembering the relation $E(\lambda)E(\lambda') = E(\min (\lambda ,\lambda'))$, we obtain, for $\operatorname{Im}(\mu)\neq 0$, 
\[
\begin{aligned}
&\int_{- \infty} ^{\infty} (\lambda - \mu) \dd E (\lambda) \left\{\int_{- \infty} ^{\infty} \frac {1}{\lambda^{\prime} - \mu} \dd E (\lambda^{\prime}) \right\} = \int_{- \infty} ^{\infty} (\lambda - \mu) d _{\lambda} \left\{\int_{- \infty} ^{\infty} \frac {1}{(\lambda^{\prime} - \mu)} d _ {\lambda^{\prime}} (E (\lambda) E (\lambda^{\prime})) \right\}
\\
&= \int_{- \infty} ^{\infty} (\lambda - \mu) d _{\lambda} \left\{\int_{- \infty} ^{\lambda} \frac {1}{\lambda^{\prime} - \mu} \dd E (\lambda^{\prime}) \right\} = \int_{- \infty} ^{\infty} \dd E (\lambda) = I,
\\
&\int_{- \infty} ^{\infty} \frac {1}{\lambda - \mu} d _{\lambda} \left\{E (\lambda) \int_{- \infty} ^{\infty} (\lambda^{\prime} - \mu) \dd E (\lambda^{\prime}) \right\} = \int_{- \infty} ^{\infty} \frac {1}{\lambda - \mu} d _{\lambda} \left\{\int_{- \infty} ^{\infty} (\lambda^{\prime} - \mu) d _ {\lambda^{\prime}} (E (\lambda) E (\lambda^{\prime})) \right\}
\\
&= \int_{- \infty} ^{\infty} \frac {1}{\lambda - \mu} d _{\lambda} \left\{\int_{- \infty} ^{\lambda} (\lambda^{\prime} - \mu) \dd E (\lambda^{\prime}) \right\} = \int_{- \infty} ^{\infty} \dd E (\lambda) = I.
\end{aligned}
\]

Therefore, we have 
\[
\begin{aligned}
& \int_{\alpha} ^{\beta} ((u - i v) I - H) ^{-1} x \dd u + \int_{\beta} ^{\alpha} ((u + i v) I - H) ^{-1} x \dd u
\\
&= \int_{- \infty} ^{\infty} \dd E (\lambda) x \left\{\int_{\alpha} ^{\beta} \frac {\dd u}{u - i v - \lambda} + \int_{\beta} ^{\alpha} \frac {\dd u}{u + i v - \lambda} \right\}
\\
&= \int_{- \infty} ^{\infty} \dd E (\lambda) x \left\{\int_{\alpha} ^{\beta} d _{u} \log (u - i v - \lambda) + \int_{\beta} ^{\alpha} d _{u} \log (u + i v - \lambda) \right\}
\end{aligned}
\]
which tends, when $v \downarrow 0$, strongly to 
\[
\begin{aligned}
&\int_{\alpha +0} ^{\beta -0} 2 \pi i \dd E (\lambda) x + \pi i (E (\beta) - E (\beta - 0)) x + \pi i (E (\alpha) - E (\alpha - 0)) x
\\
&= \pi i (E (\beta) + E (\beta - 0)) x - \pi i (E (\alpha) + E (\alpha - 0)) x.
\end{aligned}
\]

This proves formula \eqref{eq:XI.9.1}. 
\end{proof}

\begin{remarku} The eigenfunction expansion associated with the second order differential operator 
\[
- \frac {d ^{2}}{\dd x ^{2}} + q (x)
\]
\end{remarku}

with real continuous $q(x)$ in an open interval $(a, b)$ was inaugurated by H. \textsc{Weyl} \cite{weyl337}, further developed by M. H. \textsc{Stone} \cite{stone311}, and completed by E. C. \textsc{Titchmarsh} \cite{titchmarsh324} and K. \textsc{Kodaira} \cite{kodaira151} who gave a formula which determines the expansion explicitly. The expansion is exactly a concrete application of \eqref{eq:XI.9.1}. The crucial point in their theory is to give the possible boundary conditions at $x = a$ and $x = b$ so that the operator 
\[
- \frac {d ^{2}}{\dd x ^{2}} + q (x)
\]
becomes a self-adjoint operator $H$ in the Hilbert space $L^{2}(a,b)$. Their theory is very important in that it gives a unified treatment of the classical expansions in terms of special functions, such as the Fourier series expansion, the Fourier integral representation, the Hermite polynomials expansion, the Laguerre polynomials expansion and the Bessel functions expansion. We do not go into details, and refer the reader to the above cited book by \textsc{Titchmarsh} and the paper by \textsc{Kodaira}. Cf. also N.A.\textsc{Naimark} \cite{naimark222}, N.\textsc{Dunford}-J.\textsc{Schwartz} \cite{dunford48} and K.\textsc{Yosida} \cite{yosida344}. The last cited book gives an elementary treatment of the theory. 

If the continuous spectrum $C_{\sigma}(H)$ is absent, then the expansion \eqref{eq:XI.9.1} will be replaced by a series rather than the integral. We have, for instance, the following 

\begin{theorem}\label{thm:XI.9.Theorem.1} Let $H = \int \lambda dE(\lambda)$ be a self-adjoint compact operator in a Hilbert space $X$. Then \begin{enumerate}[(i)]
\item $C_{\sigma}(H)$ contains no real number except possibly $0$;
\item the eigenvalues of $H$ constitute at most a countable system of real numbers accumulating only at $0$;
\item for any eigenvalue $\lambda_{0} \neq 0$ of $H$, the corresponding eigenspace $E_{\lambda_{0}}$ is of finite dimension.
\end{enumerate}
\end{theorem}

\begin{proof}
 Suppose a closed interval $[\lambda', \lambda'']$ on the real line does not contain the number $0$. Then the range $R(E(\lambda'') - E(\lambda'))$ is of finite dimension. If otherwise, there exists, by E. \textsc{Schmidt}'s orthogonalization in Chapter~\ref{chap:III}, \ref{sec:III.5}, a countable orthonormal system $\{x_j\}$ contained in $R(E(\lambda'') - E(\lambda'))$. We have \wlim $x_j = 0$, since, by Bessel's inequality, 
\[
\sum\limits_{j=1} ^{\infty} | (f, x _{j}) | ^{2} \leqq \| f \| ^{2} \quad \text{ for any } \quad f \in X. \]

Hence, by the compactness of the operator $H$, there exists a subsequence $\{x_{j'}\}$ such that $\slim\limits_{j \to \infty} H x_{j'} = w - \lim\limits_{j \to \infty} H x_{j'} = 0$. On the other hand, we have 
\[
\begin{aligned}
&\| H x _{i} \| ^{2}  = \int \lambda^{2} \dd \| E (\lambda) x _{j} \| ^{2} = \int \lambda^{2} \dd \| E (\lambda) (E (\lambda^{\prime \prime}) - E (\lambda^{\prime})) x _{j} \| ^{2} \\
&= \int_ {\lambda^{\prime}} ^ {\lambda^{\prime \prime}} \lambda^{2} \dd \| (E (\lambda) - E (\lambda^{\prime})) x _{j} \| ^{2} \geqq \| x _{j} \| ^{2} \cdot \min (| \lambda^{\prime} | ^{2}, | \lambda^{\prime \prime} | ^{2}) = \min (| \lambda^{\prime} | ^{2}, | \lambda^{\prime \prime} | ^{2}),
\end{aligned}
\]
which is a contradiction. 

If $C_{\sigma}(H)$ contains a number $\lambda_0 \neq 0$, then, by Theorem~\ref{thm:XI.9.Theorem.1} in the preceding section, $\slim\limits_{\varepsilon \downarrow 0} (E(\lambda_0 - \varepsilon) - E(\lambda_0 - \varepsilon)) x = 0$ for any $x \in X$. As proved above, the range $R(E(\lambda_0 + \varepsilon) - E(\lambda_0 - \varepsilon))$ is of finite dimension and this dimension number is monotone decreasing as $\varepsilon \downarrow 0$. Hence we see that $(E(\lambda_0 + \varepsilon) - E(\lambda_0 - \varepsilon)) = 0$ for sufficiently small $\varepsilon > 0$, and so, by Theorem~\ref{thm:XI.9.Theorem.1} in the preceding section, $\lambda_0$ cannot be contained in $C_{\sigma}(H)$. This proves our Theorem. 
\end{proof}

\begin{corollary}\label{cor:XI.9.Corollary.1} Let $\{\lambda_j\}$ be the system of all eigenvalues of $H$ different from $0$. Then, for any $x \in X$, we have 
\begin{equation}\label{eq:XI.9.2}
x = (E (0) - E (0 - 0)) x + \slim _{n \rightarrow \infty} \sum\limits_{j=1} ^{n} (E (\lambda_{j}) - E (\lambda_{j} - 0)) x.
\end{equation}
\end{corollary}

\begin{proof}
 Clear from \eqref{eq:XI.9.1}. 
\end{proof}

\begin{corollary}[The Hilbert-Schmidt Expansion Theorem]\label{cor:XI.9.Corollary.2} For any $x \in X$, we have 
\begin{equation}\label{eq:XI.9.3}
H x = \underset {n \rightarrow \infty} {\text{ \slim }} \sum\limits_{j=1} ^{n} \lambda_{j} (E (\lambda_{j}) - E (\lambda_{j} - 0)) x.
\end{equation}
\end{corollary}

\begin{proof}
 Clear from the continuity of the operator $H$ and the fact that $H(E(0)-E(0-0))=0$ and $H(E(\lambda_{j})-E(\lambda_{j}-0))x=\lambda_{j}(E(\lambda_{j})-E(\lambda_{j}-0))x$, the latter being implied by $R(E(\lambda_{j})-E(\lambda_{j}-0))=E_{\lambda_{j}}$, the eigenspace of $H$ belonging to the eigenvalue $\lambda_{j}$. 
\end{proof}

\begin{remarku} The strong convergence in \eqref{eq:XI.9.3} is uniform in the unit circle $\{x; \| x \| \leqq 1\}$. For, we have 
\[
\begin{aligned}
\left\| H x - \int_{| \lambda | > \varepsilon} \lambda \dd E (\lambda) x \right\| ^{2} &= \left\| \int_{| \lambda | \leqq \varepsilon} \lambda \dd E (\lambda) x \right\| ^{2} = \int_{- \varepsilon} ^{\varepsilon} \lambda^{2} \dd \| E (\lambda) x \| ^{2}
\\
\quad \leqq \varepsilon^{2} \int_{- \varepsilon} ^{\varepsilon} \dd \| E (\lambda) x \| ^{2} \leqq \varepsilon^{2} \int \dd \| E (\lambda) x \| ^{2} &= \varepsilon^{2} \| x \| ^{2} \leqq \varepsilon^{2}.
\end{aligned}
\]
\end{remarku}


\yossection{10}{The Peter-Weyl-Neumann Theory}
\label{sec:XI.10}

Let $G$ be a totally bounded topological group, metrized by a distance satisfying the condition (see Chapter~\ref{chap:VIII}, \ref{sec:VIII.5}) 
\begin{equation}\label{eq:XI.10.1}
\dist (x, y) = \dist (a x b, a y b) \text{ for every } x, y, a \text{ and } b \in G.
\end{equation}

Let $f(g)$ be a complex-valued bounded uniformly continuous function defined on $G$. For any $\varepsilon > 0$, set 
\begin{equation}\label{eq:XI.10.2}
V = \left\{y \in G; \sup\limits _{x \in G} | f (x) - f (y ^{-1} x) | < \varepsilon \right\}.
\end{equation}

Then, by the continuity of $f$, the set $V$ is an open set containing the unit $e$ of the group $G$. Hence the set $U = V \cap V^{-1}$, where $V^{-1} = \{y^{-1}; y \in V\}$, is also open and $\ni e$. If we put 
\begin{align}
k (x)& = 2 ^{-1} \left(k _{1} (x) + k _{1} \left(x ^{-1}\right)\right), \text{ where } \notag \\
k _{1} (x)& = \frac {\dist \left(x , U ^{c}\right)}{\dist (x , e) + \dist \left(x , U ^{c}\right)} \quad \left(\dist \left(x, U ^{c}\right) = \inf\limits _ {y \in U ^{c}} \dist (x, y)\right), \label{eq:XI.10.3}
\end{align}
then we have the results: 

$k(x)$ is bounded uniformly continuous on $G$, and 
\begin{equation}k (x) = k \left(x ^{-1}\right), \quad 0 \leqq k (x) \leqq 1 \text{ on } G, \quad k (e) = 1 \text{ and } k (x) = 0 \text{ whenever } x \in U ^{C}. \label{eq:XI.10.4}\end{equation}

Hence, for all $x, y \in G$, we obtain 
\[
\left| k (y) \left(f (x) - f \left(y ^{-1} x\right)\right) \right| \leqq \varepsilon k (y).
\]

By taking the mean value (see Chapter~\ref{chap:VIII}, \ref{sec:VIII.5}) of both sides, we obtain 
\[
\left| M _{y} (k (y)) f (x) - M _{y} (k (y) f \left(y ^{-1} x\right)) \right| \leqq \varepsilon M _{y} (k (y)).
\]

We have $M_y(k(y)) > 0$ by $k(y) \geqq 0$ and $k(y) \not\equiv 0$. Hence 
\begin{equation}\label{eq:XI.10.5}
\left| f (x) - M _{y} \left(k _{0} (y) f \left(y ^{-1} x\right)\right) \right| \leqq \varepsilon , \quad \text{ where } \quad k _{0} (x) = k (x) / M _{x} (k (x)).
\end{equation}

Thus, by virtue of the invariance $M_y(g(y^{-1})) = M_y(g(y)) = M_y(g(ay)) = M_y(g(ya))$ of the mean value, we obtain, from \eqref{eq:XI.10.5}, 
\begin{equation}\label{eq:XI.10.6}
\left| f (x) - M _{y} \left(k _{0} \left(x y ^{-1}\right) f (y)\right) \right| \leqq \varepsilon \quad \text{ for all } \quad x \in G.
\end{equation}

\begin{proposition}\label{prop:XI.10.Proposition.1} We shall denote by $C(G)$ the set of all complex-valued bounded uniformly continuous functions $h(g)$ defined on $G$. Then $C(G)$ is a $B$-space by the norm $\|h\|_0 = \sup |h(g)|$. Then, for any $b$ and $h \in C(G)$, 
\begin{equation}\label{eq:XI.10.7}
(b \times h) (x) = M _{y} (b (x y ^{-1}) h (y))
\end{equation}
also belongs to $C(G)$. 
\end{proposition}

\begin{proof}
 By $\dist(x,z)=\dist(axc,azc)$ and the uniform continuity of the function $b(g)$, there exists, for any $\delta>0$, an $\eta=\eta(\delta)>0$ such that 
\[
\sup\limits _{y} \left| b (x y ^{-1}) - b \left(x ^{\prime} y ^{-1}\right) \right| \leqq \delta \quad \text{ whenever } \quad \dist (x, x ^{\prime}) \leqq \eta . \]

Hence, as in the case of Schwarz' inequality, we obtain 
\begin{equation}\label{eq:XI.10.8}
\begin{aligned}
& \left| M _{y} (b (x y ^{-1}) h (y)) - M _{y} (b (x ^{\prime} y ^{-1}) h (y)) \right| ^{2}
\\
&\leqq M _{y} ((b (x y ^{-1}) - b (x ^{\prime} y ^{-1})) ^{2}) \cdot M _{y} (| h (y) | ^{2}) \leqq \delta^{2} M _{y} (| h (y) | ^{2})
\end{aligned}
\end{equation}
whenever $\dist(x, x') \leqq \eta$. This proves our Proposition~\ref{prop:XI.10.Proposition.1}. 
\end{proof}

\begin{proposition}\label{prop:XI.10.Proposition.2} $C(G)$ is a pre-Hilbert space by the operation of the function sum and the scalar product 
\begin{equation}\label{eq:XI.10.9}
(b, h) = (b \times h ^{*}) (e) = M _{y} (b (y ^{-1}) \overline {{h (y ^{-1})}}) = M _{y} (b (y) \overline {{h (y)}}),
\end{equation}
$h^{*}(y) = \overline{h(y^{-1})}$. We shall denote this pre-Hilbert space by $\hat{C}(G)$. 
\end{proposition}

\begin{proof}
 Easy. 
\end{proof}

\begin{proposition}\label{prop:XI.10.Proposition.3} We shall denote the completion of the pre-Hilbert space $\hat{C}(G)$ by $\tilde{C}(G)$, and the norm in the Hilbert space $\tilde{C}(G)$ by $\|h\| = (h, h)^{1/2}$. Then the linear mapping $T$ on $\hat{C}(G)$ into $\hat{C}(G)$ defined by 
\begin{equation}\label{eq:XI.10.10}
(T h) (x) = \left(k _{0} \times h\right) (x), x \in G,
\end{equation}
can, by the continuity in $\hat{C}(G)$, be extended to a compact linear operator $\tilde{T}$ on $\tilde{C}(G)$ into $\tilde{C}(G)$. 
\end{proposition}

\begin{proof}
 By virtue of $M_y(1) = 1$, we obtain 
\begin{equation}\label{eq:XI.10.11}
\left\| h \right\| = (h, h) ^{1/2} = M _{y} (h (y) \overline {{{{h (y)}}}}) ^{1/2} \leqq \sup\limits _{y} | h (y) | = \left\| h \right\| _{0}. \end{equation}

The continuity of the operator $T$ in $\hat{C}(G)$ is clear from the Schwarz inequality for \eqref{eq:XI.10.7}, and so, by the denseness of $\hat{C}(G)$ in $\tilde{C}(G)$, we can extend $T$ to a bounded linear operator $\tilde{T}$ in $\tilde{C}(G)$. On the other hand, we see, by \eqref{eq:XI.10.8}, that $T$ is a compact operator on $\hat{C}(G)$ into $C(G)$. We prove this by the Ascoli-Arzelà theorem. Thus, by \eqref{eq:XI.10.11}, we easily see that $\tilde{T}$ is a compact operator on $\hat{C}(G)$ into $\hat{C}(G)$. 

Therefore, by the denseness of $\hat{C}(G)$ in $\tilde{C}(G)$, we see that the extended operator $\tilde{T}$ is also compact as an operator on $\tilde{C}(G)$ into $\tilde{C}(G)$. 
\end{proof}

We are now ready to prove the Peter-Weyl-Neumann theory on the representation of almost periodic functions. 

It is easy to see that, by $k_{0}(xy^{-1}) = k_{0}(yx^{-1})$, the compact operator $\tilde{T}$ is self-adjoint in the Hilbert space $\tilde{C}(G)$. Hence, by the Hilbert-Schmidt expansion theorem in the preceding section, we obtain 
\begin{equation}\label{eq:XI.10.12}
\tilde {T} h = \slim _{n \to \infty} \sum\limits_{m=1} ^{n} \lambda_{m} P _ {\lambda_{m}} h \text{ uniformly in } h \text{ satisfying } \| h \| \leqq 1,
\end{equation}
if we denote by $\{\lambda_{m}\}$ the system of all eigenvalues of $\tilde{T}$ different from $0$, and by $P_{\lambda_{m}}$ the projection upon the eigenspace of $\tilde{T}$ belonging to the eigenvalue $\lambda_{m}$. 

Since $f(g)$, introduced at the beginning of this section, belongs to $C(G)$, we have $\tilde{T}f = Tf \in C(G)$. Since the eigenspace $R(P_{\lambda_m}) = P_{\lambda_m} \cdot \tilde{C}(G)$ is of finite dimension, there exists, for every eigenvalue $\lambda_m$, a finite system $\{h_{mj}\}_{j=1,\ldots,n_m}$ of elements $\in \tilde{C}(G)$ such that each $h \in R(P_{\lambda_m}) =$ 
$P_{\lambda_m} \cdot \tilde{C}(G)$ can be represented as a uniquely determined linear combination of $h_{mj}$'s ( $j = 1, 2, \ldots, n_m$). Let 
\begin{equation}\label{eq:XI.10.13}
P _ {\lambda_{m}} h = \sum\limits_{j=1} ^ {n _{m}} c _{j} h _{mj}, \text{ where } c ^{\prime} \text{ s are complex numbers. }
\end{equation}

Then, since $h_{mj} \in R(P_{\lambda_m})$, we have $\tilde{T}h_{mj} = \lambda_m h_{mj}$, and so, by applying \eqref{eq:XI.10.8} to the operator $T$ given by \eqref{eq:XI.10.10}, we see that $h_{mj} = \lambda_m^{-1}(\tilde{T}h_{mj})$ must belong to $C(G)$. Hence, by \eqref{eq:XI.10.13}, we see that, for each eigenvalue $\lambda_m$ of $\tilde{T}$, the eigenspace $R(P_{\lambda_m}) = P_{\lambda_m} \cdot \tilde{C}(G)$ is spanned by the functions $h_{mj} \in C(G)$. 

Therefore, we have, by \eqref{eq:XI.10.12}, 
\[
(Tf)(x)=\slim_{n\to\infty}\sum_{m=1}^n\lambda_m f_m(x)
\]
in the strong topology of $\tilde C(G)$, where $f_m=P_{\lambda_m}\cdot f\in C(G)$ for each $m$.

By applying \eqref{eq:XI.10.8} to the operator $T$ given by \eqref{eq:XI.10.10}, we see that 
\begin{equation}\label{eq:XI.10.14}
(T ^{2} f) (x) = \lim\limits _{n \rightarrow \infty} M _{y} \left(k _{0} (x y ^{-1}) \cdot \sum\limits_{m=1} ^{n} \lambda_{m} f _{m} (y)\right) \text{ uniformly in } x.
\end{equation}

On the other hand, by \eqref{eq:XI.10.6} and $M_y(k_0(y)) = 1$, we obtain 
\[
\left| M _{z} \left(k _{0} (x z ^{-1}) f (z)\right) - M _{z} \left(k _{0} (x z ^{-1}) M _{y} \left(k _{0} (z y ^{-1}) f (y)\right) \right. \right| \leqq \varepsilon ,
\]
and so, combined with \eqref{eq:XI.10.6}, we have 
\begin{equation}\label{eq:XI.10.15}
\left| f (x) - M _{z} \left(k _{0} \left(x z ^{-1}\right) M _{y} \left(k _{0} \left(z y ^{-1}\right) f (y)\right)\right) \right| \leqq 2 \varepsilon .
\end{equation}

The left hand side is precisely $|f(x) - (T^2 f)(x)| \leqq 2\varepsilon$. 

Since $T \cdot R(P_{\lambda m}) \subseteq R(P_{\lambda m})$, we have proved 

\begin{theorem}\label{thm:XI.10.Theorem.1} The function $f(x)$ can be approximated uniformly on $G$ by linear combinations of the eigenfunctions of $\tilde{T}$ belonging to the eigenvalues which are different from $0$. 
\end{theorem}

We shall take a fixed eigenvalue $\lambda \neq 0$ of $\tilde{T}$, and shall denote the base $\{h_j\} \subseteq C(G)$ of the corresponding eigenspace $P_{\lambda} \cdot \tilde{C}(G)$ by $e_1(x)$, $e_2(x), \ldots, e_k(x)$. Then, by the invariance of the mean value, we obtain, for any $a \in G$, 
\[
\begin{aligned} M _{y} \left(k _{0} (x y ^{-1}) e _{j} (y a)\right) & = M _{y} \left(k _{0} (x a \cdot a ^{-1} y ^{-1}) e _{j} (y a)\right) = M _{z} \left(k _{0} (x a \cdot z ^{-1}) e _{j} (z)\right) \\ & = \left(T e _{j}\right) (x a) = \lambda e _{j} (x a). \end{aligned}
\]

Since the left hand side is equal to the result of applying the operator $T$ upon the function $e_j(ya)$ of $y$, we see that, for any given $a \in G$, the function $e_j(xa)$ of $x$ must be uniquely represented as a linear combination of the functions $e_1(x), e_2(x), \ldots, e_k(x)$. We have thus 
\begin{equation}\label{eq:XI.10.16}
e _{j} (x a) = \sum\limits_{i=1} ^{k} d _{ji} (a) e _{i} (x) \quad (j = 1, 2, \dots , k),
\end{equation}
or, in vectorial notation, 
\begin{equation}\tag{16'}\label{eq:XI.10.16'}
e (x a) = D (a) e (x).
\end{equation}

By $e(x \cdot ab) = D(ab)e(x), e(xa \cdot b) = D(b)e(xa) = D(b)D(a)e(x)$, we see, remembering the linear independence of $e_1(x), e_2(x), \ldots, e_n(x)$, that 
\begin{equation}\label{eq:XI.10.17}
D (a b) = D (b) D (a), D (e) = \text{ the unit matrix of degree } k.
\end{equation}

By applying E. \textsc{Schmidt}'s orthogonalization, we may assume that $\{e_j(x)\}$ constitutes an orthonormal system in $\tilde{C}(G)$. Then, by \eqref{eq:XI.10.16}, 
\begin{equation}\label{eq:XI.10.18}
M _{x} (e _{j} (x a) e _{i} ^{*} (x)) = d _{ji} (a)
\end{equation}
and so the elements $d_{ji}(a)$ of the matrix $D(a)$ belong to $C(G)$. By the invariance of the mean value, we see that 
\[
M _{y} \left(e _{j} (y a) \overline {{e _{i} (y a)}}\right) = M _{y} \left(e _{j} (y) \overline {{e _{i} (y)}}\right) = \delta_{ij}.
\]

Hence the matrix $D(a)$ gives a linear mapping transforming the orthonormal system $\{e_j(x)\}$ onto the orthonormal system $\{e_j(xa)\}$. Therefore $D(a)$ must be a unitary matrix. Hence the tranposed matrix $D(a)'$ of $D(a)$ is also unitary and we have 
\begin{equation}\tag{17'}\label{eq:XI.10.17prime}
D (a b) ^{\prime} = D (a) ^{\prime} D (b) ^{\prime}, D (e) ^{\prime} = \text{ the unit matrix of degree } k.
\end{equation}

Hence $D(a)'$ gives a unitary matrix representation of the group $G$ such that its matrix elements are continuous functions of $a$. Letting $x = e$ in \eqref{eq:XI.10.16}, we see that each $e_j(a)$ is a linear combination of the matrix elements of the representation $D(a)'$. 

We have thus proved 

\begin{theorem}[\textsc{Peter}-\textsc{Weyl}-\textsc{Neumann}]\label{thm:XI.10.Theorem.2} Let $G$ be a totally bounded topological group, metrized by a distance satisfying $\text{ dis }(x, y) = \text{ dis }(axb, ayb)$. Let $f(g)$ be any complex-valued bounded uniformly continuous function defined on $G$. Then, $f(g)$ can be approximated uniformly on $G$ by linear combinations of the matrix elements of unitary uniformly continuous matrix representations $D(g)'$ of the group $G$. 
\end{theorem}

Referring to A. Weil's reduction given in Chapter~\ref{chap:VIII}, \ref{sec:VIII.5}, we obtain the following 

\begin{corollaryu} Let $G$ be an abstract group, and $f(g)$ an almost periodic function on $G$. Then $f(g)$ can be approximated uniformly on $G$ by linear combinations of the matrix elements of unitary matrix representations $D(g)'$ of the group $G$. 
\end{corollaryu}

\begin{remarkn}[Remark $1$] Let the degree of a unitary matrix representation $D(g)'$ of the group $G$ be $d$, i.e., the degree of the matrix $D(g)'$ be $d$. Then each $D(g)'$ gives a linear mapping of a fixed $d$-dimensional complex Hilbert space $X_d$ onto itself. The representation $D(g)'$ is said to be irreducible if there is no proper linear subspace $\neq \{0\}$ of $X_d$ which is invariant by applying the mappings $D(g)'$, $g \in G$. Otherwise, the representation 
$D(g)'$ is called reducible, and there exists a proper linear subspace $X_{d,1} \neq \{0\}$ invariant by every $D(g)'$, $g \in G$. Then the orthogonal complement $X_{d,1}^{\perp}$ of $X_{d,1}$ in $X_d$ is, by the unitarity of the representation $D(g)'$, also invariant by every $D(g)$, $g \in G$. If we take for the base of the linear space $X_d$ the orthonormal system of vectors composed of one orthonormal system of vectors of $X_{d,1}$ and one of $X_{d,1}^{\perp}$, then, by this choice of an orthonormal base of $X_d$, the representation $D(g)'$ will be transformed into the form 
\[
UD(g)'U^{-1}=\begin{pmatrix}D_1(g)'&0\\0&D_2(g)'\end{pmatrix},
\]
where $U$ is a fixed unitary matrix.
\end{remarkn}

Hence a reducible unitary representation $D(g)'$ is completely reducible into the sum of two unitary representations $D_1(g)'$ and $D_2(g)'$ of the group $G$ acting respectively upon $X_{d,1}$ and $X_{d,1}'$. In this way, we finally can choose a fixed unitary matrix $U_d$ such that the representation $U_dD(g)'U_d^{-1}$ is the sum of irreducible unitary representations of the group $G$. Therefore, in the statements of Theorem~\ref{thm:XI.10.Theorem.2} and its Corollary, we can impose the condition that the matrix representations $D(g)'$ are all irreducible. 

\begin{remarkn}[Remark $2$] In the particular case when $G$ is the additive group of real numbers, a unitary irreducible representation $D(g)'$ is given by 
\begin{equation}\label{eq:XI.10.19}
D (g) ^{\prime} = e ^{i \alpha g}, \text{ where } \alpha \text{ is a real number and } i = \sqrt {- 1}.
\end{equation}
\end{remarkn}

For, by the commutativity of the unitary matrices $D(g)'$, $g \in G$, the representation $D(g)'$ is completely reducible into the sum of one-dimensional unitary representations $\chi(g)$, that is, complex-valued solutions $\chi(g)$ of 
\begin{equation}\label{eq:XI.10.20}
\chi \left(g _{1} + g _{2}\right) = \chi \left(g _{1}\right) \cdot \chi \left(g _{2}\right), | \chi (g _{1}) | = 1 (g _{1}, g _{2} \in G), \chi (0) = 1.
\end{equation}

It is well known that any continuous solution of \eqref{eq:XI.10.20} is of the form $\chi(g) = e^{i\alpha g}$. Hence, any continuous almost periodic function $f(g)$ on the additive group $G$ of real numbers can be approximated uniformly on $G$ by linear combinations of $e^{i\alpha g}$ with real $\alpha$'s. This constitutes the fundamental theorem in H. \textsc{Bohr}'s theory of almost periodic functions. According to the original definition by Bohr, a continuous function $f(x)$ ( $-\infty < x < \infty$) is called almost periodic if for each $\varepsilon > 0$ there exists a positive number $p = p(\varepsilon)$ such that any interval of the form $(t, t + p)$ contains at least one $\tau$ such that 
\[
\left| f (x + \tau) - f (x) \right| \leqq \varepsilon \quad \text{ for } \quad - \infty < x < \infty .
\]

See H. \textsc{Bohr} \cite{bohr25}. S. \textsc{Bochner} \cite{bochner23} has shown that a continuous function $f(x) (-\infty < x < \infty)$ is almost periodic in Bohr's sense iff the following condition is satisfied: For any sequence of real numbers $\{a_n\}$, the system of functions $\{f_{a_{n}}(x); f_{a_{n}}(x) = f(x + a_n)\}$ is totally bounded in the topology of the uniform convergence on $(-\infty, \infty)$. It was extended by J. von \textsc{Neumann} \cite{neumann231} to almost periodic functions in a group. Neumann's result contains as a special case, the Peter-Weyl theory of continuous representation of a compact Lie group (\textsc{Peter}-\textsc{Weyl} \cite{peter247}). According to our treatment, Bohr's result is easily proved by observing that $\lim |s - t| = 0$ implies $\lim \left[\sup\limits_{a,b} |f(asb) - f(atb)|\right] = 0$. 


\yossection{11}{Tannaka's Duality Theorem for Non-commutative Compact Groups}
\label{sec:XI.11}

Let $G$ be a compact (topological) group. This means that $G$ is a compact topological space as well as a group in such a way that the mapping 
\[
(\stackrel {\bullet} {x}, y) \rightarrow x y ^{-1}
\]
of the product space $G \times G$ onto $G$ is continuous. 

\begin{proposition}\label{prop:XI.11.Proposition.1} $A$ complex-valued continuous function $f(g)$ defined on a compact group $G$ is uniformly continuous in the following sense: 
\begin{equation}\label{eq:XI.11.1}
\parbox{0.92\linewidth}{\everymath{}\let\limits\nolimits for any $\varepsilon > 0$, there exists a neighbourhood $U(e)$ of the unit element $e$ of $G$ such that $|f(x) - f(y)| < \varepsilon$ whenever $xy^{-1} \in U(e)$ and also whenever $x^{-1}y \in U(e)$.}
\end{equation}
\end{proposition}


\begin{proof}
 Since $f(x)$ is continuous at every point $a \in G$, there exists a neighbourhood $V_{a}$ of $a$ such that $x \in V_{a}$ implies $|f(x) - f(a)| < \varepsilon / 2$. If we denote by $U_{a}$ the neighbourhood of $e$ defined by $U_{a} = V_{a}a^{-1} = \{va^{-1}; v \in V\}$, then $xa^{-1} \in U_{a}$ implies $|f(x) - f(a)| < \varepsilon / 2$. Let us denote by $W_{a}$ the neighbourhood of $e$ such that $W_{a}^{2} \subseteq U_{a}$, where $W_{a}^{2} = \{w_{1}w_{2}; w_{i} \in W_{a}(i = 1,2)\}$. Obviously, the system of all open sets of the form $W_{a} \cdot a$, where $a$ is an arbitrary element of $G$, covers the whole space $G$. $G$ being compact, there exists a finite set $\{a_{i}; i = 1,2,\ldots,n\}$ such that the system of open sets $W_{a_{i}} \cdot a_{i} (i = 1,2,\ldots,n)$ covers $G$. We denote by $U(e)$ the intersection of all open sets of the system $\{W_{a_{i}}\}$. Then $U(e)$ is a neighbourhood of $e$. We shall show that if $xy^{-1} \in U(e)$, then $|f(x) - f(y)| \leqq \varepsilon$. Since the system $W_{a_{i}} \cdot a_{i}$ covers $G$, there exists a number $k$ such that $ya_{k}^{-1} \in W_{a_{k}} \subseteq U_{a_{k}}$ and therefore $|f(y) - f(a_{k})| < \varepsilon / 2$. Furthermore we have $xa_{k}^{-1} = xy^{-1}ya_{k}^{-1} \in U(e) W_{a_{k}} \subseteq W_{a_{k}}^{2} \subseteq U_{a_{k}}$ so that $|f(x) - f(a_{k})| < \varepsilon / 2$. Combining these two inequalities we get $|f(x) - f(y)| < \varepsilon$. 

If we start with a neighbourhood $U_{a}$ of $e$ such that $x^{-1}a \in U_{a}$ implies $|f(x) - f(a)| < \varepsilon / 2$, we would obtain a neighbourhood $U(e)$ of $e$ such that $|f(x) - f(y)| < \varepsilon$ whenever $x^{-1}y \in U(e)$. Thus taking the intersection of these two $U(e)$'s as the $U(e)$ in the statement of the Proposition, we complete the proof. 
\end{proof}

\begin{corollaryu} $A$ complex-valued continuous function $f(x)$ on a compact group $G$ is almost periodic on $G$. 
\end{corollaryu}

\begin{proof}
 Let $U(e)$ be the neighbourhood of $e$ given in Proposition~\ref{prop:XI.11.Proposition.1}. For any $a \in G$, $U(e)a$ is a neighbourhood of $a$. The compact space $G$ is covered by the system of open sets $U(e)a$, $a \in G$, and therefore some finite subsystem $\{U(e)a_i; i = 1, 2, \ldots, n\}$ covers $G$. That is, for any $a \in G$, there exists some $a_k$ with $k \leqq n$ such that $aa_k^{-1} \in U(e)$. Hence, by $(ax)(a_kx)^{-1} = aa_k^{-1}$, we have $\sup\limits_{x}|f(ax) - f(a_kx)| < \varepsilon$. Similarly, we can find a finite system $\{b_j, j = 1, 2, \ldots, m\}$ such that, for any $b \in G$, there exists some $b_j$ with $j \leqq m$ satisfying the inequality 
\[
\sup\limits _{i,x} | f (a _{i} x b) - f (a _{i} x b _{j}) | < \varepsilon . \]

Therefore, for any pair $a, b$ of elements $\in G$, we can find $a_{k}$ and $b_{j}$ ( $k \leqq n, j \leqq m$) such that 
\[
\sup\limits _{x} | f (a x b) - f \left(a _{k} x b _{j}\right) | < 2 \varepsilon .
\]

This proves that the system of functions $\{f_{a,b}(x); f_{a,b}(x) = f(axb), a \text{ and } b \in G\}$ is totally bounded by the maximum norm $\|h\| = \sup\limits_x |h(x)|$. Hence $f(x)$ is almost periodic on $G$. 
\end{proof}

We are now able to extend the Peter-Weyl-Neumann Theory of the preceding section to complex-valued continuous functions $f(x)$ on a compact group $G$. For such a function $f(x)$ and $\varepsilon > 0$, we set 
\[
V = \left\{y \in G; \sup\limits _{x} | f (x) - f (y ^{-1} x) | < \varepsilon \right\}.
\]

Then, by the continuity of $f$, the set $V$ is an open set containing $e$. By Urysohn's theorem as applied to the normal space $G$, we see that there exists a continuous function $k_{1}(x)$ defined on $G$ such that 
\[
0 \leqq k (x) \leqq 1 \text{ on } G, k _{1} (e) = 1 \quad \text{ and } \quad k _{1} (x) = 0 \text{ whenever } x \in V ^{C}.
\]

Then the continuous function 
\begin{equation}\label{eq:XI.11.2}
k (x) = 2 ^{-1} \left(k _{1} (x) + k _{1} \left(x ^{-1}\right)\right)
\end{equation}
satisfies the condition $k(x^{-1}) = k(x)$ and 
\[
\begin{yoslines}
&0 \leqq k (x) \leqq 1 \text{ on } G, k (e) = 1 \text{ and } k (x) = 0 \text{ whenever } x \in U ^{C}, \text{ where } U = V \cup V ^{-1}.
\end{yoslines}
\]

Therefore, if we denote by $C(G)$ the $B$-space of all complex-valued continuous functions $h(x)$ defined on $G$ normed by the maximum norm, we can define a linear operator $T$ on $\hat{C}(G)$ into $\hat{C}(G)$ by 
\begin{equation}\label{eq:XI.11.3}
(T h) (x) = \left(k _{0} \times h\right) (x), x \in G.
\end{equation}

Here, as in the preceding section, 
\begin{equation}\label{eq:XI.11.4}
k _{0} (x) = k (x) / M _{x} (k (x))
\end{equation}
and $\hat{C}(G)$ is the space $C(G)$ endowed with the scalar product 
\begin{equation}\label{eq:XI.11.5}
(b, h) = M _{y} (b (y) h (y)) = (b \times h ^{*}) (e), h ^{*} (y) = \overline {{{h (y ^{-1})}}}.
\end{equation}

Thus, as in the preceding section, we obtain 

\begin{theorem}\label{thm:XI.11.Theorem.1} Any complex-valued continuous function $f(g)$ defined on a compact group $G$ is almost periodic, and $f(g)$ can be approximated uniformly on $G$ by linear combinations of the matrix elements of unitary, continuous, irreducible matrix representations of $G$. 
\end{theorem}

We shall say that two matrix representations $A_{1}(g)$ and $A_{2}(g)$ of a group $G$ are equivalent if there exists a fixed non-singular matrix $B$ such that $B^{-1}A_{1}(g)B = A_{2}(g)$ for all $g \in G$. 

\begin{proposition}[I. \textsc{Schur}'s lemma]\label{prop:XI.11.Proposition.2} If the representations $A_{1}(g), A_{2}(g)$ are irreducible and inequivalent, then there is no matrix $B$ such that 
\begin{equation}\label{eq:XI.11.6}
A _{1} (g) B = B A _{2} (g)
\end{equation}
holds identically in $g$, except $B = 0$. In \eqref{eq:XI.11.6}, the matrix $B$ is assumed to be of $n_1$ rows and $n_2$ columns, where $n_1, n_2$ are the degrees of $A_1(g), A_2(g)$, respectively. 
\end{proposition}

\begin{proof}
 Let $X_{1}$ and $X_{2}$ be the linear spaces subject to the linear transformations $A_{1}(g)$ and $A_{2}(g)$, respectively. $B$ in \eqref{eq:XI.11.6} can be interpreted as a linear mapping $x_{2} \to x_{1} = Bx_{2}$ of $X_{2}$ onto $X_{1}$. The linear subspace of $X_{1}$ consisting of all vectors $x_{1}$ of the form $Bx_{2}$ is invariant, for $A_{1}(g)x_{1} = Bx_{2}'$ with $x_{2}' = A_{2}(g)x_{2}$. By virtue of the irreducibility of $A_{1}(g)$, there are only two possibilities: either $Bx_{2} = 0$ for all $x_{2}$ in $X_{2}$, i.e. $B = 0$, or $X_{1} = BX_{2}$. On the other hand, the set of all vectors $x_{2}$ in $X_{2}$ such that $Bx_{2} = 0$ is an invariant subspace of $X_{2}$, for $BA_{2}(g)x_{2} = A_{1}(g)Bx_{2} = 0$. From the irreducibility of $A_{2}(g)$ we conclude: either $Bx_{2} = 0$ for all $x_{2}$ in $X_{2}$, i.e. $B = 0$, or $x_{2} = 0$ is the only vector in $X_{2}$ such that $Bx_{2} = 0$, so that distinct vectors in $X_{2}$ go into distinct vectors in $X_{1}$ under the linear mapping $B$. Hence if $B \neq 0$ we conclude that $B$ defines a one-one linear mapping of $X_{2}$ onto $X_{1}$. But this means that $B$ is a non-singular matrix ( $n_1 = n_2$) and so $A_{1}(g)$ and $A_{2}(g)$ would be equivalent. 
\end{proof}

\begin{proposition}[the orthogonality relations]\label{prop:XI.11.Proposition.3} Let $A_{1}(g) = (a_{ij}^{1}(g))$ and $A_{2}(g) = (a_{kl}^{2}(g))$ be unitary, continuous, irreducible matrix representations of a compact group $G$. Then we have the orthogonality relations: 
\begin{align}
M _{g} (a _{ij} ^{1} (g) \overline {{a _{kl} ^{2} (g)}})& = 0, \text{ if } A _{1} (g) \text{ is inequivalent to } A _{2} (g), \label{eq:XI.11.7} \\
M _{g} (a _{ij} ^{1} (g) \overline {{a _{kl} ^{1} (g)}})& = n _{1} ^{-1} \delta_{ik} \delta_{jl} \text{ where } n _{1} \text{ is the degree of } A _{1} (G). \notag
\end{align}
\end{proposition}

\begin{proof}
 Let $n_{1}, n_{2}$ be the degrees of $A_{1}(g), A_{2}(g)$ respectively. Take any matrix $B$ of $n_{1}$ rows and $n_{2}$ columns, and set $A(g) = A_{1}(g) BA_{2}(g^{-1})$. 

Then the matrix $A = M_g(A(g))$ satisfies $A_1(g)A = AA_2(g)$. For, we have, by the invariance of the mean value, 
\[
\begin{gathered}
A _{1} (y) A A _{2} (y ^{-1}) = M _{g} (A _{1} (y) A _{1} (g) B A _{2} (g ^{-1}) A _{2} (y ^{-1})) = M _{g} (A _{1} (y g) B A _{2} ((y g) ^{-1}) = A.
\end{gathered}
\]

By Schur's lemma, $A$ must be equal to a zero matrix. If we take $B = (b_{jl})$ in such a way that only the element $b_{jl}$ is not zero, then, by the unitarity condition $A_{2}(g^{-1}) = \overline{A_{2}(g)}'$, we obtain 
\[
M _{g} (a _{ij} ^{1} (g) \overline {{a _{kl} ^{2} (g)}}) = 0.
\]


Next we have, as above, $A_{1}(g)A = AA_{1}(g)$ for $A = M_{g}(A_{1}(g)BA_{1}(g^{-1}))$. Let $\alpha$ be any one of the eigenvalues of the matrix $A$. Then the matrix $(A - \alpha I_{n_1})$, where $I_{n_1}$ denotes the unit matrix of degree $n_1$, satisfies 
\[
A _{1} (g) \left(A - \alpha I _ {n _{1}}\right) = \left(A - \alpha I _ {n _{1}}\right) A _{1} (g).
\]

Therefore, by Schur's lemma, the matrix $(A - \alpha I_{n_1})$ is either non-singular, or $(A - \alpha I_{n_1}) = 0$. The first possibility is excluded since $\alpha$ is an eigenvalue of $A$. Thus $A = \alpha I_{n_1}$. By taking the trace (the sum of the diagonal elements) of both sides of $A = M_g(A_1(g)BA_1(g^{-1}))$, we obtain 
\[
\begin{aligned} n _{1} \alpha & = \tr (A) = M _{g} (\tr (A _{1} (g) B A _{1} (g) ^{-1}) = M _{g} (\tr (B)) \\ & = \tr (B). \end{aligned}
\]

Therefore, if we take $B = (b_{jl})$ in such a way that $b_{jl} = 1$ while the other elements are all zero, then, from $M_g(A_1(g)BA_1(g^{-1})) = n_1^{-1}$ trace $(B)I_{n_1}$, we obtain 
\[
M _{g} \left(a _{ij} ^{1} (g) \overline {{{a _{kl} ^{1} (g)}}}\right) = n _{1} ^{-1} \delta_{ik} \delta_{jl}.
\]
\end{proof}

\begin{corollaryu} There exists a set $U$ of mutually inequivalent, continuous, unitary, irreducible matrix representations $U(g) = (u_{ij}(g))$ of $G$ satisfying the following three conditions: 

\begin{enumerate}[i)]
\item for any pair of distinct points $g_1, g_2$ of $G$, there exists an $U(g) \in \mathfrak{U}$ such that $U(g_1) \neq U(g_2)$,
\item if $U(g) \in \mathfrak{U}$, then the complex conjugate representation $\bar{U}(g)$ of $U(g)$ also belongs to $\mathfrak{U}$,
\item if $U_{1}(g), U_{2}(g)$ are in $\mathfrak{U}$, then the product representation $U_{1}(g) \times U_{2}(g)$, explained below, is completely reducible to a sum of a finite number of representations $\in \mathfrak{U}$.
\end{enumerate}
\end{corollaryu}

\begin{proof}
 The definition of the product representation $U_{1}(g) \times U_{2}(g)$ is given as follows. Let $(e_{1}^{1}, e_{2}^{1}, \ldots, e_{n}^{1})$ and $(e_{1}^{2}, e_{2}^{2}, \ldots, e_{m}^{2})$ be orthonormal bases respectively of the finite dimensional complex Hilbert spaces subject to the linear mappings $U_{1}(g)$ and $U_{2}(g)$, respectively. The product space of these two Hilbert spaces is spanned by the product base consisting of $nm$ vectors $e_i^1 \times e_j^2$ ( $i = 1, 2, \ldots, n; j = 1, 2, \ldots, m$). Referring to this base, the product representation $U_1(g) \times U_2(g)$ of $U_1(g) = (u_{ij}^1(g))$ and $U_2(g) = (u_{kl}^2(g))$ is given by 
\[
\left(U _{1} (g) \times U _{2} (g)\right) \left(e _{i} ^{1} \times e _{j} ^{2}\right) = \sum\limits_{s,t} u _{si} ^{1} (g) u _{tj} ^{2} (g) \left(e _{i} ^{1} \times e _{j} ^{2}\right). \]

Let us take a maximal set $U$ of mutually inequivalent, continuous, unitary, irreducible matrix representations $U(g)$ satisfying the condition ii). Then, by Theorem~\ref{thm:XI.11.Theorem.1}, condition i) is satisfied. Condition iii) is also satisfied since the product representation of two unitary representations is also unitary and hence is completely reducible. 
\end{proof}

We are now ready to formulate T. \textsc{Tannaka}'s duality theorem. Let $R$ be the set of all Fourier polynomials: 
\begin{equation}\label{eq:XI.11.8}
x (g) = \sum \gamma_{ij} ^{(\alpha)} u _{ij} ^{(\alpha)} (g),
\end{equation}
that is, finite linear combinations of $u_{ij}^{(\alpha)}(g)$, where $(u_{ij}^{(\alpha)}(g)) \in \mathfrak{U}$ and $\gamma_{ij}^{(\alpha)}$ denote complex numbers. Then $\operatorname{Re}$ is a ring with the multiplicative unit $u(u(g) \equiv 1 \text{ on } G)$ and with complex multipliers; the sum and the multiplication in the ring $\operatorname{Re}$ being understood as the function sum and the function multiplication, respectively. Let $\mathfrak{T}$ be the set of all linear homomorphisms $T$ of the ring $\operatorname{Re}$ onto the field of complex numbers such that 
$Tu = 1, T\bar{x} = \overline{Tx}$, the bar indicating the complex-conjugate. 


$\mathfrak{T}$ is not void since each $g \in G$ induces such a homomorphism $T_g$: 
\begin{equation}\label{eq:XI.11.9}
T _{g} x = x (g).
\end{equation}

By the condition i) for $\mathfrak{U}$, we see that 
\begin{equation}\label{eq:XI.11.10}
g _{1} \neq g _{2} \quad \text{ implies } \quad T _ {g _{1}} \neq T _ {g _{2}}.
\end{equation}

\begin{proposition}\label{prop:XI.11.Proposition.4} $\mathfrak{T}$ may be considered as a group which contains $G$ as a subgroup. 
\end{proposition}

\begin{proof}
 We shall define a product $T_{1} \cdot T_{2} = T$ in $T$ as follows. Let 
\[
U ^{(\alpha)} (g) = \left(u _{ij} ^{(\alpha)} (g)\right) \quad (i, j = 1, 2, \dots , n)
\]
be members of $\mathfrak{U}$. We put, for $u_{ij}^{(\alpha)}(gh) = \sum\limits_{k} u_{ik}^{(\alpha)}(g) u_{kj}^{(\alpha)}(h)$, 
\begin{equation}\label{eq:XI.11.11}
T u _{ij} ^{(\alpha)} = \sum\limits_{k} T _{1} u _{ik} ^{(\alpha)} \cdot T _{2} u _{kj} ^{(\alpha)}. \end{equation}

By the orthogonal relations \eqref{eq:XI.11.7}, we see that the functions $u_{ij}^{(\alpha)}(g)$'s, where $(u_{ij}^{(\alpha)}(g)) \in \mathfrak{U}$, are linearly independent on $G$. Hence, we see that $T$ may be extended linearly on the whole $\mathfrak{R}$. It is easy to see that the extension $T$ is also a member of $\mathfrak{T}$ and that $\mathfrak{T}$ is a group with the unit $T_e$ ( $e =$ the unit of $G$) and $T_g^{-1} = T_{g^{-1}}$. $G$ is isomorphically embedded in this group $\mathfrak{T}$ by the correspondence $g \leftrightarrow T_g$, as will be seen from \eqref{eq:XI.11.11} and \eqref{eq:XI.11.10}. 
\end{proof}

In truth, we have 

\begin{theorem}[T. \textsc{Tannaka}]\label{thm:XI.11.Theorem.2} $\mathfrak{T} = G$, viz. every $T \in \mathfrak{T}$ is equal to a certain $T_g$: 
\begin{equation}\label{eq:XI.11.12}
T x = x (g) \quad \text{ for all } \quad x \in \mathfrak {R}.
\end{equation}
\end{theorem}

\begin{proof}
 We introduce a weak topology in the group $T$ by taking the sets of the form 
\begin{equation}\label{eq:XI.11.13}
\{T \in \mathfrak {T}; | T x _{i} - T _{e} x _{i} | < \varepsilon_{i} (i = 1, 2, \dots , n) \}
\end{equation}
as neighbourhoods of the unit $T_{e}$ of the group $\mathfrak{T}$. Then $\mathfrak{T}$ is a compact space. For, by $\sum\limits_{s}|T u_{si}^{(\alpha)}(g)|^2 = (T \cdot T)(1) = 1$ implied by ($8$) and \eqref{eq:XI.11.11}, 
$\mathfrak{T}$ is a closed subset of the topological product of compact spaces: 
\[
\prod\limits_ {x \in \mathfrak {R}} \left\{z; | z | \leqq \sup\limits _ {T \in \mathfrak {T}} | T x | \right\},
\]
and so we can apply Tychonov's theorem. It is easily seen that the isomorphic embedding $g \leftrightarrow T_{g}$ is also a topological one, because a one-one continuous map of a compact space onto a compact space is a topological map. Therefore $G$ may be considered as a closed subgroup of the compact group $T$. 

By the above weak topology of $\mathfrak{T}$, each $x(g) \in \operatorname{Re}$ gives rise to a continuous function $x(T)$ on the compact group $\mathfrak{T}$ such that $x(T_g) = x(g)$. We have only to set $x(T) = T \cdot x$. The set of all these continuous functions $x(T)$ constitutes a ring $\operatorname{Re}(\mathfrak{T})$ of complex-valued continuous functions on $\mathfrak{T}$ satisfying the conditions: 

$1$) $1 = u(T)\in \operatorname{Re} (\mathfrak{T})$ 
$2$) for any pair of distinct points $T_{1}, T_{2}$ of $\mathfrak{T}$, there exists an $x(T) \in \operatorname{Re}(\mathfrak{T})$ such that $x(T_{1}) \neq x(T_{2})$, 
$3$) for any $x(T) \in \operatorname{Re}(\mathfrak{T})$, there exists the complex-conjugate function $\overline{x}(T) = \overline{x(T)}$ which belongs to $\operatorname{Re}(\mathfrak{T})$. 

Now let us suppose that $\mathfrak{T} - G$ is not void. Since a compact space $\mathfrak{T}$ is normal, we may apply Urysohn's theorem to the effect that there exist a point $T_0 \in (\mathfrak{T} - G)$ and a continuous function $y(T)$ on $\mathfrak{T}$ such that 
\begin{equation}\label{eq:XI.11.14}
y (T) \geqq 0 \text{ on } \mathfrak {T}, y (g) = 0 \text{ on } G \text{ and } y (T _{0}) = 1.
\end{equation}

By the Stone-Weierstrass theorem in Chapter~\ref{chap:0}, as applied to the ring $\operatorname{Re}(\mathfrak{T})$ satisfying $1$), $2$) and $3$), we see that there exists, for any $\varepsilon > 0$, a function $x(g) = \sum \gamma_{ij}^{(\alpha)} u_{ji}^{(\alpha)}(g) \in \operatorname{Re}$ such that $|y(T) - \sum \gamma_{ij}^{(\alpha)} u_{ij}^{(\alpha)}(T)| < \varepsilon$ on $\mathfrak{T}$, and, in particular, $|y(g) - \sum \gamma_{ij}^{(\alpha)} u_{ij}^{(\alpha)}(g)| < \varepsilon$ on $G$. Let $u_{11}^{(\alpha_n)}(g) = u(g) = 1$. Then, by taking the mean value and remembering the orthogonality relation \eqref{eq:XI.11.7}, as applied to the compact groups $\mathfrak{T}$ and $G$, we obtain 
\begin{equation}\label{eq:XI.11.15}
\left| M _{T} (y (T)) - \gamma_{11} ^ {\left(\alpha_{0}\right)} \right| < \varepsilon , \left| M _{g} (y (g)) - \gamma_{11} ^ {\left(\alpha_{0}\right)} \right| < \varepsilon .
\end{equation}

We have thus arrived at a contradiction, since $M_T(y(T)) > 0$ and $M_g(y(g)) = 0$ by \eqref{eq:XI.11.14}. 
\end{proof}
\begin{remarkn}[Remark $1$] The above proof of Tannaka's theorem is adapted from K. \textsc{Yosida} \cite{yosida357}. The original proof is given in T. \textsc{Tannaka} \cite{tannaka320}. Since a continuous, unitary and irreducible matrix representation of a compact abelian group $G$ is precisely a continuous function $\chi(g)$ on $G$ satisfying 
\[
\chi \left(g _{1}\right) \chi \left(g _{2}\right) = \chi \left(g _{1} g _{2}\right) \quad \text{ and } \quad | \chi (g) | = 1,
\]
\end{remarkn}

Tannaka's theorem contains, as a special case, the duality theorem of L. \textsc{Pontrjagin} \cite{pontrjagin263}. For further references, see M. A. \textsc{Naimark} \cite{naimark221}. 

\begin{remarkn}[Remark $2$] To define mean values of continuous functions $f_{k}(g)$ ( $k = 1, 2, \ldots, m$) on $G$ by the method given in Chapter~\ref{chap:VIII}, \ref{sec:VIII.5}, we have only to replace the $\text{ dis }(g_1, g_2)$ there by $\text{ dis }(g_1, g_2) = \sup\limits_{g, h \in G; k=1,2,\ldots,m} |f_k(g g_1 h) - f_k(g g_2 h)|$. 
\end{remarkn}


\yossection{12}{Functions of a Self-adjoint Operator}
\label{sec:XI.12}

Let $H = \int \lambda dE(\lambda)$ be the spectral resolution of a self-adjoint operator $H$ in a Hilbert space $X$. For a complex-valued Baire function $f(\lambda)$, we consider the set 
\begin{equation}\label{eq:XI.12.1}
D (f (H)) = \left\{x \in X; \int_{- \infty} ^{\infty} | f (\lambda) | ^{2} \dd \| E (\lambda) x \| ^{2} < \infty \right\},
\end{equation}
where the integral is taken with respect to the Baire measure $m$ determined by $m((\lambda_{1}, \lambda_{2}]) = \|E(\lambda_{2})x\|^{2} - \|E(\lambda_{1})x\|^{2}$. As in the case of a continuous function $f(\lambda)$, treated in Chapter~\ref{chap:XI}, \ref{sec:XI.5}, we see that the integral 
\begin{equation}\label{eq:XI.12.2}
\int_{- \infty} ^{\infty} f (\lambda) \dd (E (\lambda) x, y), x \in D (f (H)), y \in X,
\end{equation}
exists and is finite with respect to the Baire measure $m$ determined by $m((\lambda_{1}, \lambda_{2}]) = (E(\lambda_{2})x, y) - (E(\lambda_{1})x, y)$. Further we see that \eqref{eq:XI.12.2} gives the complex conjugate of a bounded linear functional of $y$. Hence we may write \eqref{eq:XI.12.2} as $(f(H)x, y)$ by virtue of F. \textsc{Riesz}' representation theorem in Chapter~\ref{chap:III}, \ref{sec:III.6}. In this way, we can define functions of $H$: 
\begin{equation}\label{eq:XI.12.3}
f (H) = \int_{- \infty} ^{\infty} f (\lambda) \dd E (\lambda)
\end{equation}
through \eqref{eq:XI.12.1} and \eqref{eq:XI.12.2}. 

\begin{example}\label{ex:XI.12.Example.1} If $H$ is bounded self-adjoint and $f(\lambda) = \sum\limits_{j=1}^{n} \alpha_j \lambda^j$, then, as in Chapter~\ref{chap:XI}, \ref{sec:XI.5}. 
\[
f (H) = \int_{- \infty} ^{\infty} f (\lambda) \dd E (\lambda) = \sum\limits_{j=1} ^{n} \alpha_{j} H ^{j}.
\]
\end{example}

\begin{example}\label{ex:XI.12.Example.2} If $f(\lambda) = (\lambda - i)(\lambda + i)^{-1}$, then $f(H)$ is equal to the Cayley transform $U_H$ of $H$. We have, in this case, $D(f(H)) = X$ since $|f(\lambda)| = 1$ for real $\lambda$. It is easy to see, by $E(\lambda_1)E(\lambda_2) = E(\min (\lambda_1, \lambda_2))$, that if we apply the bounded self-adjoint operator $(H - iI)^{-1} = \int (\lambda - i)^{-1} dE(\lambda)$ to the operator $f(H) = \int f(\lambda) dE(\lambda)$ then the result is equal to $(H + iI)^{-1} = \int (\lambda + i)^{-1} dE(\lambda)$. Hence $f(H) = U_H$. 
\end{example}

\begin{example}\label{ex:XI.12.Example.3} As in Chapter~\ref{chap:XI}, \ref{sec:XI.5}, we have 
\begin{equation}\label{eq:XI.12.4}
\left\| f (H) x \right\| ^{2} = \int_{- \infty} ^{\infty} | f (\lambda) | ^{2} d \left\| E (\lambda) x \right\| ^{2} \text{ whenever } x \in D (f (H)).
\end{equation}
\end{example}

\begin{definitionu} Let $A$ be a not necessarily bounded linear operator, and $B$ a bounded linear operator in a Hilbert space. If 
\begin{equation}\label{eq:XI.12.5}
x \in D (A) \quad \text{ implies } \quad B x \in D (A) \quad \text{ and } \quad A B x = B A x,
\end{equation}
that is, if $AB \supseteq BA$, then we shall write $B \in (A)'$ and say that $B$ is commutative with $A$. We shall thus write the totality of bounded linear operators $B$ commutative with $A$ by $(A)'$. 
\end{definitionu}

\begin{theorem}\label{thm:XI.12.Theorem.1} For a function $f(H)$ of a self-adjoint operator $H = \int \lambda dE(\lambda)$ in a Hilbert space $X$, we have 
\begin{equation}\label{eq:XI.12.6}
(f (H)) ^{\prime} \supseteq (H) ^{\prime},
\end{equation}
that is, $f(H)$ is commutative with every bounded linear operator which commutes with $H$. (Since $E(\lambda) \in (H)'$ by Theorem~\ref{thm:XI.5.Theorem.2} in Chapter~\ref{chap:XI}, \ref{sec:XI.5}, we see, in particular, that $f(H)$ is commutative with every $E(\lambda)$.) 
\end{theorem}

\begin{proof}
 Suppose $S \in (H)'$. Then we can show that $S$ is commutative with every $E(\lambda)$. We first show that $S$ is commutative with the Cayley transform $U_{H}$ of $H$. For, if $x \in D(H)$, then we have, by $S \in (H)'$, 
\[
S (H + i I) x = (H + i I) S x, (H - i I) S x = S (H - i I) x. \]

By putting $(H + iI)x = y$, we see, from the first of the above relations, that 
\[
(H + i I) ^{-1} S y = S (H + i I) ^{-1} y \quad \text{ for all } \quad y \in X = R (H + i I).
\]

Hence 
\[
S (H - i I) (H + i I) ^{-1} = (H - i I) (H + i I) ^{-1} S, \text{ that is, } S U _{H} = U _{H} S.
\]

Therefore, $S$ is commutative with $(U_H)^n = \int_0^{2\pi} e^{ni\theta} dF(\theta) (n = 0, \pm 1, \ldots)$ and so 
\[
\int_{0} ^{2 \pi} e ^{n i \theta} d (S F (\theta) x, y) = \left(S \int_{0} ^{2 \pi} e ^{n i \theta} \dd F (\theta) x, y\right) = \int_{0} ^{2 \pi} e ^{n i \theta} \dd (F (\theta) S x, y).
\]

Hence, as in the proof of the uniqueness of the spectral resolution of a unitary operator, given in Chapter~\ref{chap:XI}, \ref{sec:XI.4}, we obtain $SF(\theta) = F(\theta)S$. This proves $SE(\lambda) = E(\lambda)S$ since $E(-\cot \theta) = F(\theta)$. Thus, for $x \in D(f(H))$,  
we obtain 
\[
(S \int f (\lambda) \dd E (\lambda) x, y) = \int f (\lambda) d (S E (\lambda) x, y) = \int f (\lambda) \dd (E (\lambda) S x, y),
\]
that is, $Sf(H) \subseteq f(H) S$. 
\end{proof}

The above Theorem~\ref{thm:XI.12.Theorem.1} admits a converse in the form of 

\begin{theorem}[\textsc{Neumann}-\textsc{Riesz}-\textsc{Mimura}]\label{thm:XI.12.Theorem.2} Let $H$ be a self-adjoint operator in a separable Hilbert space $X$. Let $T$ be a closed linear operator in $X$ such that $D(T)^{a} = X$. Then a necessary and sufficient condition in order that $T$ be a function $f(H)$ of $H$ with an everywhere finite Baire function $f(\lambda)$ is that 
\begin{equation}\label{eq:XI.12.7}
(T) ^{\prime} \supseteq (H) ^{\prime}.
\end{equation}
\end{theorem}

\begin{proof}
 We have only to show that condition \eqref{eq:XI.12.7} is sufficient. We may assume that $H$ is a bounded self-adjoint operator. If $H$ is not bounded, then we consider $H_{1} = \tan^{-1}H$. Since $|\tan^{-1}\lambda| \leqq \pi/2$, we easily see that $H_{1}$ is bounded self-adjoint. By Theorem~\ref{thm:XI.12.Theorem.1}, the operator $H = \tan H_{1}$ is commutative with every operator $\in (H_{1})'$. Thus, by hypothesis, we have 
\[
(T) ^{\prime} \supseteq (H) ^{\prime} \supseteq (H _{1}) ^{\prime}. \]

Therefore, if Theorem~\ref{thm:XI.12.Theorem.2} is proved for bounded $H_{1}$, then $T = f_{1}(H_{1}) = f_{1}(\tan^{-1}H) = f_{2}(H)$, where $f_{2}(\lambda) = f_{1}(\tan^{-1}\lambda)$. 

We may thus assume that $H$ is bounded and self-adjoint. 

The first step. For any fixed $x_0 \in D(T)$, we can find a Baire function $F(\lambda)$ such that $Tx_0 = F(H)x_0$. This may be proved as follows. Let $M(x_0)$ be the smallest closed linear subspace of $X$ spanned by $x_0, Hx_0, H^2x_0, \ldots, H^n x_0, \ldots$. We denote by $L$ the projection upon $M(x_0)$. Then $(T)' \ni L$. For, by $HM(x_0) \subseteq M(x_0)$, we obtain $HL = LHL$ and so $LH = (HL)^* = (LHL)^* = LHL = HL$, that is, $L \in (H)'$. Hence, by hypothesis, $(T)' \ni L$. 

Therefore $Tx_{0}=TLx_{0}=LTx_{0}\in M(x_{0})$, and so there exists a sequence $\{\phi_{n}(\lambda)\}$ of polynomials such that 
\begin{equation}\label{eq:XI.12.8}
T x _{0} = \slim\limits_{n \rightarrow \infty} p _{n} (H) x _{0}.
\end{equation}

Hence, by \eqref{eq:XI.12.4}, we obtain 
\[
\left\| p _{n} (H) x _{0} - p _{m} (H) x _{0} \right\| ^{2} = \int_{- \infty} ^{\infty} \left| p _{n} (\lambda) - p _{m} (\lambda) \right| ^{2} d \left\| E (\lambda) x _{0} \right\| ^{2}.
\]

As in the proof of the completeness of $L^{2}(S, \mathfrak{B}, m)$, we see that there exists a Baire function $F(\lambda)$ which is square-integrable with respect to the measure $m$ determined by $m((\lambda_{1}, \lambda_{2}]) = \|E(\lambda_{2}) x_{0}\|^{2} - \|E(\lambda_{1}) x_{0}\|^{2}$ such that 
\[
\lim\limits _{n \rightarrow \infty} \int_{- \infty} ^{\infty} | F (\lambda) - p _{n} (\lambda) | ^{2} \dd \| E (\lambda) x _{0} \| ^{2} = 0.
\]

Hence, for $F(H)$, we have 
\[
\lim\limits _{n \rightarrow \infty} \| (F (H) - p _{n} (H)) x _{0} \| ^{2} = \lim\limits _{n \rightarrow \infty} \int_{- \infty} ^{\infty} | F (\lambda) - p _{n} (\lambda) | ^{2} \dd \| E (\lambda) x _{0} \| ^{2} = 0.
\]

This proves that $Tx_0 = F(H)x_0$. Since $F(\lambda)$ is finite almost everywhere with respect to the measure $m$ determined by $m((\lambda_1, \lambda_2]) = \|E(\lambda_2)x_0\|^2 - \|E(\lambda_1)x_0\|^2$, we may assume that $F(\lambda)$ is a Baire function which takes a finite value at every $\lambda$; we may put $F(\lambda) = 0$ for those $\lambda$ for which $|F(\lambda)| = \infty$. 

The second step. Since $X$ is separable and $D(T)^a = X$, we may choose a countable sequence $\{g_n\} \subseteq D(T)$ such that $\{g_n\}$ is strongly dense in $X$. We set 
\begin{equation}\label{eq:XI.12.9}
f _{1} = g _{1}, f _{2} = g _{2} - L _{1} g _{2}, \dots , f _{n} = g _{n} - \sum\limits_{k=1} ^{n-1} L _{k} g _{n}, \text{ where } L _{k}
\end{equation}
is the projection on the closed linear subspace $M(f_k)$. 


By the first step, we have $(T)' \ni L_k$ and so 
\begin{equation}\label{eq:XI.12.10}
L _{k} g _{n} \in D (T) \text{ which implies } f _{n} \in D (T) \quad (n = 1, 2, \dots).
\end{equation}

We may show that 
\begin{equation}\label{eq:XI.12.11}
L _{i} L _{k} = 0 \quad (\text{ for } i \neq k)
\end{equation}
and 
\begin{equation}\label{eq:XI.12.12}
I = \sum\limits_{k=1} ^{\infty} L _{k}.
\end{equation}

Suppose that \eqref{eq:XI.12.11} is proved for $i, k < n$. Then, for $i < n$, 
\begin{equation}L _{i} f _{n} = L _{i} g _{n} - L _{i} \left(\sum\limits_{k=1} ^{n-1} L _{k} g _{n}\right) = L _{i} g _{n} - L _{i} ^{2} g _{n} = L _{i} g _{n} - L _{i} g _{n} = 0, L _{i} H ^ {k ^{\prime}} f _{n} = H ^ {k ^{\prime}} L _{i} f _{n} = 0. \notag\end{equation}

Thus $M(f_{n})$ is orthogonal to $M(f_{i})$. This proves $L_{i}L_{n}=L_{n}L_{i}=0$. 

Next we put $\sum\limits_{k=1}^{\infty} L_k = P$ and show that $P g_n = g_n$ ( $n = 1, 2, \ldots$). Then, since $\{g_n\}$ is dense in $X$, we obtain $P = I$. But, by ($9$), we have 
\[
P g _{n} = P f _{n} + \sum\limits_{k=1} ^{n-1} P L _{k} g _{n}.
\]

We also have $Pf_{n} = f_{n}$ by $f_{n} \in M(f_{n})$. Hence, by $PL_{k} = L_{k}$ implied by \eqref{eq:XI.12.11}, we obtain $Pg_{n} = g_{n}$ ( $n = 1, 2, \ldots$). 

The third step. Take a sequence $\{c_n\}$ of positive numbers such that 
\[
\slim _{k \to \infty} \sum\limits_{n=1} ^{k} c _{n} f _{n}, \slim _{k \to \infty} \sum\limits_{n=1} ^{k} c _{n} T f _{n}
\]
both exist. For instance, we may take $c_{n} = 2^{-n} (\|f_{n}\| + \|Tf_{n}\|)^{-1}$. Since $T$ is a closed operator, we have 
\begin{equation}\label{eq:XI.12.13}
x _{0} = \sum\limits_{n=1} ^{\infty} c _{n} f _{n} \in D (T) \text{ and } T x _{0} = y _{0} = \sum\limits_{n=1} ^{\infty} c _{n} T f _{n}.
\end{equation}

Hence, by the first step, 
\begin{equation}\label{eq:XI.12.14}
T x _{0} = F (H) x _{0}.
\end{equation}

Let $B \in (H)'$ be a bounded self-adjoint operator. Then, by hypothesis, $B \in (T)'$. By Theorem~\ref{thm:XI.12.Theorem.1}, $F(H)$ is commutative with $B$. Hence 
\begin{equation}\label{eq:XI.12.15}
F (H) B x _{0} = B F (H) x _{0} = B T x _{0} = T B x _{0}.
\end{equation}

Let $e_n(\lambda)$ be the defining function of the set $\{\lambda; |F(\lambda)| \leqq n\}$, and put 
\[
B = c _{m} ^{-1} P _{n} H ^{k} L _{m}, \quad \text{ where } \quad P _{n} = e _{n} (H).
\]

Then we can show that 
\begin{equation}\label{eq:XI.12.16}
T P _{n} = F (H) P _{n}.
\end{equation}

In fact, we have $L_{m}x_{0} = \sum\limits_{n=1}^{\infty} c_{n}L_{m}f_{n} = c_{m}f_{m}$, by \eqref{eq:XI.12.11} and $f_{m} \in M(f_{m})$. Hence, by \eqref{eq:XI.12.15}, 
\[
\begin{aligned}
&F (H) P _{n} H ^{k} f _{m}  = F (H) c _{m} ^{-1} P _{n} H ^{k} L _{m} x _{0} = F (H) B x _{0} = T B x _{0} = T c _{m} ^{-1} P _{n} H ^{k} L _{m} x _{0} = T P _{n} H ^{k} f _{m},
\end{aligned}
\]
that is, for $h$ spanned by $H^{k}f_{m}$ with fixed $m$, we have 
\begin{equation}\tag{16'}\label{eq:XI.12.16'}
F (H) P _{n} h = T P _{n} h.
\end{equation}

But such $h$'s are dense in $M(f_m)$ and so, by \eqref{eq:XI.12.12}, we see that if we let $m$ take all positive integers then the $h$'s are dense in $X$. Hence we have proved \eqref{eq:XI.12.16} for those $h$'s which are dense in $X$. 

Now $P_{n}$ is bounded by \eqref{eq:XI.12.4}. By the operational calculus given below, the operator $F(H)P_{n}$ is equal to the function $F_{n}(H)$ where 
\[
\begin{aligned}
&F _{n} (\lambda)  = F (\lambda) e _{n} (\lambda) = F (\lambda), \text{ for } | F (\lambda) | \leqq n, = 0, \text{ for } | F (\lambda) | > n.
\end{aligned}
\]

Thus $F_{n}(H) = F(H)P_{n}$ is bounded. 

Let $h^{*} \in X$ be arbitrary, and let $h^{*} = \slim\limits_{j \to \infty} h_{j}$ where $h_{j}^{\prime}s$ are linear combinations of elements of the form $H^{h}f_{m}$. Such a choice of $h_{j}$ is possible as proved above. By the continuity of the operator $F(H)P_{n}$, we have 
\[
F (H) P _{n} h ^{*} = \slim _{j \rightarrow \infty} F (H) P _{n} h _{j}.
\]

Hence, by $\slim\limits_{j \to \infty} P_n h_j = P_n h^*$ and \eqref{eq:XI.12.16}, we see that the closed operator $T$ satisfies \eqref{eq:XI.12.16}. 

The fourth step. Let $y \in D(F(H))$ and set $y_n = P_n y$. Since $F(\lambda)$ is finite everywhere, we must have $\slim\limits_{n \to \infty} P_n = I$. Thus $\slim\limits_{n \to \infty} y_n = \slim\limits_{n \to \infty} P_n y = y$. Hence, by \eqref{eq:XI.12.16}, we obtain $T \supseteq F(H)$. For, $\slim\limits_{n \to \infty} F(H) P_n y = \slim\limits_{n \to \infty} F_n(H) y = F(H) y$ whenever $y \in D(F(H))$. 

Let $y \in D(T)$ and set $y_{n} = P_{n}y$. Then 
\[
\begin{gathered}
T y _{n} = T P _{n} y = F (H) P _{n} y (\text{ by } (1 6)), T P _{n} y = P _{n} T y (\text{ by } P _{n} = e _{n} (H) \in (H) ^{\prime}).
\end{gathered}
\]

By the operational calculus given below, the function $F(H)$ of $H$ is a closed operator. Hence, letting $n \to \infty$ in the above relations, we see that $F(H) \geqq T$. We have thus proved that $T = F(H)$. 
\end{proof}

\begin{theorem}\label{thm:XI.12.Theorem.3} \begin{enumerate}[(i)]
\item Let $\bar{f}(\lambda)$ be the complex-conjugate function of $f(\lambda)$. Then $D(\bar{f}(H)) = D(f(H))$ and, for $x, y \in D(f(H)) = D(\bar{f}(H))$, we have 
\begin{equation}\label{eq:XI.12.17}
(f (H) x, y) = (x, \bar {f} (H) y).
\end{equation}
\item If $x \in D(f(H))$, $y \in D(g(H))$, then 
\begin{equation}\label{eq:XI.12.18}
(f (H) x, g (H) y) = \int_{- \infty} ^{\infty} f (\lambda) \bar {g} (\lambda) \dd (E (\lambda) x, y).
\end{equation}
\item $(\alpha f)(H)x = \alpha f(H)x$ if $x\in D(f(H))$. If $x\in D(f(H))\cap D(g(H))$, then 
\begin{equation}\label{eq:XI.12.19}
(f + g) (H) x = f (H) x + g (H) x.
\end{equation}
\item If $x \in D(f(H))$, then the condition $f(H)x \in D(g(H))$ is equivalent to the condition $x \in D(f \cdot g(H))$, where $f \cdot g(\lambda) = f(\lambda)g(\lambda)$, and we have 
\begin{equation}\label{eq:XI.12.20}
g (H) f (H) x = (g \cdot f) (H) x.
\end{equation}
\item If $f(\lambda)$ is finite everywhere, then $f(H)$ is a normal operator and 
\begin{equation}\label{eq:XI.12.21}
f (H) ^{*} = \bar {f} (H).
\end{equation}

In particular, $f(H)$ is self-adjoint if $f(\lambda)$ is real-valued and finite everywhere.
\end{enumerate}
\end{theorem}

\begin{proof}
\begin{enumerate}[(i)]
\item $D(f(H)) = D(\bar{f}(H))$ is clear and 
\[ \begin{aligned}
&(f (H) x, y)  = \int_{- \infty} ^{\infty} f (\lambda) \dd (E (\lambda) x, y) = \int_{- \infty} ^{\infty} f (\lambda) d (x, E (\lambda) y) = \overline {{(f (H) y , x)}} = (x, \overline {{f}} (H) y).
\end{aligned}
\]
\item We know, by Theorem~\ref{thm:XI.12.Theorem.1}, that $E(\lambda)$ is commutative with $g(H)$. Thus 
\[
\begin{aligned}
(f (H) x, g (H) y) &= \int_{- \infty} ^{\infty} f (\lambda) \dd (E (\lambda) x, g (H) y) = \int_{- \infty} ^{\infty} f (\lambda) d (x, E (\lambda) g (H) y)
\\
&= \int_{- \infty} ^{\infty} f (\lambda) \dd \overline {{(g (H) E (\lambda) y , x)}} = \int_{- \infty} ^{\infty} f (\lambda) \dd \left(\int_{- \infty} ^{\infty} \overline {{g (\mu)}} \dd \overline {{(E (\mu) E (\lambda) y , x)}}\right)
\\
&= \int_{- \infty} ^{\infty} f (\lambda) \dd \left(\int_{- \infty} ^{\lambda} \overline {{g (\mu)}} \dd \overline {{(y , E (\mu) x)}}\right) = \int_{- \infty} ^{\infty} f (\lambda) \overline {{g (\lambda)}} \dd (E (\lambda) x, y).
\end{aligned}
\]
\item is clear.
\item Let $x$ satisfy $\int_{-\infty}^{\infty}|f(\lambda)|^2\dd\|E(\lambda)x\|^2 < \infty$. Then, by $E(\lambda)E(\mu) = E(\min (\lambda ,\mu))$, the condition $\int_{-\infty}^{\infty}|g(\lambda)|^2\dd\|E(\lambda)f(H)x\|^2 < \infty$ implies, by virtue of the commutativity of $E(\lambda)$ with $f(H)$, 
\[
\begin{gathered}
\infty > \int_{- \infty} ^{\infty} | g (\lambda) | ^{2} \dd \| E (\lambda) f (H) x \| ^{2} = \int_{- \infty} ^{\infty} | g (\lambda) | ^{2} \dd \| f (H) E (\lambda) x \| ^{2}
\\
= \int_{- \infty} ^{\infty} | g (\lambda) | ^{2} \dd \left(\int_{- \infty} ^{\infty} | f (\mu) | ^{2} \dd \| E (\mu) E (\lambda) x \| ^{2}\right)
\\
= \int_{- \infty} ^{\infty} | g (\lambda) | ^{2} \dd \left(\int_{- \infty} ^{\lambda} | f (\mu) | ^{2} \dd \| E (\mu) x \| ^{2}\right) = \int_{- \infty} ^{\infty} | g (\lambda) f (\lambda) | ^{2} \dd \| E (\lambda) x \| ^{2}.
\end{gathered}
\]

Since the above calculation may be traced conversely, we see that, under the hypothesis $x \in D(f(H))$, the two conditions $f(H)$ $x \in D(g(H))$ and $x \in D(f \cdot g(H))$ are equivalent and we have 
\[
\begin{aligned}
&(g (H) f (H) x, y)  = \int_{- \infty} ^{\infty} g (\lambda) \dd (E (\lambda) f (H) x, y) = \int_{- \infty} ^{\infty} g (\lambda) \dd \left(\int_{- \infty} ^{\lambda} f (\mu) \dd (E (\mu) x, y)\right) \\
&= \int_{- \infty} ^{\infty} g (\lambda) f (\lambda) \dd (E (\lambda) x, y) = ((g \cdot f) (H) x, y).
\end{aligned}
\]
\item Let us put $h(\lambda)=|f(\lambda)|+\alpha$, $k(\lambda)=h(\lambda)^{-1}$, $g(\lambda)=f(\lambda)h(\lambda)^{-1}$, where $\alpha$ is any positive integer. Then $k(\lambda)$ and $g(\lambda)$ are both bounded functions. Hence $D(k(H))=D(g(H))=X$. Thus, by (iv), 
\begin{equation}\label{eq:XI.12.22}
f (H) = h (H) g (H) = g (H) h (H).
\end{equation}

We have, by (i) and $D(k(H)) = X$, $k(H)^* = k(H)$, i.e., $k(H)$ is self-adjoint. By (iv) we have $x = h(H)k(H)x$ for all $x \in X$ and $x = k(H)h(H)x$ for all $x \in D(h(H))$. Hence $h(H) = k(H)^{-1}$. Thus, by Theorem~\ref{thm:VII.3.Theorem.1} in Chapter~\ref{chap:VII}, \ref{sec:VII.3}, $h(H)$ is self-adjoint. Therefore $D(f(H)) = D(h(H))$ is dense in $X$ and so we may define $f(H)^*$. We shall show that $f(H)^* = \overline{f}(H)$. Let a pair $\{y, y^*\}$ of elements $\in X$ be such that $(f(H)x, y) = (x, y^*)$ for all $x \in D(f(H))$. Then, by $g(H)^* = \overline{g}(H)$ (implied by (i)) and \eqref{eq:XI.12.22}, 
\[
(f (H) x, y) = (g (H) h (H) x, y) = (h (H) x, \bar {g} (H) y).
\]

Hence, by $x \in D(f(H)) = D(h(H))$ and the self-adjointness of $h(H)$, we obtain 
\[
\bar {g} (H) y \in D (h (H)) \quad \text{ and } \quad h (H) \bar {g} (H) y = y ^{*}.
\]

Thus, again by \eqref{eq:XI.12.22}, we obtain $\bar{f}(H)y=y^{*}$ so that $f(H)^{*}=\bar{f}(H)$. Therefore, by (iv), we see that $f(H)$ is normal, that is, $f(H)^{*}f(H)=f(H)f(H)^{*}$. \qedhere
\end{enumerate}
\end{proof}

\begin{corollaryu} If $f(\lambda)$ is finite everywhere, then $f(H)$ is closed. 
\end{corollaryu}

\begin{proof}
 Clear from $f(H)^{**} = \bar{f}(H)^* = \bar{\bar{f}}(H) = f(H)$. 
\end{proof}

A historical note. Theorem~\ref{thm:XI.12.Theorem.2} was first proved by J. \textsc{von} \textsc{Neumann} \cite{neumann234} for the case of a bounded self-adjoint operator $T$. Cf. F. \textsc{Riesz} \cite{rickart276}. The general case of a closed linear operator $T$ was proved by Y. \textsc{Mimura} \cite{mimura201}. The exposition given above is adapted from Y. \textsc{Mimura} \cite{mimura201} and B. Sz. \textsc{Nagy} \cite{nagy217}. 


\yossection{13}{Stone's Theorem and Bochner's Theorem}
\label{sec:XI.13}

As an example of functions of a self-adjoint operator, we give 

\begin{theorem}[M. H. \textsc{Stone}]\label{thm:XI.13.Theorem.1} Let $\{U_t\}$, $-\infty < t < \infty$, be a one-parameter group of class $(C_0)$ of unitary operators in a Hilbert space $X$. Then 
\begin{equation}\label{eq:XI.13.1}
U _{t} = f _{t} (H), \text{ where } f _{t} (\lambda) = \exp (i t \lambda) \text{ and } i H = A, H ^{*} = H,
\end{equation}
is the infinitesimal generator of the group $U_{t}$. 
\end{theorem}


Conversely, for any self-adjoint operator $H$ in $X$, $U_{t} = f_{t}(H)$ defines a one-parameter group of class $(C_0)$ of unitary operators. 

\begin{proof}
 We have, by the representation theorem of the semi-group theory, 
\[
U _{t} x = \slim _{n \to \infty} \exp (t i H (I - n ^{-1} i H) ^{-1}) x. \]

Since the function $g(t) = \exp (ti\lambda (1 - n^{-1}i\lambda)^{-1})$ is smaller than $\exp ((-nt\lambda^2) / (n^2 +\lambda^2))$ in absolute value, we have, for $H = \int \lambda dE(\lambda)$, 
\[
\exp (t i H (I - n ^{-1} i H) ^{-1}) = \int_{- \infty} ^{\infty} \exp \left(\frac {t i \lambda}{1 - n ^{-1} i \lambda}\right) \dd E (\lambda)
\]
and, moreover, 
\[
\begin{aligned}
\lim\limits _{n \rightarrow \infty}& \int_{- \infty} ^{\infty} \left| \exp \left(\frac {t i \lambda}{1 - n ^{-1} i \lambda}\right) - \exp (t i \lambda) \right| ^{2} \dd \| E (\lambda) x \| ^{2}
\\
&= \lim\limits _{n \rightarrow \infty} \int_{- \infty} ^{\infty} \left| \exp \left(\frac {- t \lambda^{2}}{n - i \lambda}\right) - 1 \right| ^{2} \dd \| E (\lambda) x \| ^{2} = 0.
\end{aligned}
\]

This proves that $U_{t} = f_{t}(H) = \int_{-\infty}^{\infty}\exp (it\lambda)dE(\lambda).$ 

For the converse part of Theorem~\ref{thm:XI.13.Theorem.1}, we observe that, by the operational calculus of the preceding section, 
\[
f _{t} (H) ^{*} = f _{-t} (H) \quad \text{ and } \quad f _{t} (H) f _{s} (H) = f _{t+s} (H), f _{0} (H) = I.
\]

We also have the strong continuity of $f_{t}(H)$ at $t = 0$ by 
\[
\left\| f _{t} (H) x - x \right\| ^{2} = \int_{- \infty} ^{\infty} | \exp (i t \lambda) - 1 | ^{2} \dd \| E (\lambda) x \| ^{2} \rightarrow 0 \text{ as } t \rightarrow 0.
\]

Hence $U_{t} = f_{t}(H)$ is a one-parameter group of class $(C_0)$ of unitary operators. 
\end{proof}

\begin{remarku} For the original proof, see M. H. \textsc{Stone} \cite{stone312}, Cf. also J. \textsc{von} \textsc{Neumann} \cite{neumann235}. Another proof given by E. \textsc{Hopf} \cite{hirsch109} is based on a theorem due to S. \textsc{Bochner}: 
\end{remarku}

\begin{theorem}[\textsc{Bochner}]\label{thm:XI.13.Theorem.2} $A$ complex-valued continuous function $f(t)$, $-\infty < t < \infty$, is representable as 
\begin{equation}\label{eq:XI.13.2}
f (t) = \int_{- \infty} ^{\infty} e ^{i t \lambda} \dd v (\lambda) \text{ with a non-decreasing right-continuous bounded function } v(\lambda),
\end{equation}
right-continuous bounded function $v(\lambda)$, 


iff $f(t)$ is positive definite in the following sense: 
\begin{equation}\label{eq:XI.13.3}
\int_{- \infty} ^{\infty} \int_{- \infty} ^{\infty} f (t - s) \varphi (t) \overline {{\varphi (s)}} \dd t \dd s \geqq 0
\end{equation}
for every continuous function $\varphi$ with compact support. 
\end{theorem}


The proof of Theorem~\ref{thm:XI.13.Theorem.1} given by E. \textsc{Hopf} starts with the fact that $f(t) = (U_{t}x, x)$ satisfies condition ($3$) as may be seen from 
\[
\begin{aligned} & \int_{- \infty} ^{\infty} \int_{- \infty} ^{\infty} (U _{t-s} x, x) \varphi (t) \overline {{{{\varphi (s)}}}} \dd t \dd s = \int_{- \infty} ^{\infty} \int_{- \infty} ^{\infty} (U _{t} x, U _{s} x) \varphi (t) \overline {{{{\varphi (s)}}}} \dd t \dd s \\ & \qquad = \left(\int_{- \infty} ^{\infty} \varphi (t) U _{t} x \dd t, \int_{- \infty} ^{\infty} \varphi (s) U _{s} x \dd s\right) \geqq 0. \end{aligned}
\]

We shall show that Bochner's theorem is a consequence of Stone's theorem. 

Deduction of Theorem~\ref{thm:XI.13.Theorem.2} from Theorem~\ref{thm:XI.13.Theorem.1}. Consider the totality $\mathfrak{F}$ of complex-valued functions $x(t), -\infty < t < \infty$, such that $x(t) = 0$ except possibly for a finite set of values of $t$; the finite set may vary with $x$. $\mathfrak{F}$ is a pre-Hilbert space by 
\begin{equation}\label{eq:XI.13.4}
\begin{aligned}
&(x + y) (t) = x (t) + y (t), (\alpha x) (t) = \alpha x (t) \text{ and } (x, y) = \sum\limits_{- \infty < t, s < \infty} f (t - s) x (t) \overline {{y (s)}} \text{ for } x, y \in \mathfrak {F},
\end{aligned}
\end{equation}
excepting the axiom that $(x, x) = 0$ implies $x = 0$. That $(x, x) \geqq 0$ for any $x \in F$ is a simple consequence of the positive definiteness of the function $f(t)$. 

Let us set $\mathfrak{N} = \{x \in \mathfrak{F}; (x, x) = 0\}$. Then the factor space $\mathfrak{F}/\mathfrak{N}$ is a pre-Hilbert space with respect to the scalar product $(\bar{x}, \bar{y}) = (x, y)$ where $\bar{x}$ is the residue class mod $\mathfrak{N}$ containing $x \in \mathfrak{F}$. Let $X$ be the completion of the pre-Hilbert space $X = \mathfrak{F}/\mathfrak{N}$. The operator $U_{\tau}$ defined by 
\begin{equation}\label{eq:XI.13.5}
(U _{\tau} x) (t) = x (t - \tau), x \in \mathfrak {F},
\end{equation}
surely satisfies the conditions 
\begin{equation}\label{eq:XI.13.6}
(U _{\tau} x, U _{\tau} y) = (x, y), U _{\tau} U _{\sigma} = U _{\tau + \sigma} \text{ and } U _{0} = I.
\end{equation}

Therefore, it is easy to see that $\{U_{\tau}\}$ naturally defines a unitary operator $\hat{U}_{\tau}$ in $X$ in such a way that $\{\hat{U}_{\tau}\}, -\infty < \tau < \infty$, constitutes a one-parameter semi-group of class $(C_{0})$ of unitary operators in $X$; the strong continuity in $t$ of $\hat{U}_{t}$ follows from the continuity of the function $f(t)$. 

Hence, by Stone's theorem, $\hat{U}_t = \int_{-\infty}^{\infty} e^{it\lambda} dE(\lambda)$. Let $x_0(t) \in \mathfrak{F}$ be defined by $x_0(\tau) = 1$ and $x_0(t) = 0$ whenever $t \neq \tau$. Then, by \eqref{eq:XI.13.4} and \eqref{eq:XI.13.5}, $f(\tau) = (U_\tau x_0, x_0)$. Therefore, 
\[
f (\tau) = \int_{- \infty} ^{\infty} e ^{i \tau \lambda} \dd \| E (\lambda) \bar {x} _{0} \| ^{2},
\]
which proves Bochner's theorem. 

\begin{remarku} The idea of using a positive definite function to define a pre-Hilbert space as in \eqref{eq:XI.13.4} was systematically applied by B. Sz. \textsc{Nagy} \cite{nagy219} to various interesting problems concerning the Hilbert space. 
\end{remarku}


\yossection{14}{A Canonical Form of a Self-adjoint Operator with Simple Spectrum}
\label{sec:XI.14}

Let $H = \int \lambda dE(\lambda)$ be a self-adjoint operator in a Hilbert space $X$ with simple spectrum as defined in Chapter~\ref{chap:XI}, \ref{sec:XI.8}. Thus there exists an element $y \in X$ such that the set $\{(E(\beta) - E(\alpha))y; \alpha < \beta\}$ spans a dense linear subspace of $X$. We put 
\begin{equation}\label{eq:XI.14.1}
\sigma (\lambda) = (E (\lambda) y, y).
\end{equation}

Then $\sigma(\lambda)$ is monotone non-decreasing, right-continuous and bounded. We shall denote by $\sigma(B)$ the Baire measure determined on Baire sets on $\mathbb{R}^{1}$ from $\sigma((a,b]) = \sigma(b) - \sigma(a)$. We denote by $L_{\sigma}^{2}(-\infty,\infty)$ the totality of complex-valued Baire-measurable functions $f(\lambda), -\infty < \lambda < \infty$, such that $\|f\|_{\sigma} = \left(\int_{-\infty}^{\infty}|f(\lambda)|^{2}\sigma(\dd\lambda)\right)^{1/2} < \infty$. Then $L_{\sigma}^{2}(-\infty,\infty)$ is a Hilbert space by the scalar product $(f,g)_{\sigma} = \int_{-\infty}^{\infty}f(\lambda)\overline{g(\lambda)}\sigma(\dd\lambda)$ with the convention to consider $f = g$ in $L_{\sigma}^{2}$ iff $f(\lambda) = g(\lambda)\sigma$ -a.e. 

\begin{theoremu} With any $f(\lambda) \in L_{\sigma}^{2}(-\infty, \infty)$ we associate a vector $\hat{f}$ of $X$ defined by 
\begin{equation}\label{eq:XI.14.2}
\hat {f} = \int_{- \infty} ^{\infty} f (\lambda) \dd E (\lambda) y.
\end{equation}
\end{theoremu}

Then the correspondence $f(\lambda) \to \hat{f}$ is a one-one linear isometric mapping of $L_{\sigma}^{2}(-\infty, \infty)$ onto $X$. Let this mapping be denoted by $V$, i.e., $\hat{f} = Vf$. Then the operator $H_{1} = V^{-1}HV$ in $L_{\sigma}^{2}(-\infty, \infty)$ is precisely the 
operator of multiplication by $\lambda$: 
\begin{equation}\label{eq:XI.14.3}
\begin{aligned}
D (H _{1}) &= D (V ^{-1} H V) = \{f (\lambda); f (\lambda) \text{ and } \lambda f (\lambda) \text{ both } \in L _{\sigma} ^{2} (- \infty , \infty) \}
\\
\text{ and } (H _{1} f) (\lambda) &= \lambda f (\lambda) \text{ whenever } f (\lambda) \in D (H _{1}).
\end{aligned}
\end{equation}

\begin{proof}
 We have, by $E(\lambda)E(\mu) = E(\min (\lambda ,\mu))$, 
\[
\begin{aligned} (E (\lambda) y, \hat {f}) & = \int_{- \infty} ^{\infty} \overline {{f (\mu)}} d _{\mu} (E (\lambda) y, E (\mu) y) = \int_{- \infty} ^{\infty} \overline {{f (\mu)}} d _{\mu} (E (\mu) E (\lambda) y, y) \\ & = \int_{- \infty} ^{\lambda} \overline {{f (\mu)}} \dd (E (\mu) y, y) = \int_{- \infty} ^{\lambda} \overline {{f (\mu)}} \sigma (\dd \mu), \end{aligned}
\]
and so 
\begin{equation}\label{eq:XI.14.4}
(\hat {f}, \hat {g}) = \int_{- \infty} ^{\infty} f (\lambda) \dd (E (\lambda) y, \hat {g}) = \int_{- \infty} ^{\infty} f (\lambda) \overline {{g (\lambda)}} \sigma (\dd \lambda) = (f, g) _{\sigma}. \end{equation}

Therefore, $V$ maps $D(V) = L_{\sigma}^{2}(-\infty, \infty)$ onto $\{\hat{f}; \hat{f} = \int_{-\infty}^{\infty} f(\lambda) dE(\lambda) y, f \in L_{\sigma}^{2}(-\infty, \infty)\}$ one-to-one, linearly and isometrically. Hence, in particular, $R(V)$ is a closed linear subspace of $X$. But $R(V)$ surely contains the elements of the form $\int_{\alpha}^{\beta} dE(\lambda) y = (E(\beta) - E(\alpha)) y, -\infty < \alpha < \beta < \infty$, and so, by the hypothesis that the spectrum of $H$ is simple, we see that $R(V) = R(V)^a = X$. Thus the first half of the Theorem is proved. 

Next we have 
\[
\begin{aligned}
&E (\lambda) \hat {f}  = E (\lambda) \int_{- \infty} ^{\infty} f (\mu) \dd E (\mu) y = \int_{- \infty} ^{\infty} f (\mu) d _{\mu} (E (\lambda) E (\mu) y) = \int_{- \infty} ^{\lambda} f (\mu) \dd E (\mu) y,
\end{aligned}
\]
and so, by \eqref{eq:XI.14.4}, 
\begin{equation}\label{eq:XI.14.5}
(E (\lambda) \hat {f}, \hat {g}) = \int_{- \infty} ^{\lambda} f (\mu) \overline {{g (\mu)}} \sigma (\dd \mu).
\end{equation}

Hence the condition $\int_{-\infty}^{\infty}\lambda^{2}\dd(E(\lambda)\hat{f},\hat{f})<\infty$, equivalent to $\hat{f}\in D(H)$, is equivalent to the condition $\int_{-\infty}^{\infty}\lambda^{2}|f(\lambda)|^{2}\sigma(\dd\lambda)<\infty$. Moreover, in the last case, we have, by \eqref{eq:XI.12.20} in Chapter~\ref{chap:XI}, \ref{sec:XI.12}, 
\[
H V f = H \hat {f} = \int_{- \infty} ^{\infty} \lambda \dd E (\lambda) \hat {f},
\]
and hence, by \eqref{eq:XI.14.4} and \eqref{eq:XI.14.5}, 
\[
\begin{aligned} & {(H _{1} f, g) _{\sigma} = (V ^{-1} H V f, g) _{\sigma} = (H V \cdot f, V \cdot g) = (H \cdot \hat {f}, \hat {g})} \\ & {\qquad = \int_{- \infty} ^{\infty} \lambda \dd (E (\lambda) \hat {f}, \hat {g}) = \int_{- \infty} ^{\infty} \lambda f (\lambda) \overline {{g (\lambda)}} \sigma (\dd \lambda).} \end{aligned}
\]

On the other hand, we have 
\[
(H _{1} f, g) _{\sigma} = \int_{- \infty} ^{\infty} (H _{1} f) (\lambda) g (\lambda) \sigma (\dd \lambda).
\]

Therefore we must have 
\[
(H _{1} f) (\lambda) = \lambda f (\lambda) \sigma \text{-almost everywhere. }
\]
\end{proof}

\begin{remarku} There is a close connection between self-adjoint operators with simple spectrum and the Jacobi matrices. See M. H. \textsc{Stone} \cite{stone311}, p. $275$. For a canonical form of a self-adjoint operator with not necessarily simple spectrum, see J. von Neumann's reduction theory: \textsc{Neumann} [$9$]. 
\end{remarku}


\yossection{15}{The Defect Indices of a Symmetric Operator. The Generalized Resolution of the Identity}
\label{sec:XI.15}

\begin{definition}\label{def:XI.15.Definition.1} Let $U = U_H = (H - iI)(H + iI)^{-1}$ be the Cayley transform of a closed symmetric operator $H$ in a Hilbert space $X$. Let $X_H^+ = D(U_H)^\perp$, $X_H^- = R(U_H)^\perp$, and let $m = \dim(X_H^+)$, $n = \dim(X_H^-)$ be the dimension numbers of $X_H^+$, $X_H^-$ respectively. Then $H$ is said to be of the defect indices $(m, n)$. $H$ is self-adjoint iff it is of the defect indices $(0, 0)$ (see Chapter~\ref{chap:VII}, \ref{sec:VII.4}). 
\end{definition}

\begin{proposition}\label{prop:XI.15.Proposition.1} The defect indices $(m, n)$ of a closed symmetric operator $H$ may be defined as follows: $m$ is the dimension number of the linear subspace $\{x \in X; H^{*}x = ix\}$; $n$ is the dimension number of the linear subspace $\{x \in X; H^{*}x = -ix\}$. 
\end{proposition}

\begin{proof}
 Clear from Theorem~\ref{thm:VII.4.Theorem.3} in Chapter~\ref{chap:VII}, \ref{sec:VII.4}. 
\end{proof}

\begin{example}\label{ex:XI.15.Example.1} Let $X = L^2(0,1)$. Let $D$ be the totality of absolutely continuous functions $x(t) \in L^2(0,1)$ such that $x(0) = x(1) = 0$ and $x'(t) \in L^2(0,1)$. Then the operator $T_1$ defined by $T_1x(t) = i^{-1}x'(t)$ on $D = D(T_1)$ is of the defect indices ($1$, $1$). 
\end{example}

\begin{proof}
 As was shown in Example~\hyperref[ex:VII.3.Example.4]{$4$} of Chapter~\ref{chap:VII}, \ref{sec:VII.3}, $T_1^* = T_2$ is defined by 
\[
T _{2} x (t) = i ^{-1} x ^{\prime} (t) \text{ on } D (T _{2}) = \{x (t) \in L ^{2} (0, 1); x (t) \text{ is }
\]
absolutely continuous such that $x'(t) \in L^2(0, 1)$. 

Thus the solution $y \in L^2(0,1)$ of $T_1^*y = T_2y = iy$ is a distribution solution of the differential equation 
\begin{equation}\label{eq:XI.15.1}
y ^{\prime} (t) = - y (t) \left(y, y ^{\prime} \in L ^{2} (0, 1)\right).
\end{equation}

Then $z(t) = y(t) \exp(t)$ is a distribution solution of the differential equation 
\begin{equation}\label{eq:XI.15.2}
z ^{\prime} (t) = 0 \left(z, z ^{\prime} \in L ^{2} (0, 1)\right).
\end{equation}

We shall show that there exists a constant $C$ such that $z(t) = C$ for a.e. $t \in (0,1)$. To this purpose, take any function $x_0(t) \in C_0^\infty(0,1)$ such that 
\[
\begin{aligned}
&\int_{0} ^{1} x _{0} (t) \dd t  = 1 \text{ and put } x (t) - x _{0} (t) \int_{0} ^{1} x (t) \dd t  = u (t), w (t) = \int_{0} ^{t} u (s) \dd s,
\end{aligned}
\]
where $x(t)$ is an arbitrary function from $C_0^\infty(0,1)$. We have $w \in C_0^\infty(0,1)$ by $\int_{0}^{1} u(s) \dd s = 0$. Therefore, by \eqref{eq:XI.15.2}, we have 
\[
- \int_{0} ^{1} z (t) w ^{\prime} (t) \dd t = - \int_{0} ^{1} z (t) u (t) \dd t = 0,
\]
that is, 
\[
\int_{0} ^{1} z (t) x (t) \dd t = C \int_{0} ^{1} x (t) \dd t, \text{ where } C = \int_{0} ^{1} z (t) x _{0} (t) \dd t.
\]

This proves, by the arbitrariness of $x(t) \in C_0^\infty(0,1)$, that $z(t) = C$ for a.e. $t \in (0,1)$. 

Hence any solution of $T^{*}y = iy$ is of the form $y(t) = C\exp (-t)$. In the same manner we see that any solution of $T^{*}y = -iy$ is of the form $y(t) = C\exp (t)$. Thus $T$ is of the defect indices ($1$, $1$). 
\end{proof}

\begin{definition}\label{def:XI.15.Definition.2} $A$ symmetric operator $H$ in a Hilbert space $X$ is called maximal symmetric if there is no proper symmetric extension of $H$. 
\end{definition}

\begin{proposition}\label{prop:XI.15.Proposition.2} $A$ maximal symmetric operator $H$ is closed and $H = H^{**}$. A self-adjoint operator $H$ is maximal symmetric. 
\end{proposition}

\begin{proof}
 By Proposition~\ref{prop:VII.3.Proposition.1} in Chapter~\ref{chap:VII}, \ref{sec:VII.3}, $H^{**}$ is a closed symmetric extension of $H$. Thus the first half of Proposition~\ref{prop:XI.15.Proposition.2} is clear. Let $H_0$ be a symmetric extension of a self-adjoint operator $H$. Then, from $H \subseteq H_0$, $H_0 \subseteq H_0^*$, we obtain $H_0 \subseteq H_0^* \subseteq H^*$ and so, by $H = H^*$, $H \subseteq H_0 \subseteq H$. This proves that a self-adjoint operator $H$ is maximal symmetric. 
\end{proof}

\begin{corollary}\label{cor:XI.15.Corollary.1} Every maximal symmetric extension $H_0$ of a given symmetric operator $H$ is also an extension of $H^{**}$. 
\end{corollary}

\begin{proof}
 The relation $H \subseteq H_0$ implies the relation $H_0^* \subseteq H^*$ and $H^{**} \subseteq H_0^{**}$. Since, by the preceding proposition, $H_0 = H_0^{**}$, Corollary~\ref{cor:XI.15.Corollary.1} is true. 
\end{proof}

\begin{corollary}\label{cor:XI.15.Corollary.2} If $H$ is a symmetric operator such that $H^{*} = H^{**}$, then the self-adjoint operator $H^{*}$ is the only maximal symmetric extension of $H$. 
\end{corollary}

\begin{proof}
 Being self-adjoint, $H^{**}$ is maximal symmetric. Thus any maximal symmetric extension $H_{0}$ of $H$, which is also a symmetric extension of $H^{**}$ by Corollary~\ref{cor:XI.15.Corollary.1}, is identical to $H^{**} = H^{*}$. 
\end{proof}

We are thus in a position to state 

\begin{definition}\label{def:XI.15.Definition.3} $A$ symmetric operator $H$ such that $H^{*} = H^{**}$ is said to be essentially self-adjoint. A self-adjoint operator $H$ is said to be hypermaximal. The latter terminology is due to J. \textsc{von} \textsc{Neumann}. 
\end{definition}

\begin{example}\label{ex:XI.15.Example.2} Let $X = L^2(-\infty, \infty)$, and define an operator $H$ by $Hx(t) = t \cdot x(t)$ for $x(t) \in C_0^0(-\infty, \infty)$. $H$ is surely a symmetric operator in $X$. It is easy to see that $H^*$ is the coordinate operator defined in Example~\hyperref[ex:VII.3.Example.2]{$2$} in Chapter~\ref{chap:VII}, \ref{sec:VII.3} so that $H$ is essentially self-adjoint. The operator $H$ defined by $Hx(t) = i^{-1}x'(t)$ for $x(t) \in C_0^1(-\infty, \infty)$ is also essentially self-adjoint in $X = L^2(-\infty, \infty)$. For, in this case, $H^*$ is the momentum operator defined in Example~\hyperref[ex:VII.3.Example.3]{$3$} in Chapter~\ref{chap:VII}, \ref{sec:VII.3}. 
\end{example}

\begin{theorem}\label{thm:XI.15.Theorem.1} Let the defect indices $(m, n)$ of a closed symmetric operator $H$ satisfy 
\[
m = m ^{\prime} + p, n = n ^{\prime} + p \quad (p > 0).
\]
\end{theorem}

Then there exists a closed symmetric extension $H'$ of $H$ with the defect indices $(m', n')$. 

\begin{proof}
 Let $\{\varphi_1, \varphi_2, \ldots, \varphi_p, \varphi_{p+1}, \ldots, \varphi_{p+m}\}, \{\psi_1, \psi_2, \ldots, \psi_p, \psi_{p+1}, \ldots, \psi_{p+n'}\}$ be complete orthonormal systems of $X_H^+ = D(U_H)^\perp$, $X_H^- = R(U_H)^\perp$, respectively. Define an isometric extension $V$ of $U_H$ by 
\begin{align}
V x& = U _{H} x \quad \text{ for } \quad x \in D (U _{H}), \notag \\
V \cdot \sum\limits_{i=1} ^{p} \alpha_{i} \varphi_{i}& = \sum\limits_{i=1} ^{p} \alpha_{i} \psi_{i}. \notag
\end{align}

We have $R(I - U_{H})^{a} = X$ by Theorem~\ref{thm:VII.4.Theorem.1} in Chapter~\ref{chap:VII}, \ref{sec:VII.4}. Thus, by $R(I - V) \supseteq R(I - U_{H})$ and Theorem~\ref{thm:VII.4.Theorem.2} in Chapter~\ref{chap:VII}, \ref{sec:VII.4}, there exists a uniquely determined closed symmetric extension $H'$ of $H$ such that $V = (H' - iI)(H' + iI)^{-1}$. The defect indices of $H'$ are $(m', n')$ as may be seen from $\dim(D(V)^{\perp}) = m'$, $\dim(R(V)^{\perp}) = n'$. 
\end{proof}

\begin{corollaryu} $A$ closed symmetric operator $H$ with the defect indices $(m, n)$ is maximal symmetric iff either $m = 0$ or $n = 0$. 
\end{corollaryu}

\begin{proof}
 The ``only if'' part is clear from Theorem~\ref{thm:XI.15.Theorem.1}. If, for instance, $m = 0$, then, by $D(U_{H}) = X$, we must have $U_{H_{0}} = U_{H}$ for a closed symmetric extension $H_{0}$ of $H$. This proves $H_{0} = i(I + U_{H_{0}})(I - U_{H_{0}})^{-1} = i(I + U_{H})(I - U_{H})^{-1} = H$. Also, in the case $n = 0$, we see that there is no proper closed symmetric extension of $H$. 
\end{proof}

\begin{example}\label{ex:XI.15.Example.3} Suppose that $\{\varphi_1, \varphi_2, \ldots, \varphi_n, \ldots\}$ is a complete ortho-normal system of a (separable) Hilbert space $X$. Then, by 
\[
U \cdot \sum\limits_{i=1} ^{\infty} \alpha_{i} \varphi_{i} = \sum\limits_{i=1} ^{\infty} \alpha_{i} \varphi_{i+1}, \sum\limits_{i=1} ^{\infty} | \alpha_{i} | ^{2} < \infty
\]
we define a closed isometric operator $U$ such that $D(U) = X$ and $\dim (R(U)^\perp) = 1$. If $R(I - U)^a \neq X$, then there exists an $x \neq 0$ such that $x \in R(I - U)^\perp$. Consequently $((I - U)x, x) = 0$ and so $(Ux, x) = \|x\|^2 = \|Ux\|^2$. This implies that 
\begin{equation}\left\| (I - U) x \right\| ^{2} = \left\| x \right\| ^{2} - (U x, x) - (x, U x) + \left\| U x \right\| ^{2} = \| x \| ^{2} - \| x \| ^{2} - \| x \| ^{2} + \| x \| ^{2} = 0, \notag\end{equation}
that is, $Ux = x$ and so, by the above definition of $U$, $x$ must be zero. This contradiction shows that $R(I - U)^a = X$. Hence, by Theorem~\ref{thm:VII.4.Theorem.2} in Chapter~\ref{chap:VII}, \ref{sec:VII.4}, $U$ is the Cayley transform of a closed symmetric operator $H$. $H$ is of the defect indices ($0$, $1$) since $D(U) = X$ and $R(U)^\perp$ is spanned by $\varphi_1$. Thus $H$ is maximal symmetric without being self-adjoint. 
\end{example}

\begin{theorem}[M. A. \textsc{Naimark} {\cite{naimark223}}]\label{thm:XI.15.Theorem.2} Let a closed symmetric operator $H_{1}$ in a Hilbert space $X_{1}$ be of the defect indices $(m, n)$. Then we can construct a Hilbert space $X$, containing $X_{1}$ as a closed linear subspace, and a closed symmetric operator $H$ in $X$ with the defect indices $(m + n, m + n)$ such that 

$H_{1} = P(X_{1})HP(X_{1})$, where $P(X_{1})$ is the projection of $X$ onto $X_{1}$ 
\end{theorem}

\begin{proof}
 Consider a Hilbert space $X_{2}$ of the same dimensionality as $X_{1}$. We construct a closed symmetric operator $H_{2}$ in $X_{2}$ with the defect indices $(n, m)$. For instance, we may take $H_{2} = -H_{1}$ supposing that $X_{2}$ coincides with $X_{1}$. Then we would have $\{x \in X_{2}; H_{2}^{*}x = ix\} = \{x \in X_{1}; H_{1}^{*}x = -ix\}, \{x \in X_{2}; H_{2}^{*}x = -ix\} = \{x \in X_{1}; H_{1}^{*}x = ix\}$ and so the defect indices of $H_{2}$ must be $(n, m)$. 

We then consider an operator $H$ defined by 
\[
H \{x, y \} = \left\{H _{1} x, H _{2} y \right\} \text{ for } \{x, y \} \in X _{1} \times X _{2} \text{ with } x \in D (H _{1}), y \in D (H _{2}).
\]

It is easy to see that $H$ is a closed symmetric operator in the product Hilbert space $X = X_{1} \times X_{2}$. The condition 
\[
H ^{*} \{x, y \} = i \{x, y \} \quad (\text{ or the condition } H ^{*} \{x, y \} = - i \{x, y \})
\]
means that $H_1^* x = ix, H_2^* y = iy$ (or $H_1^* x = -ix, H_2^* y = -iy$). Hence we see that the defect indices of $H$ are $(m + n, m + n)$. 
\end{proof}

\begin{corollaryu} Let, by Theorem~\ref{thm:XI.15.Theorem.1}, $\hat{H}$ be a self-adjoint extension of $H$, and let $\hat{H} = \int \lambda \dd\hat{E} (\lambda)$ be the spectral resolution of $\hat{H}$. Since $H$ and a fortiori $\hat{H}$ are extensions of $H_{1}$ when $H_{1}$ is considered as an operator in $X$, we have the result: 
\begin{gather}\text{ if } x \in D (H _{1}) \subseteq X _{1} = P (X _{1}) X, \text{ then } x = P (X _{1}) x \in D (\hat {H}) \notag \\H _{1} x = P (X _{1}) \hat {H} x = P (X _{1}) \hat {H} P (X _{1}) x = \int_{- \infty} ^{\infty} \lambda \dd F (\lambda) x, \text{ where } F (\lambda) = P (X _{1}) \hat {E} (\lambda) P (X _{1}). \label{eq:XI.15.3}
\end{gather}
\end{corollaryu}

The system $\{F(\lambda); -\infty < \lambda < \infty\}$ surely satisfies the conditions: 

$F(\lambda)$ is a self-adjoint operator in $X_{1}$, 
\begin{gather}\lambda_{1} < \lambda_{2} \text{ implies that } (F (\lambda_{1}) x, x) \leqq (F (\lambda_{2}) x, x) \text{ for every } x \in X _{1}, F (\lambda + 0) = F (\lambda), \label{eq:XI.15.4} \\F (- \infty) x = \slim _{\lambda \downarrow - \infty} F (\lambda) x = 0, F (\infty) x = \slim _{\lambda \uparrow \infty} F (\lambda) x = x \text{ for all } x \in X. \notag
\end{gather}

\begin{remarkn}[Remark $1$] For the closed symmetric operator $H_{1}$, we have 
\[
H _{1} x = \int_{- \infty} ^{\infty} \lambda \dd F (\lambda) x \text{ for all } x \in D (H _{1}),
\]
\end{remarkn}

where $\{F(\lambda);-\infty<\lambda<\infty\}$ satisfies \eqref{eq:XI.15.4}. In this sense, $H_{1}$ admits a generalized spectral resolution. 

\begin{example}\label{ex:XI.15.Example.4} Let $X_{1} = L^{2}(-\infty,0)$, and let $D_{1}$ be the totality of absolutely continuous functions $x(t) \in L^{2}(-\infty,0)$ such that $x(0) = 0$ and $x'(t) \in L^{2}(-\infty,0)$. Then the operator $H_{1}$ defined by $H_{1}x(t) = i^{-1}x'(t)$ on $D_{1} = D(H_{1})$ is a maximal symmetric operator with the defect indices ($0$, $1$). This we see as in the above Example~\hyperref[ex:XI.15.Example.1]{$1$}. Let $X_{2} = L^{2}(0,\infty)$, and let $D_{2}$ be the totality of absolutely continuous functions $x(t) \in L^{2}(0,\infty)$ such that $x(0) = 0$ and $x'(t) \in L^{2}(0,\infty)$. Then the operator $H_{2}$ defined by $H_{2}x(t) = i^{-1}x'(t)$ on $D_{2} = D(H_{2})$ is a maximal symmetric operator with the defect indices ($1$, $0$). In this case, $X = X_{1} \times X_{2} = L^{2}(-\infty,\infty)$, and the operator $H$ in Theorem~\ref{thm:XI.15.Theorem.2} is given concretely by $Hx(t) = i^{-1}x'(t)$ for those $x(t) \in L^{2}(-\infty,\infty)$ for which $x(0) = 0$ and $x'(t) \in L^{2}(-\infty,\infty)$. 
\end{example}

\begin{remarkn}[Remark $2$] Since $\hat{H} = \int \lambda \dd\hat{E} (\lambda)$ is an extension of $H_{1}$, we have, for $x \in D(H_{1})$, 
\begin{equation}\label{eq:XI.15.5}
\begin{aligned} \| H _{1} x \| ^{2} & = \| \hat {H} x \| ^{2} = \int_{- \infty} ^{\infty} \lambda^{2} \dd \| \hat {E} (\lambda) x \| ^{2} = \int_{- \infty} ^{\infty} \lambda^{2} \dd (\hat {E} (\lambda) x, x) \\ & = \int_{- \infty} ^{\infty} \lambda^{2} \dd (\hat {E} (\lambda) P (X _{1}) x, P (X _{1}) x) = \int_{- \infty} ^{\infty} \lambda^{2} \dd (F (\lambda) x, x). \end{aligned}
\end{equation}
\end{remarkn}

However, the condition $\int_{-\infty}^{\infty} \lambda^2 \dd(F(\lambda)x, x) < \infty$ does not necessarily imply that $x \in D(H_1)$. Concerning this point we have 

\begin{theorem}\label{thm:XI.15.Theorem.3} In the case of a maximal symmetric operator $H_{1}$, we have, for the corresponding generalized spectral resolution $\int \lambda dF(\lambda)$ in $X_{1}$, 

$x \in D(H_1)$ is equivalent to the condition $\int_{-\infty}^{\infty} \lambda^2 \dd(F(\lambda)x, x) < \infty$. 
\end{theorem}


\begin{proof}
 The reasoning in Remark $2$ shows that $\int_{-\infty}^{\infty} \lambda^2 \dd(F(\lambda)x, x) < \infty$ implies that $\int_{-\infty}^{\infty} \lambda^2 \dd\|\hat{E}(\lambda)x\|^2 < \infty$, i.e., $x \in D(\hat{H})$. The operator 
\[
H ^{\prime} = P (X _{1}) \hat {H} P (X _{1}),
\]
when considered as an operator in $X_{1}$, is a symmetric extension of $H_{1}$ and $D(H') = D(\hat{H}) \cap X_{1}$. Thus, by the maximality of $H_{1}$, we must have $H_{1} = H'$. Thus $D(H_{1}) = D(H') = D(\hat{H}) \cap X_{1}$ and so, by $H' = P(X_{1})\hat{H}P(X_{1})$, the condition $\int_{-\infty}^{\infty}\lambda^{2}\dd\|\hat{E}(\lambda)x\|^{2} < \infty$ with $x \in X_{1}$ implies 
$\int_{-\infty}^{\infty}\lambda^{2}\dd(F(\lambda)x,x) < \infty$; and conversely, the condition $x\in D(H_1)$ implies $\int_{-\infty}^{\infty}\lambda^{2}\dd(F(\lambda)x,x) = \int_{-\infty}^{\infty}\lambda^{2}\dd\|\hat{E} (\lambda)x\|^{2} < \infty .$ 
\end{proof}

\begin{remarkn}[Remark $3$] Since, by Theorem~\ref{thm:XI.15.Theorem.1}, any closed symmetric operator can be extended to a maximal symmetric operator $H_{1}$, we can apply our Theorem~\ref{thm:XI.15.Theorem.3} to the effect that ($5$) is valid. For a detailed exposition of the generalized spectral resolution, see N. I. \textsc{Achieser}-I. M. \textsc{Glazman} \cite{achieser1} or B. Sz. \textsc{Nagy} \cite{nagy219}. The spectral representation of a self-adjoint operator in a Hilbert space can be extended to a certain class of linear operators in a Banach space with an appropriate modification. This result is due to N. \textsc{Dunford} and may be considered as an ``elementary divisor theory'' in infinite dimensional spaces. N. \textsc{Dunford}-J. \textsc{Schwartz} \cite{dunford49}. 
\end{remarkn}


\yossection{16}{The Group-ring \texorpdfstring{$L^1$}{L1} and Wiener's Tauberian Theorem}
\label{sec:XI.16}

The Gelfand representation admits another important application in functional analysis, namely, an operator-theoretical treatment of the Tauberian theorem of N. \textsc{Wiener}. 

The linear space $L^{1}(-\infty,\infty)$ is a ring with respect to the function sum and the product $\times$ defined by 
\begin{equation}\label{eq:XI.16.1}
(f \times g) (t) = (f * g) (t) = \int_{- \infty} ^{\infty} f (t - s) g (s) \dd s.
\end{equation}

For, by the Fubini-Tonelli theorem, 
\[
\begin{aligned}
&\int_{- \infty} ^{\infty} \left| \int_{- \infty} ^{\infty} f (t - s) g (s) \dd s \right| \dd t  \leqq \int_{- \infty} ^{\infty} | f (t - s) | \dd t \int_{- \infty} ^{\infty} | g (s) | \dd s = \int_{- \infty} ^{\infty} | f (t) | \dd t \int_{- \infty} ^{\infty} | g (t) | \dd t.
\end{aligned}
\]

We have thus proved 
\begin{equation}\label{eq:XI.16.2}
\left\| f \times g \right\| \leqq \left\| f \right\| \cdot \left\| g \right\|, \text{ where } \| \| \text{ is the norm in } L ^{1} (- \infty , \infty).
\end{equation}

Therefore we have proved 

\begin{proposition}\label{prop:XI.16.Proposition.1} We can introduce formally a multiplicative unit $e$ in the ring $L^{1}$ consisting of all elements $\tilde{z}$ given formally by 
\begin{equation}\label{eq:XI.16.3}
\tilde {z} = \lambda e + x, x \in L ^{1} (- \infty , \infty).
\end{equation}
\end{proposition}

In fact $L^1$ is a normed ring by the following rule: 
\begin{equation}\label{eq:XI.16.4}
\begin{gathered}
(\lambda_{1} e + x _{1}) + (\lambda_{2} e + x _{2}) = (\lambda_{1} + \lambda_{2}) e + (x _{1} + x _{2}), \alpha (\lambda e + x) = \alpha \lambda e + \alpha x,
\\
(\lambda_{1} e + x _{1}) \times (\lambda_{2} e + x _{2}) = \lambda_{1} \lambda_{2} e + \lambda_{1} x _{2} + \lambda_{2} x _{1} + x _{1} \times x _{2}, \| \lambda e + x \| = | \lambda | + \| x \|.
\end{gathered}
\end{equation}

This normed ring $L^{1}$ is called the group-ring of the group $\mathbb{R}^{1}$ of real numbers written additively. We shall find all maximal ideals of this ring $L^{1}$. One is $I_{0}=L^{1}(-\infty,\infty)$ from which $L^{1}$ is obtained by adjunction of the unit $e$. We shall find all maximal ideals $I\neq I_{0}$. 

For any maximal ideal $I$ of the normed ring $L^1$, we shall denote by $(\tilde{z}, I)$ the complex number which corresponds to the element $\tilde{z}$ by the ring homomorphic mapping $L^1 \to L^1 / I$. Thus $(\tilde{z}, I_0) = \lambda$ for $\tilde{z} = \lambda e + x, x \in I_0$. 

Let $I$ be a maximal ideal of $L^1$ different from $I_0$. Then there exists a function $x \in L^1(-\infty, \infty) = I_0$ such that $(x, I) \neq 0$. We set 
\begin{equation}\label{eq:XI.16.5}
\chi (\alpha) = (x _{\alpha}, I) / (x, I), \text{ where } x _{\alpha} (t) = x (t + \alpha).
\end{equation}

Then $\chi(0) = 1$, and, by $|(x_{\alpha}, I)| \leqq \|x_{\alpha}\| = \|x\|$, we see that $|\chi(\alpha)| \leqq \|x\| / |(x, I)|$. Thus the function $\chi(\alpha)$ is bounded in $\alpha$. Moreover, by 
\[
\left| \chi (\alpha + \delta) - \chi (\alpha) \right| \leqq \left\| x _{\alpha + \delta} - x _{\alpha} \right\| / |(x, I)|,
\]
the function $\chi (\alpha)$ is continuous in $\alpha$. For, by \eqref{eq:X.1.2} in Chapter~\ref{chap:X}, \ref{sec:X.1}, we have 
\[
\lim\limits _{\delta \rightarrow 0} \int_{- \infty} ^{\infty} | x (\alpha + \delta + t) - x (\alpha + t) | \dd t = 0.
\]

On the other hand, we have, by $(x_{\alpha + \beta} \times x)(t) = (x_{\alpha} \times x_{\beta})(t)$, 
\[
\left(x _{\alpha + \beta}, I\right) (x, I) = \left(x _{\alpha}, I\right) \left(x _{\beta}, I\right)
\]
so that 
\begin{equation}\label{eq:XI.16.6}
\chi (\alpha + \beta) = \chi (\alpha) \chi (\beta).
\end{equation}

From this we can prove that there exists a uniquely determined real number $\xi = \xi(I)$ such that 
\begin{equation}\label{eq:XI.16.7}
\chi (\alpha) = \exp (i \cdot \xi (I) \cdot \alpha).
\end{equation}

In fact, by $\chi(n\alpha) = \chi(\alpha)^n$ and the boundedness of the function $\chi$, we obtain $|\chi(\alpha)| \leqq 1$. Consequently, by $\chi(\alpha) \chi(-\alpha) = \chi(0) = 1$, we must have $|\chi(\alpha)| = 1$. Thus, as a continuous solution of \eqref{eq:XI.16.6} whose absolute value is one, the function $\chi(\alpha)$ must be given in form \eqref{eq:XI.16.7}. That the value $\xi(I)$ is determined by $I$ independently of the choice of $x \in L^1(-\infty, \infty)$ may be seen from $\chi(\alpha)(y, I) = (y_\alpha, I)$ which is implied by $x_\alpha \times y = x \times y_\alpha$. 

Every continuous solution of \eqref{eq:XI.16.6} with the absolute value identically equal to one is called a continuous unitary character of the group $\mathbb{R}^{1}$ of real numbers written additively. 

We have thus constructed the (continuous unitary) character $\chi(\alpha)$ of the group $\mathbb{R}^{1}$ with respect to the given maximal ideal $I \neq I_{0}$. 

We next show how to reconstruct the ideal $I$ with respect to this character or, what amounts to the same thing, to reconstruct the value $(\tilde{z}, I)$ from $\chi$. 

For any $y(t) \in L^{1}(-\infty, \infty)$, we have 
\[
(x \times y) (t) = \int_{- \infty} ^{\infty} x (t - s) y (s) \dd s = \int_{- \infty} ^{\infty} x _{-s} (t) y (s) \dd s.
\]

Hence, by \eqref{eq:XI.16.5} and the continuity of the ring homomorphism $L^1 \to L^1 / I$, we obtain 
\[
(x \times y, I) = (x, I) (y, I) = (x, I) \int_{- \infty} ^{\infty} \chi (- s) y (s) \dd s.
\]

Thus, by $(x, I) \neq 0$, we obtain 
\begin{equation}\label{eq:XI.16.8}
(y, I) = \int_{- \infty} ^{\infty} y (s) \exp (- i \cdot \xi (I) \cdot s) \dd s.
\end{equation}

Therefore, for any $\tilde{z} = \lambda e + x$ with $x \in L^{1}(-\infty, \infty)$, 
\begin{equation}\label{eq:XI.16.9}
(\tilde {z}, I) = (\lambda e, I) + (x, I) = \lambda + \int_{- \infty} ^{\infty} x (s) \exp (- i \cdot \xi (I) \cdot s) \dd s.
\end{equation}

Conversely, any (continuous, unitary) character $\chi(\alpha) = \exp(i \cdot \xi \cdot \alpha)$ of the group $\mathbb{R}^{1}$ defines a ring homomorphism 
\begin{equation}\label{eq:XI.16.10}
\tilde {z} \rightarrow \lambda + \int_{- \infty} ^{\infty} x (s) \exp (- i \cdot \xi \cdot s) \dd s (\tilde {z} = \lambda e + x, x \in L ^{1} (- \infty , \infty))
\end{equation}
of $L^{1}$ onto the ring of complex numbers. For, we have, by the Fubini-Tonelli theorem. 
\begin{equation}\label{eq:XI.16.11}
\begin{aligned}
&\int_{- \infty} ^{\infty} (x _{1} \times x _{2}) (t) \exp (- i \cdot \xi \cdot t) \dd t = \int_{- \infty} ^{\infty} x _{1} (t) \exp (- i \cdot \xi \cdot t) \dd t \int_{- \infty} ^{\infty} x _{2} (t) \exp (- i \cdot \xi \cdot t) \dd t.
\end{aligned}
\end{equation}

We have thus proved 

\begin{theorem}[\textsc{Gelfand} {\cite{gelfand78}} and \textsc{Raikov} {\cite{raikov267}}]\label{thm:XI.16.Theorem.1} There exists a one-to-one correspondence between the set of all maximal ideals $I \neq L^{1}(-\infty, \infty)$ of the group ring $L^{1}$ of the group $\mathbb{R}^{1}$ and the set of all continuous, unitary characters $\chi(\alpha)$ of this group $\mathbb{R}^{1}$. This correspondence is defined by formula \eqref{eq:XI.16.9}. 
\end{theorem}

We shall show that the normed ring $L^{1}$ is semi-simple or, what amounts to the same thing (see Chapter~\ref{chap:XI}, \ref{sec:XI.2}), that the following theorem is true. 

\begin{theorem}\label{thm:XI.16.Theorem.2} The normed ring $L^{1}$ has no generalized nilpotent elements other than $0$. 
\end{theorem}

\begin{proof}
 Let $x \in L^{1}(-\infty, \infty)$ and $y \in L^{2}(-\infty, \infty)$. Then, by Schwarz' inequality, 
\[
\left| \int_{- \infty} ^{\infty} x (t - s) y (s) \dd s \right| \leqq \left(\int_{- \infty} ^{\infty} | x (t - s) | \dd s \int_{- \infty} ^{\infty} | x (t - s) | | y (s) | ^{2} \dd s\right) ^{1/2},
\]
and so, by the Fubini-Tonelli theorem, we see that the left hand side belongs to $L^2(-\infty, \infty)$ and 
$\|x \times y\|_{2} \leqq \|x\| \|y\|_{2}, \text{ where } \|\|_{2} \text{ is the norm in } L^{2_1}(-\infty, \infty).$ 


Hence we can define a bounded linear operator $T_{x}$ on $L^{2}(-\infty, \infty)$ into $L^{2}(-\infty, \infty)$ by 
\begin{equation}\label{eq:XI.16.12}
(T _{x} y) (t) = \int_{- \infty} ^{\infty} x (t - s) y (s) \dd s \quad \text{ whenever } \quad x \in L ^{1} (- \infty , \infty);
\end{equation}
and, moreover, we have 
\begin{align}
\left\| T _{x} \right\| _{2}& \leqq \left\| x \right\|, \label{eq:XI.16.13} \\
T _{x} ^{*}& = T _ {x ^{*}} \text{ where } x ^{*} (t) = \overline {{x (- t)}}. \label{eq:XI.16.14}
\end{align}

Thus, by applying the Fubini-Tonelli theorem again, we obtain 
\begin{equation}\label{eq:XI.16.15}
T_{x}T_{x}^{*} = T_{x\times x^{*}} = T_{x^{*}\times x} = T_{x}^{*}T_{x}, \text{ i.e., } T_{x} \text{ is a normal operator.}
\end{equation}

Hence, by Chapter~\ref{chap:XI}, \ref{sec:XI.3}, we have $\| T_x \|_2 = \lim\limits_{n \to \infty} (\|T_x^n \|_2)^{1/n}$. Therefore, by 
\[
\left| \underbrace {\left| x \times x \times x \times \cdots \times x \right|} _ {n \text{ factor }} \right| \geqq \left\| T _{x} ^{n} \right\| _{2},
\]
we see that, if $x$ is a generalized nilpotent element, then $\| T_x \|_2 = 0$. From this fact we easily prove that $x$ must be a zero vector of $L^1(-\infty, \infty)$. 

Let now $\tilde{z} = \lambda e + x$ with $x \in L^{1}(-\infty, \infty)$ be a generalized nilpotent element of the normed ring $L^{1}$. Then, for any maximal ideal $I$ of $L^{1}$, we must have $(\tilde{z}, I) = \lambda + (x, I) = 0$. Hence the Fourier transform 
\[
(2 \pi) ^{-1/2} \int_{- \infty} ^{\infty} x (t) \exp (- i \cdot \xi \cdot t) \dd t
\]
must be identically equal to $-(2\pi)^{-1/2}\lambda$. From this we prove that $\lambda$ must be $0$. For, we have (the Riemann-Lebesgue theorem) 
\begin{align}
\left| \int_{- \infty} ^{\infty} x (t) \exp (- i \cdot \xi \cdot t) \dd t \right|& = \left| 2 ^{-1} \int_{- \infty} ^{\infty} \left[ x (t) - x \left(t + \frac {\pi}{\xi}\right) \right] \exp (- i \cdot \xi \cdot t) \dd t \right| \notag \\
&\leqq 2 ^{-1} \int_{- \infty} ^{\infty} \left| x (t) - x \left(t + \frac {\pi}{\xi}\right)\right| \dd t \rightarrow 0 \text{ as } \xi \rightarrow \infty . \notag
\end{align}

Thus a generalized nilpotent element $\tilde{z} \in L^1$ must be of the form $\tilde{z} = x$ with $x \in L^1(-\infty, \infty)$ and so, by what we have proved above, $\tilde{z} = x = 0$. 
\end{proof}

We are now in a position to state and prove the Tauberian theorem of N. \textsc{Wiener} \cite{wiener340}: 

\begin{theorem}\label{thm:XI.16.Theorem.3} Let $x(t) \in L^{1}(-\infty, \infty)$ be such that its Fourier transform 
\[
(2 \pi) ^{-1/2} \int_{- \infty} ^{\infty} x (t) \exp (- i \cdot \xi \cdot t) \dd t
\]
does not vanish for any real $\xi$. Then, for any $y(t) \in L^{1}(-\infty, \infty)$ and $\varepsilon > 0$, we can find the real numbers $\beta$'s, the complex numbers $\alpha$'s and a positive integer $N$ in such a way that 
\begin{equation}\label{eq:XI.16.16}
\int_{- \infty} ^{\infty} \left| y (t) - \sum\limits_{j=1} ^{N} \alpha_{j} x (t - \beta_{j}) \right| \dd t < \varepsilon .
\end{equation}
\end{theorem}

\begin{proof}
 It suffices to find $\tilde{z} \in L^1$ in such a way that 
\begin{equation}\label{eq:XI.16.17}
\left\| y - x \times \tilde {z} \right\| \leqq \varepsilon / 2. \end{equation}

We first show that 
\begin{equation}\lim\limits _{\alpha \rightarrow \infty} \int_{- \infty} ^{\infty} | y (t) - y ^{(\alpha)} (t) | \dd t = 0, \text{ where } y ^{(\alpha)} (t) = \frac {1}{\pi} \int_{- \infty} ^{\infty} y (t - s) \frac {1 - \cos \alpha s}{\alpha s ^{2}} \dd s. \label{eq:XI.16.18}\end{equation}

For, by $\int_{-\infty}^{\infty}(1 - \cos \alpha s)(\alpha s^2)^{-1}\dd s = \pi (\alpha >0)$, the left side is 
\[
\leqq (\pi) ^{-1} \lim\limits _{\alpha \rightarrow \infty} \int_{- \infty} ^{\infty} \frac {(1 - \cos s)}{s ^{2}} \dd s \left\{\int_{- \infty} ^{\infty} \left| y (t) - y \left(t - \frac {s}{\alpha}\right)\right| \dd t \right\} = 0.
\]


Next we have 
\begin{equation}\label{eq:XI.16.19}
(2 \pi) ^{-1/2} \int_{- \infty} ^{\infty} \left(\frac {2}{\pi}\right) ^{1/2} \frac {1 - \cos \alpha u}{\alpha u ^{2}} e ^{- i u \xi} \dd u = \begin{cases}1 - | \xi | / \alpha & (| \xi | < \alpha), \\ 0 & (| \xi | \geqq \alpha).\end{cases}
\end{equation}

For, we have 
\[
\begin{aligned}
&(2 \pi) ^{-1/2} \int_{- \alpha} ^{\alpha} (1 - | \xi | / \alpha) e ^{i \xi u} \dd \xi  = \left(\frac {2}{\pi}\right) ^{1/2} \int_{0} ^{\alpha} \left(1 - \frac {\xi}{\alpha}\right) \cos u \xi \cdot \dd \xi \\
&= \left(\frac {2}{\pi}\right) ^{1/2} \frac {1}{u} \int_{0} ^{\alpha} \left(1 - \frac {\xi}{\alpha}\right) d (\sin u \xi) = \left(\frac {2}{\pi}\right) ^{1/2} \int_{0} ^{\alpha} \frac {\sin u \xi}{u \alpha} \dd \xi = \left(\frac {2}{\pi}\right) ^{1/2} \frac {1 - \cos u \alpha}{u ^{2} \alpha},
\end{aligned}
\]
and so we have only to apply Plancherel's theorem in Chapter~\ref{chap:VI}, \ref{sec:VI.2}. Hence, by the Parseval relation of the Fourier transform, 
\begin{equation}\label{eq:XI.16.20}
y ^{(\alpha)} (t) \text{ satisfies } \int_{- \infty} ^{\infty} y ^{(\alpha)} (t) \exp (- i \cdot \xi \cdot t) \dd t = 0 \text{ if } | \xi | \geqq \alpha .
\end{equation}

Therefore, we may assume that $y \in L^{1}(-\infty, \infty)$ in \eqref{eq:XI.16.17} satisfies the condition 
\begin{equation}\label{eq:XI.16.21}
\int_{- \infty} ^{\infty} y (t) \exp (- i \cdot \xi \cdot t) \dd t = 0 \text{ whenever } | \xi | \geqq \alpha .
\end{equation}

Next we choose a positive number $\beta$ and a sufficiently large positive number $\gamma$ such that $[-\beta - \gamma, -\beta + \gamma]$ and $[\beta - \gamma, \beta + \gamma]$ both contain $[- \alpha, \alpha]$. Let $C_1(\xi)$ and $C_2(\xi)$ be the defining functions of the intervals $[- \gamma, \gamma]$ and $[- \beta, \beta]$, respectively. Then 
\begin{equation}\label{eq:XI.16.22}
\begin{yoslines}
&u (\xi) = (2 \beta) ^{-1} \int_{- \infty} ^{\infty} C _{1} (\xi - \eta) C _{2} (\eta) \dd \eta = 1 \text{ for } \xi \in [ - \alpha , \alpha ], \\
&\quad = 0 \text{ for sufficiently large } | \xi |, 0 \leqq u (\xi) \leqq 1 \text{ for all real } \xi .
\end{yoslines}
\end{equation}

By the Parseval relation of the Fourier transform, we have 
\begin{equation}\label{eq:XI.16.23}
\hat {\boldsymbol {u}} (t) = \frac {1}{2 \beta} (2 \pi) ^{1/2} \hat {C} _{1} (t) \hat {C} _{2} (t).
\end{equation}

Thus, by Plancherel's theorem, we see that $\hat{\pmb{u}}(t)$ belongs to $L^1(-\infty, \infty) \cap L^2(-\infty, \infty)$. Hence we may apply Theorem~\ref{thm:XI.16.Theorem.1} so that 
$f(t) = \hat{u}(t) = (2\pi)^{-1/2}(u,I_t)$ with a certain maximal ideal $I_{t}\neq I_{0}$ 


Moreover, by Plancherel's theorem, the inverse Fourier transform of $f(t)$ is equal to $u(\xi)$, i.e., we have 
\begin{equation}\label{eq:XI.16.24}
u (\xi) = (2 \pi) ^{-1/2} (f, I _{- \xi}).
\end{equation}

We next set $\tilde{g} = e - (2\pi)^{-1/2}f$. Then, by what was proved above, 
\begin{equation}\label{eq:XI.16.25}
\begin{aligned}
&0 \leqq (\tilde {g}, I _{\xi}) \leqq 1; \quad (\tilde {g}, I _{\xi}) = 0 \text{ whenever } \xi \in [ - \alpha , \alpha ]; \\
&\qquad (\tilde{g}, I_{\xi}) = 1 \text{ for all sufficiently large } |\xi|.
\end{aligned}
\end{equation}


On the other hand, we have, for $x^{*}(t) = \overline{x(-t)}$, the relation $(\overline{x,I_{\xi}}) = (x^{*},I_{\xi})$. Thus, by hypothesis, we have $(x^{*}\times x,I_{\xi}) = |(x,I_{\xi})|^2 > 0$ for all real $\xi$. Hence the element 
\[
\tilde {g} + x ^{*} \times x \in L ^{1}
\]
satisfies the condition that $(\tilde{g} + x^{*}\times x, I) > 0$ for all maximal ideals $I$ of $L^1$. Consequently, the inverse $(\tilde{g} + x^{*}\times x)^{-1}$ in $L^1$ does exist. 

We define 
\begin{equation}\label{eq:XI.16.26}
\tilde {z} = (\tilde {g} + x \times x ^{*}) ^{-1} \times x ^{*} \times y.
\end{equation}

Then, by \eqref{eq:XI.16.4}, the element $x \times \tilde{z}$ belongs to $L^1(-\infty, \infty)$. Moreover, we have for every real number $\xi$, 
\[
(x \times \tilde {z}, I _{\xi}) = (x, I _{\xi}) (\tilde {z}, I _{\xi}) = (x, I _{\xi}) \frac {(x ^{*} , I _{\xi}) (y , I _{\xi})}{(\tilde {g} , I _{\xi}) + (x , I _{\xi}) (x ^{*} , I _{\xi})}.
\]

Hence, by \eqref{eq:XI.16.25} and the hypothesis that $(y, I_{\xi}) = 0$ whenever $|\xi| \geqq \alpha$, we obtain 
\[
(x \times \tilde {z}, I _{\xi}) = (y, I _{\xi}) \text{ for all real } \xi .
\]

The normed ring $L^1$ is semi-simple by Theorem~\ref{thm:XI.16.Theorem.2}. Hence we must have $x \times \tilde{z} = y$. We have thus proved Theorem~\ref{thm:XI.16.Theorem.3}.
\end{proof}

\begin{corollaryu} Let $k_{1}(t)$ belong to $L^{1}(-\infty, \infty)$, and let its Fourier transform vanish for no real argument. Let $k_{2}(t)$ belong to $L^{1}(-\infty, \infty)$. Let $f(t)$ be Baire measurable and bounded over $(-\infty, \infty)$. Let, for a constant $C$, 
\begin{equation}\label{eq:XI.16.27}
\lim\limits _{t \rightarrow \infty} \int_{- \infty} ^{\infty} k _{1} (t - s) f (s) \dd s = C \int_{- \infty} ^{\infty} k _{1} (t) \dd t.
\end{equation}
\end{corollaryu}

Then 
\begin{equation}\label{eq:XI.16.28}
\lim\limits _{t \rightarrow \infty} \int_{- \infty} ^{\infty} k _{2} (t - s) f (s) \dd s = C \int_{- \infty} ^{\infty} k _{2} (t) \dd t.
\end{equation}

\begin{proof}
 We may plainly suppose that $C$ = $0$. And \eqref{eq:XI.16.28} holds for $k_{2}(t)$ of the form $k_{2}(t) = (x \times k_{1})(t)$ with $x(t) \in L^{1}(-\infty, \infty)$. It is easy to prove that \eqref{eq:XI.16.28} holds for $k_{2}(t)$ if $k_{2}(t) = \slim\limits_{n \to \infty} k^{(n)}(t)$ in $L^{1}(-\infty, \infty)$ with $k^{(n)}(t) \in L^{1}(-\infty, \infty)$ for which \eqref{eq:XI.16.28} holds $(n = 1, 2, \ldots)$. Hence, by Theorem~\ref{thm:XI.16.Theorem.3}, we see that \eqref{eq:XI.16.28} holds for every $k_{2}(t) \in L^{1}(-\infty, \infty)$. 
\end{proof}

\begin{remarku} N. \textsc{Wiener} (\cite{wiener339}, \cite{wiener340} and \cite{wiener341}) has applied the above Corollary to a unified treatment of classical results concerning the limit relations in series and integrals which include a new proof of the prime number theorem. Cf. also H. R. \textsc{Pitt} \cite{pitt262}. The above proof of Theorem~\ref{thm:XI.16.Theorem.3} is adapted from M. \textsc{Fukamiya} \cite{fukamiya72} and I. E. \textsc{Segal} \cite{segal301}. Cf. also M. A. \textsc{Naimark} \cite{naimark221} and C. E. \textsc{Rickart} \cite{rickart271} and the bibliographies cited in these books. For the sake of comprehension of the scope of the above Corollary, we shall reproduce Wiener's deduction of the special Tauberian theorem. The special Tauberian theorem, as formulated by J. E. \textsc{Littlewood} reads: 
\end{remarku}

\begin{theorem}\label{thm:XI.16.Theorem.4} Let $\sum\limits_{n=0}^{\infty} a_n x^n$ converge to $s(x)$ for $|x| < 1$, and let 
\begin{equation}\label{eq:XI.16.29}
\lim\limits _{x \rightarrow 1-0} s (x) = C.
\end{equation}
\end{theorem}

Let, moreover, 
\begin{equation}\label{eq:XI.16.30}
\sup\limits _{n \geqq 1} n | a _{n} | = K < \infty .
\end{equation}

Then 
\begin{equation}\label{eq:XI.16.31}
\sum\limits_{n=0} ^{\infty} a _{n} = C.
\end{equation}

\begin{proof}
 Put $f(x) = \sum\limits_{n=0}^{[x]} a_n$. Then we see, by 
\begin{align}
\left| f (x) - s \left(e ^{-1/x}\right) \right|& = \left| \sum\limits_{n=1} ^{[x]} a _{n} \left(1 - e ^{-n/x}\right) - \sum\limits_{[x]+1} ^{\infty} a _{n} e ^{-n/x} \right| \notag \\
&\leqq \sum\limits_{n=1} ^{[x]} \frac {K}{n} \frac {n}{x} + \sum\limits_{[x]+1} ^{\infty} \frac {K}{n} e ^{-n/x} \leqq 2 K + K \int_{[x]} ^{\infty} e ^{-u/x} u ^{-1} \dd u \notag \\
&\leqq 2 K + K \int_{1} ^{\infty} e ^{-u} u ^{-1} \dd u = \text{ constant }, \notag
\end{align}
that $f(x)$ is bounded. 

Hence, by partial integration, 
\[
s \left(e ^{-x}\right) = \sum\limits_{n=0} ^{\infty} a _{n} e ^{-nx} = \int_{-0} ^{\infty} e ^{-ux} \dd f (u) = \int_{0} ^{\infty} x e ^{-ux} f (u) \dd u.
\]

Thus we have 
\begin{equation}\tag{31'}\label{eq:XI.16.31'}
C = \lim\limits _{x \rightarrow 0} \int_{0} ^{\infty} x e ^{-ux} f (u) \dd u = \lim\limits _{\xi \rightarrow \infty} \int_{- \infty} ^{\infty} e ^{- \xi} e ^ {- e ^{\eta - \xi}} f (e ^{\eta}) e ^{\eta} \dd \eta .
\end{equation}

This may be written as 
\begin{equation}\label{eq:XI.16.32}
\lim\limits _{t \rightarrow \infty} \int_{- \infty} ^{\infty} k _{1} (t - s) f \left(e ^{s}\right) \dd s = C \int_{- \infty} ^{\infty} k _{1} (t) \dd t, \text{ where } k _{1} (t) = e ^{-t} e ^ {- e ^{-t}},
\end{equation}
since 
\[
\int_{- \infty} ^{\infty} k _{1} (t) \dd t = \int_{- \infty} ^{\infty} e ^{-t} e ^ {- e ^{-t}} \dd t = \int_{0} ^{\infty} e ^{-x} \dd x = 1.
\]

Furthermore, we have 
\[
\int_{- \infty} ^{\infty} k _{1} (t) e ^{-iut} \dd t = \int_{0} ^{\infty} x ^{iu} e ^{-x} \dd x = \Gamma (1 + i u) \neq 0.
\]

Thus we can apply the Corollary and obtain, for 
\[
k _{2} (t) = 0 (t < 0); k _{2} (t) = e ^{-t} (t > 0),
\]
the limit relation 
\[
\begin{aligned}
&C= C \int_{0} ^{\infty} e ^{-t} \dd t = C \int_{- \infty} ^{\infty} k _{2} (t) \dd t = \lim\limits _{t \rightarrow \infty} \int_{- \infty} ^{\infty} k _{2} (t - s) f (e ^{s}) \dd s = \lim\limits _{x \rightarrow \infty} x ^{-1} \int_{0} ^{x} f (y) \dd y.
\end{aligned}
\]

It follows that, if $\lambda > 0$, 
\begin{equation}\label{eq:XI.16.33}
\begin{aligned}
C &= \frac {(1 + \lambda) C - C}{\lambda} = \lim\limits _{x \rightarrow \infty} \frac {1}{\lambda x} \left\{\int_{0} ^{(1+ \lambda) x} f (y) \dd y - \int_{0} ^{x} f (y) \dd y \right\}
\\
&= \lim\limits _{x \rightarrow \infty} \frac {1}{\lambda x} \int_{x} ^{(1+ \lambda) x} f (y) \dd y = \lim\limits _{x \rightarrow \infty} \left\{f (x) + \frac {1}{\lambda x} \int_{x} ^{(1+ \lambda) x} [ f (y) - f (x) ] \dd y \right\}.
\end{aligned}
\end{equation}

On the other hand, we have, by \eqref{eq:XI.16.30}, 
\[
\begin{aligned}
&\left| \frac {1}{\lambda x} \int_{x} ^{(1+ \lambda) x} [ f (y) - f (x) ] \dd y \right|  \leqq \frac {1}{\lambda x} \int_{x} ^{(1+ \lambda) x} \sum\limits_{[x]+1} ^{[y]} \frac {K}{n} \dd y \leqq \sum\limits_{[x]+1} ^{[ (1+ \lambda) x ]} \frac {K}{[ x ]} \leqq \frac {[ \lambda x ] K}{[ x ]} \leqq 2 \lambda K
\end{aligned}
\]
for sufficiently large value of $x$. Hence, by \eqref{eq:XI.16.33}, 
\[
\varlimsup\limits_{x \rightarrow \infty} | f (x) - C | \leqq 2 \lambda K,
\]
and since $\lambda$ is any positive number, we get 
\[
\lim\limits _{x \rightarrow \infty} f (x) = C.
\]

Thus we have proved \eqref{eq:XI.16.31}. 
\end{proof}
